% Traduction et production de textes injonctifs liés aux médias et à la communication — Allemand Tle A — Cours signé OnBuch+
% Programme officiel MINESEC. Compiler avec : tectonic AL21-traduction-et-production-de-textes-injonctifs-lies-aux-media.tex
\newcommand{\DOCMATIERE}{Allemand}
\newcommand{\DOCNIVEAU}{Tle A}
\newcommand{\DOCMODULE}{Chapitre 4 : Média et communication}
\newcommand{\DOCLECON}{AL21}
\newcommand{\DOCTITRE}{Traduction et production de textes injonctifs liés aux médias et à la communication}
% OnBuch+ — préambule « cours signés » ENRICHI (compilé avec tectonic / XeLaTeX)
% Système visuel « Pop » d'OnBuch+ : fond crème (jamais blanc), bordures encre
% épaisses, ombres « dures » décalées, titres Archivo Black en capitales.
% Version MATHÉMATIQUES : graphes (pgfplots/tikz), boîtes pédagogiques
% (définition, propriété, méthode, attention, le savais-tu, exercice résolu,
% exercice, TP, bilan), page de garde et signature OnBuch+ ancrée en bas.
%
% Macros attendues dans main.tex AVANT \input{preamble} :
%   \DOCMATIERE  \DOCNIVEAU  \DOCMODULE  \DOCLECON (numéro)  \DOCTITRE
\documentclass[11pt]{article}

\usepackage{fontspec}
\usepackage[ngerman,french]{babel} % français = langue principale ; l'allemand via \all{..} / textebox
\usepackage[a4paper,margin=2cm,top=3.4cm,bottom=2.8cm,headheight=1.4cm,footskip=1.3cm]{geometry}
\usepackage{amsmath,amssymb,amsfonts}
\usepackage{mathtools} % \xmapsto, \coloneqq... souvent utilisés par les modèles
\usepackage{cancel} % simplifications barrées (bilans de forces)
% \abs{z} (module, valeur absolue) et \norm{u} (norme), utilisés par les modèles
\providecommand{\abs}{}\let\abs\relax\DeclarePairedDelimiter{\abs}{\lvert}{\rvert}
\providecommand{\norm}{}\let\norm\relax\DeclarePairedDelimiter{\norm}{\lVert}{\rVert}
\usepackage[table]{xcolor} % [table] : fournit aussi \rowcolors (alternance de lignes)
\usepackage{pagecolor}
\usepackage{tikz}
\usetikzlibrary{arrows.meta,positioning,calc,shapes.geometric,shapes.misc,shapes.arrows,decorations.pathmorphing,decorations.pathreplacing,decorations.markings,patterns,shadows,mindmap,trees,fit,backgrounds,matrix,intersections,angles,quotes}
\usepackage{pgfplots}
\usepackage[version=4]{mhchem} % équations nucléaires \ce{^{235}_{92}U}
\usepackage[european,siunitx]{circuitikz} % schémas électriques
\pgfplotsset{compat=1.18}
\usepgfplotslibrary{fillbetween,groupplots} % aires entre courbes (calcul intégral)
\usepackage{enumitem}
\usepackage{fancyhdr}
% [most] charge toutes les bibliothèques tcolorbox (listings, minted, raster,
% documentation...) et gonfle inutilement la mémoire TeX sur un document long
% (~300 boîtes) : on ne charge que ce qui est réellement utilisé.
\usepackage[skins,breakable]{tcolorbox}
\usepackage{graphicx}
\usepackage{colortbl}
\usepackage{booktabs}
\usepackage{tabularx}
\usepackage{diagbox} % case d'en-tête diagonale des tableaux à double entrée
% Colonnes extensibles centrée / gauche / droite (C, L, R), souvent utilisées par les modèles
\newcolumntype{C}{>{\centering\arraybackslash}X}
\newcolumntype{L}{>{\raggedright\arraybackslash}X}
\newcolumntype{R}{>{\raggedleft\arraybackslash}X}
\usepackage{array}
\usepackage{multicol}
\usepackage{multirow}
\usepackage{siunitx}
% parse-numbers=false : accepte aussi \SI{2,5 \times 10^{-3}}{...}
\sisetup{locale=FR,output-decimal-marker={,},group-minimum-digits=4,parse-numbers=false}
\DeclareSIUnit\ohm{\ensuremath{\symup{\Omega}}} % U+2126 absent de la police
\DeclareSIUnit\division{div} % réglages d'oscilloscope (ms/div, V/div)
\DeclareSIUnit\tr{tr} % tours (vitesse de rotation, ex. \SI{1500}{\tr\per\minute})
\DeclareSIUnit\molar{M} % concentration molaire (ex. \SI{3}{\molar}, \SI{1}{\micro\molar})
\DeclareSIUnit\an{an} % année (le modèle confond parfois avec \year, un compteur TeX) — ex. \SI{5}{\kilo\meter\per\an}
\DeclareSIUnit\FCFA{FCFA} % franc CFA (ex. \SI{500}{\FCFA\per\kilo\gram})
\DeclareSIUnit\degreeBrix{\textdegree Brix} % teneur en sucre d'un sirop/jus (réfractométrie)
\DeclareSIUnit\month{mois} % le modèle utilise parfois \month, un compteur TeX, comme unité de durée
\usepackage{qrcode}
% \degree : commande fréquemment utilisée par les modèles pour un angle ou une
% température en dehors de \SI{}{} (non fournie par les paquets chargés).
\providecommand{\degree}{^{\circ}}
% \qed (fin de démonstration), utilisé par les modèles sans amsthm
\providecommand{\qed}{\hfill$\square$}
% \barre : double barre des valeurs interdites dans un tableau de variations (tkz-tab)
\providecommand{\barre}{\ensuremath{\|}}
% environnement proof (amsthm non chargé) utilisé par les modèles
\ifdefined\proof\else\newenvironment{proof}[1][Démonstration]{\par\noindent\textit{#1.}\ }{\qed\par}\fi
\usepackage{lastpage}
\usepackage{hyperref}
\hypersetup{hidelinks}

% ============================================================
% Palette « Pop » OnBuch+ — mêmes hex que la classe P (plus_ui.dart)
% ============================================================
\definecolor{poporange}{HTML}{F59321}
\definecolor{popdark}{HTML}{E07A0C}
\definecolor{popcream}{HTML}{FFF6E8}
\definecolor{popink}{HTML}{1C1714}
\definecolor{poppurple}{HTML}{7A5AE0}
\definecolor{popgreen}{HTML}{2E9E6A}
\definecolor{poppink}{HTML}{E0457B}
\definecolor{popblue}{HTML}{2F7DE1}
\definecolor{popgold}{HTML}{C98A12}
\definecolor{popmuted}{HTML}{8A7F76}
\definecolor{popline}{HTML}{EADFD2}
\definecolor{popgray}{HTML}{8A7F76}
\colorlet{popgrey}{popgray}
% Alias tolérants (le modèle nomme parfois les couleurs différemment)
\colorlet{popwhite}{white}
\colorlet{popblack}{popink}
\colorlet{popred}{poppink}
\colorlet{popyellow}{popgold}
% Teintes claires (fonds de tableaux/encadrés) — le modèle nomme parfois
% les couleurs avec un suffixe L (ex. popblueL) sans les avoir définies.
\colorlet{poporangeL}{poporange!20}
\colorlet{popdarkL}{popdark!20}
\colorlet{poppurpleL}{poppurple!20}
\colorlet{popgreenL}{popgreen!20}
\colorlet{poppinkL}{poppink!20}
\colorlet{popblueL}{popblue!20}
\colorlet{popgoldL}{popgold!20}
\colorlet{popredL}{popred!20}
\colorlet{popgrayL}{popgray!20}
% Teintes claires pour fonds de tableaux / zones
\colorlet{popblueL}{popblue!10!white}
\colorlet{poporangeL}{poporange!14!white}
\colorlet{popgreenL}{popgreen!12!white}
\colorlet{poppurpleL}{poppurple!12!white}
\colorlet{poppinkL}{poppink!12!white}
\colorlet{popgoldL}{popgold!14!white}

\pagecolor{popcream}

% ============================================================
% Typographie
% ============================================================
\defaultfontfeatures{Ligatures=TeX}
\setmainfont{PlusJakartaSans-Medium.ttf}[Path=../../../fonts/,BoldFont=PlusJakartaSans-Bold.ttf]
% Symboles, API et emojis absents de la police principale : repli sur DejaVu Sans (caractères actifs)
\newfontface\symfont{DejaVuSans.ttf}[Path=../../../fonts/]
\makeatletter
\newcommand\symactive[1]{\catcode#1=13 \begingroup\lccode`\~=#1 \lowercase{\endgroup\def~}{{\symfont\char#1}}}
\newcommand\symblank[1]{\catcode#1=13 \begingroup\lccode`\~=#1 \lowercase{\endgroup\def~}{}}
\newcommand\symhyph[1]{\catcode#1=13 \begingroup\lccode`\~=#1 \lowercase{\endgroup\def~}{\mbox{-}}}
\newcount\symc
\newcommand\symrange[2]{\symc=#1 \@whilenum\symc<#2 \do{\expandafter\symactive\expandafter{\the\symc}\advance\symc by 1 }}
\newcommand\symblankrange[2]{\symc=#1 \@whilenum\symc<#2 \do{\expandafter\symblank\expandafter{\the\symc}\advance\symc by 1 }}
\makeatother
\symrange{"0250}{"02FF}\symactive{"2713}\symactive{"2714}\symactive{"2717}\symactive{"2718}\symactive{"2610}\symactive{"2611}\symactive{"2612}
\symhyph{"2011}\symhyph{"2E1A}\symblankrange{"1F300}{"1FAFF}\symblank{"FE0F}
% Maths en sans-serif (Fira Math) avec chiffres et lettres droites de Plus
% Jakarta, pour que les formules se fondent dans le texte.
\usepackage[math-style=ISO,bold-style=ISO]{unicode-math}
\setmathfont{FiraMath-Regular.otf}[Path=../../../fonts/]
\setmathfont{PlusJakartaSans-Medium.ttf}[Path=../../../fonts/,range={up/num,up/{Latin,latin},\mathup}]
\usepackage{icomma} % virgule décimale sans espace en mode math
% Caractères mathématiques Unicode absents des polices de texte (Plus Jakarta,
% Archivo Black) : rendus par leur équivalent mathématique, en texte comme en
% formule — sinon ils s'affichent en carré vide (ex. « Arithmétique dans ℤ »).
\usepackage{newunicodechar}
\newunicodechar{ℝ}{\ensuremath{\mathbb{R}}}
\newunicodechar{ℤ}{\ensuremath{\mathbb{Z}}}
\newunicodechar{ℂ}{\ensuremath{\mathbb{C}}}
\newunicodechar{ℕ}{\ensuremath{\mathbb{N}}}
\newunicodechar{ℚ}{\ensuremath{\mathbb{Q}}}
\newunicodechar{𝔻}{\ensuremath{\mathbb{D}}}
\newunicodechar{α}{\ensuremath{\alpha}}
\newunicodechar{β}{\ensuremath{\beta}}
\newunicodechar{γ}{\ensuremath{\gamma}}
\newunicodechar{δ}{\ensuremath{\delta}}
\newunicodechar{ε}{\ensuremath{\varepsilon}}
\newunicodechar{θ}{\ensuremath{\theta}}
\newunicodechar{λ}{\ensuremath{\lambda}}
\newunicodechar{μ}{\ensuremath{\mu}}
\newunicodechar{σ}{\ensuremath{\sigma}}
\newunicodechar{τ}{\ensuremath{\tau}}
\newunicodechar{φ}{\ensuremath{\varphi}}
\newunicodechar{ψ}{\ensuremath{\psi}}
\newunicodechar{ω}{\ensuremath{\omega}}
\newunicodechar{Σ}{\ensuremath{\Sigma}}
% majuscules produites par \MakeUppercase dans les titres (α-aminés -> Α)
\newunicodechar{Α}{\ensuremath{\alpha}}
\newunicodechar{Β}{\ensuremath{\beta}}
\newunicodechar{Γ}{\ensuremath{\Gamma}}
\newunicodechar{Δ}{\ensuremath{\Delta}}
\newunicodechar{Ε}{\ensuremath{\varepsilon}}
\newunicodechar{Θ}{\ensuremath{\Theta}}
\newunicodechar{Λ}{\ensuremath{\Lambda}}
\newunicodechar{Μ}{\ensuremath{\mu}}
\newunicodechar{Τ}{\ensuremath{\tau}}
\newunicodechar{Φ}{\ensuremath{\Phi}}
\newunicodechar{Ψ}{\ensuremath{\Psi}}
\newunicodechar{Ω}{\ensuremath{\Omega}}
\newunicodechar{↦}{\ensuremath{\mapsto}}
\newunicodechar{∈}{\ensuremath{\in}}
\newunicodechar{∉}{\ensuremath{\notin}}
\newunicodechar{∧}{\ensuremath{\wedge}}
\newunicodechar{⇔}{\ensuremath{\Leftrightarrow}}
\newunicodechar{⟺}{\ensuremath{\Longleftrightarrow}}
\newunicodechar{⇒}{\ensuremath{\Rightarrow}}
\newunicodechar{⟹}{\ensuremath{\Longrightarrow}}
\newunicodechar{∘}{\ensuremath{\circ}}
\newunicodechar{∪}{\ensuremath{\cup}}
\newunicodechar{∩}{\ensuremath{\cap}}
\newunicodechar{≡}{\ensuremath{\equiv}}
\newunicodechar{⊂}{\ensuremath{\subset}}
\newunicodechar{⊥}{\ensuremath{\perp}}
\newunicodechar{∃}{\ensuremath{\exists}}
\newunicodechar{∀}{\ensuremath{\forall}}
\newunicodechar{∛}{\ensuremath{\sqrt[3]{\,}}}
\newunicodechar{∜}{\ensuremath{\sqrt[4]{\,}}}
\newunicodechar{⃗}{\ensuremath{{}^{\to}}}
\newunicodechar{ⁿ}{\ifmmode^{n}\else\textsuperscript{\ensuremath{n}}\fi}
\newunicodechar{ˣ}{\ifmmode^{x}\else\textsuperscript{\ensuremath{x}}\fi}
\newunicodechar{ᵇ}{\ifmmode^{b}\else\textsuperscript{\ensuremath{b}}\fi}
\newunicodechar{ʳ}{\ifmmode^{r}\else\textsuperscript{\ensuremath{r}}\fi}
\newunicodechar{ᵘ}{\ifmmode^{u}\else\textsuperscript{\ensuremath{u}}\fi}
\newunicodechar{ʸ}{\ifmmode^{y}\else\textsuperscript{\ensuremath{y}}\fi}
\newunicodechar{ᵃ}{\ifmmode^{a}\else\textsuperscript{\ensuremath{a}}\fi}
\newunicodechar{ᵉ}{\ifmmode^{e}\else\textsuperscript{\ensuremath{e}}\fi}
\newunicodechar{ᵏ}{\ifmmode^{k}\else\textsuperscript{\ensuremath{k}}\fi}
\newunicodechar{ᵖ}{\ifmmode^{p}\else\textsuperscript{\ensuremath{p}}\fi}
\newunicodechar{ᴺ}{\ifmmode^{N}\else\textsuperscript{\ensuremath{N}}\fi}
\newunicodechar{ᶜ}{\ifmmode^{c}\else\textsuperscript{\ensuremath{c}}\fi}
\newunicodechar{ʰ}{\ifmmode^{h}\else\textsuperscript{\ensuremath{h}}\fi}
\newunicodechar{ᵠ}{\ifmmode^{q}\else\textsuperscript{\ensuremath{q}}\fi}
\newunicodechar{ₙ}{\ifmmode_{n}\else\textsubscript{\ensuremath{n}}\fi}
\newunicodechar{ₐ}{\ifmmode_{a}\else\textsubscript{\ensuremath{a}}\fi}
\newunicodechar{ᵢ}{\ifmmode_{i}\else\textsubscript{\ensuremath{i}}\fi}
\newunicodechar{ᵣ}{\ifmmode_{r}\else\textsubscript{\ensuremath{r}}\fi}
\newunicodechar{ₖ}{\ifmmode_{k}\else\textsubscript{\ensuremath{k}}\fi}
\newunicodechar{ᵦ}{\ifmmode_{\beta}\else\textsubscript{\ensuremath{\beta}}\fi}
\newunicodechar{ₓ}{\ifmmode_{x}\else\textsubscript{\ensuremath{x}}\fi}
\newfontfamily\popdisplay{ArchivoBlack-Regular.ttf}[Path=../../../fonts/]
\newfontfamily\poplogo{SpaceGrotesk-Bold.ttf}[Path=../../../fonts/]
\newfontfamily\popsemi{PlusJakartaSans-SemiBold.ttf}[Path=../../../fonts/]
\newfontfamily\popxbold{PlusJakartaSans-ExtraBold.ttf}[Path=../../../fonts/]
\newfontfamily\popxboldls{PlusJakartaSans-ExtraBold.ttf}[Path=../../../fonts/,LetterSpace=3.0]

\setlist[enumerate]{leftmargin=1.7em,itemsep=5pt,topsep=4pt}
\setlist[itemize]{leftmargin=1.7em,itemsep=4pt,topsep=4pt,label={\color{poporange}\textbullet}}
\setlength{\parindent}{0pt}
\setlength{\parskip}{5pt}
\renewcommand{\arraystretch}{1.25}

% pgfplots : style Pop par défaut
\pgfplotsset{
  every axis/.append style={axis line style={popink,line width=1pt},tick style={popink},
    grid style={popline,line width=0.5pt},label style={font=\small},tick label style={font=\footnotesize}},
  popcurve/.style={poporange,line width=1.8pt,no markers,smooth},
}
% Même style disponible dans un \draw TikZ simple (hors pgfplots)
\tikzset{popcurve/.style={poporange,line width=1.8pt}}
\tikzset{popgrid/.style={popline,very thin}}
\colorlet{popcurve}{poporange} % le nom du style est parfois utilisé comme couleur

% ============================================================
% Pill / squiggle
% ============================================================
\newcommand{\poppill}[3]{%
  \tikz[baseline=-0.55ex]\node[draw=popink,line width=1.2pt,rounded corners=2.6pt,%
    fill=#2,inner xsep=6.5pt,inner ysep=3pt]{%
    \popxboldls\fontsize{7}{7}\selectfont\color{#3}\MakeUppercase{#1}};%
}
\newcommand{\popsquiggle}[1]{%
  \tikz[baseline=-0.4ex]{
    \draw[#1,line width=2.6pt,line cap=round]
      plot [smooth] coordinates {(0,0.09) (0.34,0.01) (0.68,0.21) (1.02,0.01) (1.36,0.10)};
  }%
}
% Mot-clé surligné (marqueur orange sous le texte)
\newcommand{\cle}[1]{{\popxbold\color{popdark}#1}}

% ============================================================
% En-tête / pied
% ============================================================
\pagestyle{fancy}
\fancyhf{}
\renewcommand{\headrulewidth}{0pt}
\renewcommand{\footrulewidth}{0pt}
\lhead{\raisebox{-0.15ex}{\poplogo\fontsize{13}{13}\selectfont\color{popink}OnBuch\color{poporange}+}}
\rhead{\poppill{\DOCMATIERE\ \textbullet\ \DOCNIVEAU\ \textbullet\ Leçon \DOCLECON}{white}{popink}}
\lfoot{\popxbold\fontsize{7.6}{7.6}\selectfont\color{popmuted}\MakeUppercase{Cours signé OnBuch+}\ \textbullet\ {\popsemi\DOCTITRE}}
\rfoot{\tikz[baseline=-0.6ex]\node[rounded corners=8.5pt,fill=popink,minimum height=17pt,minimum width=17pt,inner xsep=5pt,inner sep=0]{\popxbold\fontsize{8}{8}\selectfont\color{popcream}\thepage\,/\,\pageref*{LastPage}};}

% ============================================================
% Titres
% ============================================================
\newcommand{\chaptitre}[2]{%
  \begin{tcolorbox}[enhanced,colback=poporange,colframe=popink,boxrule=2.6pt,arc=11pt,
      drop shadow=popink,left=15pt,right=15pt,top=12pt,bottom=11pt,before skip=3pt,after skip=18pt]
    \poppill{#2}{popink}{popcream}\par\vspace{9pt}
    {\raggedright\hyphenpenalty=10000\exhyphenpenalty=10000\popdisplay\fontsize{22}{24}\selectfont\color{popcream}\MakeUppercase{#1}\par}
  \end{tcolorbox}
}
% Section numérotée : pastille orange avec le numéro + titre Archivo
\newcounter{popsec}
\newcommand{\coursec}[1]{%
  \stepcounter{popsec}\setcounter{popsub}{0}%
  \par\vspace{16pt}\noindent
  \begin{minipage}{\linewidth}
  \tikz[baseline=-0.7ex]\node[draw=popink,line width=1.6pt,fill=poporange,rounded corners=4pt,minimum size=22pt,inner sep=0]{\popdisplay\fontsize{12}{12}\selectfont\color{popcream}\thepopsec};%
  \hspace{8pt}{\popdisplay\fontsize{15}{17}\selectfont\color{popink}\MakeUppercase{#1}}\par
  \vspace{4pt}\noindent{\color{poporange}\rule{36pt}{3.6pt}}
  \end{minipage}\par\nopagebreak\vspace{8pt}
}
\newcounter{popsub}
\newcommand{\courssub}[1]{%
  \stepcounter{popsub}%
  \par\vspace{9pt}\noindent{\popxbold\fontsize{11.5}{13}\selectfont\color{popink}{\color{poporange}\thepopsec.\thepopsub}\hspace{6pt}#1}\par\nopagebreak\vspace{4pt}
}

% ============================================================
% Boîtes pédagogiques — même recette : fond blanc, bordure épaisse colorée,
% ombre dure encre, badge plein. Titre optionnel [#1].
% ============================================================
\tcbset{popbase/.style={breakable,enhanced,colback=white,boxrule=2.3pt,arc=9pt,
  shadow={3pt}{-3pt}{0pt}{popink},left=14pt,right=14pt,top=11pt,bottom=11pt,before skip=11pt,after skip=15pt,
  fontupper=\popsemi\fontsize{10.6}{14.5}\selectfont\color{popink},
  fontlower=\popsemi\fontsize{10.6}{14.5}\selectfont\color{popink}}}
% Coche dessinée (le glyphe ✓ n'existe pas dans les polices du document)
\newcommand{\popcheck}[1]{\tikz[baseline=-0.5ex]\draw[#1,line width=1.6pt,line cap=round,line join=round](-0.13,0.0)--(-0.04,-0.09)--(0.13,0.1);}
\providecommand{\coche}{\popcheck{popgreen}}
\providecommand{\resultat}[1]{\par\smallskip\noindent\fcolorbox{popgreen}{popgreenL}{\parbox{\dimexpr\linewidth-2\fboxsep-2\fboxrule\relax}{\textbf{Résultat.} #1}}\par\smallskip}
\newcommand{\popboxhead}[3]{\poppill{#1}{#2}{white}\ifx\relax#3\relax\else\hspace{6pt}{\popxbold\fontsize{10.6}{12}\selectfont\color{popink}#3}\fi\par\vspace{7pt}}

\newtcolorbox{aretenir}[1][]{popbase,colframe=popgreen,before upper={\popboxhead{À retenir}{popgreen}{#1}}}
% Texte en allemand (document de lecture, dialogue, extrait) : cadre
% d'étude. Un titre optionnel (source, auteur) s'affiche dans l'en-tête.
\newtcolorbox{textebox}[1][]{popbase,colframe=popblue,before upper={\popboxhead{Text}{popblue}{#1}\begin{otherlanguage}{ngerman}},after upper={\end{otherlanguage}}}
% Marques ✓ / ✗ (forme correcte / fautive) souvent utilisées par les modèles
\providecommand{\cross}{\textcolor{poppink}{\ensuremath{\times}}}
\providecommand{\xmark}{\cross}\providecommand{\crossmark}{\cross}\providecommand{\cmark}{\ensuremath{\checkmark}}
% Ligne de réponse / justification à compléter (inventée par les modèles)
\providecommand{\justification}[1]{\par\noindent\textit{Justification~:} \dotfill\par}
% \textipa (package tipa non chargé) : la police ne couvre pas l'API -> simple passage
\providecommand{\textipa}[1]{#1}
% Soulignement (ulem non chargé) : \uline, \ul
\providecommand{\uline}[1]{\underline{#1}}\providecommand{\ul}[1]{\underline{#1}}
% \glossaire{\item ...} inventée par les modèles : liste de glossaire
\providecommand{\glossaire}[1]{\begin{itemize}#1\end{itemize}}
% \alert (beamer) inventée par les modèles : texte mis en évidence
\providecommand{\alert}[1]{\textbf{\textcolor{poppink}{#1}}}
% variantes ASCII d'ordinaux écrites en macro
\providecommand{\iere}{\textsuperscript{re}}\providecommand{\ieme}{\textsuperscript{e}}\providecommand{\ere}{\textsuperscript{re}}\providecommand{\eme}{\textsuperscript{e}}\providecommand{\er}{\textsuperscript{er}}\providecommand{\ie}{\textsuperscript{e}}
% Ordinaux français écrits en macro par les modèles : 1\ière, 3\ième
\newcommand{\ière}{\textsuperscript{re}}\newcommand{\ième}{\textsuperscript{e}}
% \exemple (inventée par les modèles) : étiquette « Exemple : »
\providecommand{\exemple}{\textbf{Exemple~:}\ }
% \square utilisable aussi hors mode math (cases à cocher)
\AtBeginDocument{\let\squareMath\square\renewcommand{\square}{\ensuremath{\squareMath}}}
% Mot ou phrase en allemand à l'intérieur d'un paragraphe français
\newcommand{\all}[1]{\ifmmode\text{\textit{#1}}\else\begingroup\selectlanguage{ngerman}\itshape #1\endgroup\fi}
\let\esp\all % alias toléré
\newtcolorbox{exemplebox}[1][]{popbase,colframe=poporange,before upper={\popboxhead{Exemple}{poporange}{#1}}}
\newtcolorbox{definition}[1][]{popbase,colframe=popblue,before upper={\popboxhead{Définition}{popblue}{#1}}}
\newtcolorbox{propriete}[1][]{popbase,colframe=poppurple,before upper={\popboxhead{Propriété}{poppurple}{#1}}}
\newtcolorbox{methode}[1][]{popbase,colframe=popgold,before upper={\popboxhead{Méthode}{popgold}{#1}}}
\newtcolorbox{attention}[1][]{popbase,colframe=poppink,before upper={\popboxhead{Attention}{poppink}{#1}}}
\newtcolorbox{experience}[1][]{popbase,colframe=popblue,colback=popblueL,before upper={\popboxhead{Expérience}{popblue}{#1}}}
% « Le savais-tu ? » : carte violette pleine (contexte, culture, Cameroun)
\newtcolorbox{savaistu}[1][]{popbase,colback=poppurple,colframe=popink,
  fontupper=\popsemi\fontsize{10.4}{14.2}\selectfont\color{white},
  before upper={\poppill{Le savais-tu\,?}{white}{poppurple}\ifx\relax#1\relax\else\hspace{6pt}{\popxbold\color{white}#1}\fi\par\vspace{7pt}}}
% Exercice résolu : énoncé en haut, \tcblower puis la solution sur fond crème
\newcounter{popexo}\newcounter{popexr}
\newtcolorbox{exoresolu}[1][]{popbase,colframe=poporange,colbacklower=popcream,
  segmentation style={popink,line width=1.2pt,dashed},
  before upper={\stepcounter{popexr}\popboxhead{Exercice résolu \thepopexr}{poporange}{#1}},
  before lower={\poppill{Solution}{popink}{popcream}\par\vspace{6pt}}}
% Exercice (non résolu en place) — la correction est dans la partie Corrigés
\newtcolorbox{exercice}[1][]{popbase,colframe=popink,
  before upper={\stepcounter{popexo}\popboxhead{Exercice \thepopexo}{popink}{#1}}}
% Corrigé
\newtcolorbox{corrige}[1][]{popbase,colframe=popgreen,colback=popgreenL,
  before upper={\popboxhead{Corrigé}{popgreen}{#1}}}
% Objectifs / prérequis / bilan
\newtcolorbox{objectifs}{popbase,colframe=popink,colback=poporangeL,
  before upper={\popboxhead{Objectifs de la leçon}{popink}{}}}
\newtcolorbox{prerequis}{popbase,colframe=popink,colback=popblueL,
  before upper={\popboxhead{Prérequis}{popblue}{}}}
\newtcolorbox{bilan}{popbase,colframe=popink,colback=white,boxrule=2.8pt,
  before upper={\popboxhead{Fiche bilan}{poporange}{L'essentiel à connaître pour l'examen}}}
% Figure encadrée avec légende
\newtcolorbox{popfigure}[1][]{enhanced,breakable=false,colback=white,colframe=popline,boxrule=1.4pt,arc=8pt,
  left=10pt,right=10pt,top=10pt,bottom=8pt,before skip=10pt,after skip=12pt,halign=center,
  lower separated=false,halign lower=center,
  fontlower=\popsemi\fontsize{9}{11.5}\selectfont\color{popmuted},
  #1}
\newcommand{\legende}[1]{\par\vspace{6pt}{\popsemi\fontsize{9}{11.5}\selectfont\color{popmuted}#1}}

% ============================================================
% Signature OnBuch+
% ============================================================
\newcommand{\poplogoinline}[1]{{\poplogo\fontsize{#1}{#1}\selectfont\color{popink}OnBuch\color{poporange}+}}
\newcommand{\signatureonbuch}{%
  \begin{tcolorbox}[enhanced,colback=popink,colframe=popink,boxrule=2.2pt,arc=10pt,
      drop shadow=poporange,left=14pt,right=14pt,top=10pt,bottom=10pt,before skip=0pt,after skip=0pt]
    \tikz[baseline=-0.7ex]\node[circle,fill=popgreen,draw=popcream,line width=1.3pt,minimum size=22pt,inner sep=0]{\popcheck{white}};%
    \hspace{9pt}\begin{minipage}[c]{0.62\linewidth}
      {\popxbold\fontsize{12}{13}\selectfont\color{popcream}Signé\ \textbullet\ {\poplogo OnBuch\color{poporange}+}}\par\vspace{2pt}
      {\raggedright\popsemi\fontsize{8}{10.5}\selectfont\color{popline}\MakeUppercase{Équipe pédagogique OnBuch+}\\\MakeUppercase{Conforme au programme officiel MINESEC}\par}
    \end{minipage}\hfill
    {\poplogo\fontsize{22}{22}\selectfont\color{popcream}OnBuch\color{poporange}+}
  \end{tcolorbox}%
}

% ============================================================
% Page de garde
% #1 titre · #2 matière · #3 niveau · #4 module · #5 n° leçon · #6 durée ·
% #7 description / objectifs (texte ou itemize)
% ============================================================
\newcommand{\pagegarde}[7]{%
  \thispagestyle{empty}
  \newgeometry{a4paper,margin=1.7cm,top=1.6cm,bottom=1.4cm}%
  \begin{tikzpicture}[remember picture,overlay]
    \fill[poporange!18] ([xshift=-3.2cm,yshift=-3.6cm]current page.north east) circle (4.6cm);
    \fill[poppurple!14] ([xshift=2.4cm,yshift=7.6cm]current page.south west) circle (3.8cm);
  \end{tikzpicture}%
  {\poplogo\fontsize{30}{30}\selectfont\color{popink}OnBuch\color{poporange}+}\hfill
  \raisebox{0.6ex}{\poppill{Cours signé \textbullet\ Premium}{popink}{popcream}}\par
  \vspace{1pt}\popsquiggle{poppurple}\par
  \vspace{8pt}
  \begin{tcolorbox}[enhanced,colback=poporange,colframe=popink,boxrule=2.8pt,arc=13pt,
      drop shadow=popink,left=18pt,right=18pt,top=12pt,bottom=13pt,after skip=10pt]
    \poppill{#2 \textbullet\ #3}{popink}{popcream}\hspace{5pt}\poppill{#4}{white}{popink}\par\vspace{8pt}
    {\popdisplay\fontsize{14}{14}\selectfont\color{popink}LEÇON #5}\par\vspace{4pt}
    {\raggedright\hyphenpenalty=10000\exhyphenpenalty=10000\popdisplay\fontsize{25}{27}\selectfont\color{popcream}\MakeUppercase{#1}\par}
    \vspace{8pt}
    {\popxbold\fontsize{9.5}{9.5}\selectfont\color{popink}\MakeUppercase{Durée conseillée : #6}}
  \end{tcolorbox}
  \begin{tcolorbox}[enhanced,colback=white,colframe=popink,boxrule=2.2pt,arc=10pt,
      drop shadow=popink,left=16pt,right=16pt,top=10pt,bottom=10pt,after skip=8pt]
    \poppill{Dans cette leçon}{poppurple}{white}\par\vspace{5pt}
    \popsemi\fontsize{9.3}{12.6}\selectfont\color{popink}#7
  \end{tcolorbox}
  \noindent\begin{minipage}[t]{0.487\linewidth}
    \begin{tcolorbox}[enhanced,colback=popgreen,colframe=popink,boxrule=2.2pt,arc=10pt,
        drop shadow=popink,left=11pt,right=11pt,top=10pt,bottom=10pt,height=3.1cm,valign=center]
      \tcbox[colback=white,colframe=popink,boxrule=1.3pt,arc=5pt,boxsep=0pt,nobeforeafter,
        left=3pt,right=3pt,top=3pt,bottom=3pt,valign=center]{\qrcode[height=1.9cm]{https://play.google.com/store/apps/details?id=com.luvvix.onbuch&hl=fr}}%
      \hspace{8pt}\begin{minipage}[c]{0.5\linewidth}
        {\popdisplay\fontsize{11}{12.5}\selectfont\color{white}\MakeUppercase{Télécharge OnBuch}}\par\vspace{4pt}
        {\raggedright\popsemi\fontsize{8.4}{11}\selectfont\color{white}Gratuit \textbullet\ Google Play\\Tuteur IA, annales, concours, résultats\par}
      \end{minipage}
    \end{tcolorbox}
  \end{minipage}\hfill
  \begin{minipage}[t]{0.487\linewidth}
    \begin{tcolorbox}[enhanced,colback=poppurple,colframe=popink,boxrule=2.2pt,arc=10pt,
        drop shadow=popink,left=11pt,right=11pt,top=10pt,bottom=10pt,height=3.1cm,valign=center]
      \tikz[baseline=-0.7ex]\node[circle,fill=white,draw=popink,line width=1.3pt,minimum size=1.6cm,inner sep=0]{\popdisplay\fontsize{24}{24}\selectfont\color{poppurple}+};%
      \hspace{8pt}\begin{minipage}[c]{0.6\linewidth}
        {\popdisplay\fontsize{11}{12.5}\selectfont\color{white}\MakeUppercase{Passe à OnBuch+}}\par\vspace{4pt}
        {\raggedright\popsemi\fontsize{8.4}{11}\selectfont\color{white}Tous les cours signés, Tuteur IA illimité, fiches PDF — onglet OnBuch+ de l'app\par}
      \end{minipage}
    \end{tcolorbox}
  \end{minipage}
  \vspace{8pt}
  \popcontenu
  \vfill
  \signatureonbuch
  \restoregeometry
  \newpage
}

% Rangée « Ce cours contient » (page de garde)
\newcommand{\poptile}[3]{% #1 couleur #2 gros texte #3 légende
  \begin{tcolorbox}[enhanced,colback=#1,colframe=popink,boxrule=2pt,arc=9pt,shadow={2.5pt}{-2.5pt}{0pt}{popink},
      left=6pt,right=6pt,top=9pt,bottom=9pt,halign=center,valign=center,height=1.75cm,width=\linewidth,nobeforeafter]
    {\popdisplay\fontsize{12.5}{13}\selectfont\color{white}\MakeUppercase{#2}}\par\vspace{3pt}
    {\centering\popsemi\fontsize{7.8}{9.5}\selectfont\color{white}#3\par}
  \end{tcolorbox}}
\newcommand{\popcontenu}{%
  \par\noindent{\popxboldls\fontsize{7.5}{7.5}\selectfont\color{popmuted}CE COURS SIGNÉ CONTIENT}\par\vspace{7pt}
  \noindent\begin{minipage}[t]{0.235\linewidth}\poptile{popblue}{Cours}{complet, illustré, pas à pas}\end{minipage}\hfill
  \begin{minipage}[t]{0.235\linewidth}\poptile{poporange}{Méthodes}{et exercices résolus}\end{minipage}\hfill
  \begin{minipage}[t]{0.235\linewidth}\poptile{poppink}{Exercices}{type examen + corrigés}\end{minipage}\hfill
  \begin{minipage}[t]{0.235\linewidth}\poptile{popgreen}{Fiche bilan}{l'essentiel à retenir}\end{minipage}\par}

% Sommaire
\newtcolorbox{sommaire}{enhanced,colback=white,colframe=popink,boxrule=2.2pt,arc=10pt,
  drop shadow=popink,left=18pt,right=18pt,top=13pt,bottom=13pt,before skip=6pt,after skip=20pt,
  before upper={\poppill{Sommaire}{poppurple}{white}\par\vspace{9pt}\popsemi\fontsize{10.6}{14.5}\selectfont\color{popink}}}

% Signature de fin de cours — poussée tout en bas de la dernière page
\newcommand{\findecours}{%
  \par\vfill\nopagebreak
  \begin{center}\popsquiggle{poporange}\end{center}\vspace{4pt}
  \signatureonbuch
}
\AtBeginDocument{\def\llbracket{\mathopen{[\mkern-3.5mu[}}\def\rrbracket{\mathclose{]\mkern-3.5mu]}}} % intervalles d'entiers (glyphe ⟦ absent de la police)
% Glyphes absents de Fira Math : redéfinis à partir de glyphes disponibles
\AtBeginDocument{%
  \def\setminus{\mathbin{\backslash}}\let\smallsetminus\setminus
  \def\vdots{\mathord{\vbox{\baselineskip3.2pt\lineskiplimit0pt\kern4pt\hbox{.}\hbox{.}\hbox{.}}}}%
  \def\ddots{\mathinner{\mkern1mu\raise6.5pt\hbox{.}\mkern2mu\raise3.6pt\hbox{.}\mkern2mu\raise0.7pt\hbox{.}\mkern1mu}}%
  \def\triangle{\mathord{\scalebox{0.9}{$\symup{\Delta}$}}}%
  \def\nabla{\mathord{\raisebox{\depth}{\scalebox{1}[-1]{$\symup{\Delta}$}}}}%
  \def\checkmark{\popcheck{popdark}}%
  \def\star{\mathbin{\ast}}\def\bigstar{\mathord{\ast}}%
  \def\diamond{\mathbin{\scalebox{0.6}{$\Diamond$}}}%
  \def\bigcup{\mathop{\mathchoice{\raisebox{-0.3em}{\scalebox{1.7}{$\cup$}}}{\scalebox{1.25}{$\cup$}}{\cup}{\cup}}\displaylimits}%
  \def\bigcap{\mathop{\mathchoice{\raisebox{-0.3em}{\scalebox{1.7}{$\cap$}}}{\scalebox{1.25}{$\cap$}}{\cap}{\cap}}\displaylimits}%
  \def\lesseqqgtr{\mathrel{\lessgtr}}\def\bigoplus{\mathop{\oplus}\displaylimits}%
}
\providecommand{\ssi}{si et seulement si} % abréviation fréquente
\AtBeginDocument{\providecommand{\wideparen}{\overparen}} % arc de cercle

%% FIGFIX-BEGIN v4 — correctifs globaux des figures (docs/onbuch-generation/tools/figfix)
\usepackage{adjustbox}\usepackage{etoolbox}
% toute figure TikZ/pgfplots plus large que la ligne est réduite à la largeur disponible
\BeforeBeginEnvironment{tikzpicture}{\begin{adjustbox}{max width=\linewidth}}
\AfterEndEnvironment{tikzpicture}{\end{adjustbox}}
% légendes pgfplots lisibles par-dessus les courbes
\pgfplotsset{every axis/.append style={legend style={fill=white,fill opacity=0.92,text opacity=1,draw=black!15,inner sep=3pt}}}
% étiquettes d'arêtes (syntaxe edge["..."]) avec fond blanc
\tikzset{every edge quotes/.append style={fill=white,inner sep=1.2pt}}
%% FIGFIX-END
\begin{document}

\pagegarde{Traduction et production de textes injonctifs liés aux médias et à la communication}{Allemand}{Tle A}{Chapitre 4 : Média et communication}{AL21}{2 h}{Cette leçon mixte entraîne les élèves de Tle A à lire, produire et traduire des textes injonctifs (affiches, règlements, consignes, slogans) sur le thème des médias et de la communication. Elle révise l'impératif, le discours indirect et le passif, et développe les techniques de thème et de version exigées au Baccalauréat.
\vspace{4pt}
\begin{multicols}{2}\small
\begin{itemize}
\item À la fin de cette leçon, je sais identifier les caractéristiques d'un texte injonctif (affiche, règlement, consigne, slogan) dans les médias.
\item Je sais employer le lexique allemand des médias et de la communication dans mes productions.
\item Je sais former et utiliser l'impératif allemand (du, ihr, Sie) pour rédiger des consignes et des slogans.
\item Je sais rapporter des propos au discours indirect avec le subjonctif I et les verbes de déclaration.
\item Je sais construire le passif (Vorgangspassiv) pour présenter une information sans nommer l'agent.
\item Je sais traduire de courtes phrases et un texte du français vers l'allemand (thème).
\end{itemize}\end{multicols}}

\begin{sommaire}
\begin{multicols}{2}
\begin{enumerate}
\item Texte support et compréhension
\item Lexique thématique : médias et communication
\item Grammaire 1 : l'impératif et les formes d'injonction
\item Grammaire 2 : discours indirect et passif
\item Méthode de traduction : thème et version
\item Production guidée : affiche, règlement, slogan
\item Activité d'intégration
\item Méthodes et exercices résolus
\item Exercices
\item Corrigés des exercices
\item Fiche bilan
\end{enumerate}
\end{multicols}
\end{sommaire}

\begin{objectifs}
\begin{itemize}
\item À la fin de cette leçon, je sais identifier les caractéristiques d'un texte injonctif (affiche, règlement, consigne, slogan) dans les médias.
\item Je sais employer le lexique allemand des médias et de la communication dans mes productions.
\item Je sais former et utiliser l'impératif allemand (du, ihr, Sie) pour rédiger des consignes et des slogans.
\item Je sais rapporter des propos au discours indirect avec le subjonctif I et les verbes de déclaration.
\item Je sais construire le passif (Vorgangspassiv) pour présenter une information sans nommer l'agent.
\item Je sais traduire de courtes phrases et un texte du français vers l'allemand (thème).
\item Je sais traduire un texte injonctif de l'allemand vers le français (version).
\item Je sais rédiger une affiche ou un règlement court sur les médias en allemand correct.
\item Je sais relire et corriger ma production (orthographe, majuscules, ordre des mots, cas).
\end{itemize}
\end{objectifs}

\begin{prerequis}
\begin{itemize}
\item Conjugaison des verbes forts et faibles au présent (classe de Première)
\item Ordre des mots : verbe en position 2, inversion après un complément en tête
\item Les cas allemands (nominatif, accusatif, datif) et les articles
\item Vocabulaire de base des médias vu au chapitre 4 (Zeitung, Fernsehen, Radio, Internet)
\item Notions de base sur le subjonctif I introduites en début d'année de Terminale
\end{itemize}
\end{prerequis}

\newpage

\coursec{Texte support et compréhension}

\courssub{Situation d'accroche et problématique}

Au lycée de Yaoundé, le club des jeunes reporters prépare une grande affiche pour la semaine des médias. Sur le carton, on lit déjà en lettres rouges : \all{Schaltet ein! Informiert euch!} « Allumez vos postes ! Informez-vous ! ». En Allemagne, en Autriche et en Suisse, les écoles organisent chaque année des \all{Medientage} (journées des médias) où les élèves rédigent slogans, consignes et règlements en allemand. Pour participer à un échange scolaire avec un lycée de Hambourg, les élèves camerounais doivent traduire le règlement de leur club et rédiger des slogans percutants.

\textbf{Problématique :} comment donner des ordres, rapporter des propos et mettre en valeur l'information sans nommer d'auteur ? C'est tout l'art des \cle{textes injonctifs} liés aux médias. Avant de produire, il faut d'abord apprendre à \emph{lire} ce type de texte : c'est l'objectif de cette première section.

\courssub{Lecture du texte support}

\begin{textebox}[Achtung, Medienkompetenz! — Affiche d'un lycée de Hambourg]
Liebe Schülerinnen und Schüler!

Unsere Schule veranstaltet vom 12. bis zum 16. Mai die große Medienwoche. Informiert euch täglich, aber überprüft die Quellen! Glaubt nicht alles, was im Internet steht! Vergleicht verschiedene Sender und Zeitungen, bevor ihr eine Nachricht teilt!

Fake News werden heute oft sehr schnell verbreitet – seid kritisch! Viele Meldungen in den sozialen Netzwerken werden von Robotern geschrieben und von Tausenden von Nutzern geteilt. Deshalb gilt bei uns die Regel: Jede Information wird von zwei unabhängigen Quellen bestätigt, bevor sie im Schülerforum veröffentlicht wird.

Der Direktor erklärte gestern, die Medienwoche werde am Montag um 9 Uhr feierlich eröffnet. Er fügte hinzu, alle Klassen sollten ein eigenes Plakat gestalten, und die besten Arbeiten würden in der Aula ausgestellt werden.

Das Programm: Am Dienstag besucht uns eine Journalistin des Norddeutschen Rundfunks. Am Mittwoch wird ein Workshop zum Thema „Netiquette“ angeboten. Am Donnerstag diskutieren wir über Werbung und ihre Tricks. Am Freitag wird der beste Slogan gewählt.

Beteiligt euch zahlreich! Bringt eure Handys mit, aber schaltet sie während der Workshops ab! Schaltet ein, denkt nach, diskutiert mit!

\textit{Das Organisationsteam der Medienwoche}
\end{textebox}

\textbf{Traduction de référence :} « Chers élèves ! Notre école organise du 12 au 16 mai la grande semaine des médias. Informez-vous chaque jour, mais vérifiez les sources ! Ne croyez pas tout ce qui se trouve sur Internet ! Comparez différentes chaînes et journaux avant de partager une nouvelle ! Les fausses informations sont souvent diffusées très vite aujourd'hui – soyez critiques ! Beaucoup de messages sur les réseaux sociaux sont écrits par des robots et partagés par des milliers d'utilisateurs. C'est pourquoi chez nous la règle vaut : toute information est confirmée par deux sources indépendantes avant d'être publiée sur le forum des élèves. Le directeur a expliqué hier que la semaine des médias serait inaugurée lundi à 9 heures. Il a ajouté que toutes les classes devraient concevoir leur propre affiche et que les meilleurs travaux seraient exposés dans la salle des fêtes. Le programme : mardi, une journaliste de la radio du Nord de l'Allemagne nous rendra visite. Mercredi, un atelier sur le thème de la « nétiquette » sera proposé. Jeudi, nous débattrons de la publicité et de ses astuces. Vendredi, le meilleur slogan sera élu. Participez nombreux ! Apportez vos portables, mais éteignez-les pendant les ateliers ! Allumez, réfléchissez, discutez ! »

\textbf{Glossaire :}

\begin{center}
\begin{tabularx}{\linewidth}{|l|X|}
\hline
\rowcolor{popblueL} \textbf{Allemand} & \textbf{Français} \\
\hline
die Medienkompetenz (nur Sg.) & l'esprit critique face aux médias \\
\hline
der Sender, die Sender & la chaîne, la station (de radio ou de télévision) \\
\hline
die Sendung, die Sendungen & l'émission \\
\hline
verbreiten (verbreitete, hat verbreitet) & diffuser, répandre \\
\hline
die Quelle, die Quellen & la source \\
\hline
überprüfen & vérifier, contrôler \\
\hline
die Fake News (Pl.) & les fausses informations \\
\hline
einschalten (schaltete ein, hat eingeschaltet) & allumer (un appareil) \\
\hline
abschalten (schaltete ab, hat abgeschaltet) & éteindre (un appareil) \\
\hline
die Nachricht, die Nachrichten & la nouvelle, l'information \\
\hline
die Meldung, die Meldungen & le message, la dépêche \\
\hline
der Nutzer, die Nutzer & l'utilisateur \\
\hline
veröffentlichen & publier \\
\hline
die Aula, die Aulen & la salle des fêtes (d'école) \\
\hline
sich beteiligen an (+ Dat.) & participer à \\
\hline
\end{tabularx}
\end{center}

\begin{attention}[Prononciation et orthographe]
\all{überprüfen} se prononce « u-beur-pru-fen » ; \all{Quelle} « kvé-le » ; \all{verbreiten} « féur-braï-ten ». Attention : \all{die Nachricht} signifie « la nouvelle / l'information », jamais « la nuit ». Et \all{die Sendung} (l'émission) ne se confond pas avec \all{der Sender} (la chaîne) : la chaîne \emph{diffuse} l'émission.
\end{attention}

\courssub{Repérage du type de texte}

\begin{definition}[Le texte injonctif]
Un \cle{texte injonctif} est un texte qui cherche à \textbf{faire agir} le lecteur : consigne, règlement, slogan, affiche, mode d'emploi, panneau. Il utilise typiquement :
\begin{itemize}
\item l'\textbf{impératif} : \all{Überprüft die Quellen!} « Vérifiez les sources ! » ;
\item l'\textbf{infinitif normatif} (panneaux, règlements) : \all{Handys abschalten!} « Éteindre les portables ! » ;
\item le \textbf{présent de vérité générale} : \all{Jede Information wird bestätigt.} « Toute information est confirmée. »
\end{itemize}
\end{definition}

Appliquons cette grille à notre affiche :

\begin{center}
\begin{tabularx}{\linewidth}{|X|X|}
\hline
\rowcolor{popblueL} \textbf{Critère} & \textbf{Réponse pour notre texte} \\
\hline
Émetteur & l'équipe d'organisation de la semaine des médias (un lycée de Hambourg) \\
\hline
Destinataire & les élèves (\all{Liebe Schülerinnen und Schüler!}) \\
\hline
Objectif & informer (programme), ordonner (règles), persuader (participer) \\
\hline
Ton & injonctif, chaleureux et mobilisateur \\
\hline
Marques typiques & impératifs à la 2\textsuperscript{e} pers. du pluriel, passifs, discours indirect \\
\hline
\end{tabularx}
\end{center}

\begin{methode}[Lire un texte injonctif en 3 étapes]
\begin{enumerate}
\item \textbf{Identifier l'émetteur et le destinataire} : qui parle ? à qui ? (ici : l'équipe d'organisation s'adresse aux élèves).
\item \textbf{Relever les verbes d'action et leur mode} : impératif, infinitif, présent passif ? (ici : \all{informiert}, \all{überprüft}, \all{wird verbreitet}...).
\item \textbf{Dégager l'intention} : informer, interdire, ordonner ou persuader ? Un même texte peut cumuler plusieurs intentions.
\end{enumerate}
\end{methode}

\courssub{Compréhension globale}

\begin{exemplebox}[Questions de compréhension globale]
\begin{enumerate}
\item \textbf{Qui s'adresse à qui ?} — \textit{Das Organisationsteam der Medienwoche wendet sich an die Schülerinnen und Schüler der Schule.} (L'équipe d'organisation s'adresse aux élèves.)
\item \textbf{Quelles actions sont demandées ?} — Les élèves doivent s'informer, vérifier les sources, comparer les médias, rester critiques, participer, apporter leurs portables mais les éteindre pendant les ateliers.
\item \textbf{Quel est l'événement annoncé ?} — La semaine des médias, du 12 au 16 mai, avec une inauguration le lundi à 9 heures.
\end{enumerate}
\end{exemplebox}

\courssub{Compréhension de détail : relevés grammaticaux}

\textbf{a) Les impératifs.} Le texte contient de nombreux impératifs à la 2\textsuperscript{e} personne du pluriel (\all{ihr}-forme), formés sur le radical du présent + \textbf{-t} (ou \textbf{-et} si le radical se termine en \emph{-t}, \emph{-d} ou en consonne + \emph{-m}/\emph{-n}) :

\begin{center}
\begin{tabular}{|l|l|l|}
\hline
\rowcolor{popblueL} \textbf{Impératif (texte)} & \textbf{Infinitif} & \textbf{Français} \\
\hline
Informiert euch! & sich informieren & Informez-vous ! \\
\hline
überprüft die Quellen! & überprüfen & vérifiez les sources ! \\
\hline
Glaubt nicht alles! & glauben & Ne croyez pas tout ! \\
\hline
Vergleicht die Sender! & vergleichen & Comparez les chaînes ! \\
\hline
seid kritisch! & sein (irrégulier) & soyez critiques ! \\
\hline
Beteiligt euch! & sich beteiligen & Participez ! \\
\hline
Bringt eure Handys mit! & mitbringen & Apportez vos portables ! \\
\hline
schaltet sie ab! & abschalten & éteignez-les ! \\
\hline
Schaltet ein! & einschalten & Allumez ! \\
\hline
denkt nach! & nachdenken & réfléchissez ! \\
\hline
diskutiert mit! & mitdiskutieren & discutez avec nous ! \\
\hline
\end{tabular}
\end{center}

Notez les verbes à \cle{particule séparable} : à l'impératif, la particule va en fin de phrase : \all{Schaltet ein!} « Allumez ! » ; \all{Bringt mit!} « Apportez ! ».

\textbf{b) Les tournures passives.} Le passif met l'information en avant sans nommer l'auteur :

\begin{center}
\begin{tabularx}{\linewidth}{|X|X|}
\hline
\rowcolor{popblueL} \textbf{Phrase du texte} & \textbf{Traduction} \\
\hline
Fake News werden oft schnell verbreitet. & Les fausses informations sont souvent diffusées vite. \\
\hline
Viele Meldungen werden von Robotern geschrieben. & Beaucoup de messages sont écrits par des robots. \\
\hline
Jede Information wird von zwei Quellen bestätigt. & Toute information est confirmée par deux sources. \\
\hline
Ein Workshop wird angeboten. & Un atelier est proposé. \\
\hline
Der beste Slogan wird gewählt. & Le meilleur slogan est élu. \\
\hline
\end{tabularx}
\end{center}

\textbf{c) Le discours indirect (Konjunktiv I).} Pour rapporter les propos du directeur sans les citer mot à mot, l'allemand utilise le subjonctif I ; quand la forme du subjonctif I se confond avec l'indicatif, on la remplace par le subjonctif II :

\all{Der Direktor erklärte, die Medienwoche werde am Montag eröffnet.} « Le directeur a expliqué que la semaine des médias serait inaugurée lundi. » (\all{werde} = Konjunktiv I de \all{werden}, 3\textsuperscript{e} pers. du sg.)

\all{Er fügte hinzu, alle Klassen sollten ein Plakat gestalten, und die besten Arbeiten würden in der Aula ausgestellt werden.} « Il a ajouté que toutes les classes devraient concevoir une affiche et que les meilleurs travaux seraient exposés dans la salle des fêtes. » (\all{sollten} et \all{würden ... ausgestellt werden} = formes de remplacement au Konjunktiv II, car le Konjunktiv I de la 3\textsuperscript{e} pers. du pluriel serait identique à l'indicatif).

Ces trois points (impératif, passif, discours indirect) seront approfondis dans les sections grammaire de cette leçon.

\begin{exoresolu}[Relevé grammatical guidé]
\textbf{Énoncé.} Dans la phrase \all{„Schaltet eure Handys ab, bevor der Workshop beginnt!“}, identifiez : 1) le mode du verbe \all{abschalten} ; 2) la place de la particule ; 3) le temps du verbe \all{beginnen}.
\tcblower
\textbf{Solution.} 1) \all{schaltet ... ab} est un \textbf{impératif} 2\textsuperscript{e} personne du pluriel (radical \all{schalt-} + \textbf{-et}, car le radical se termine en \emph{-t}). 2) La particule séparable \all{ab} est \textbf{détachée et rejetée en fin de proposition principale} : \all{Schaltet eure Handys ab!}. 3) \all{beginnt} est au \textbf{présent} dans la subordonnée introduite par \all{bevor} (avant que) ; le verbe conjugué est en \textbf{dernière position} de la subordonnée. Traduction : « Éteignez vos portables avant que l'atelier (ne) commence ! »
\end{exoresolu}

\begin{attention}[Pièges des francophones]
\begin{itemize}
\item Ne traduis pas \all{Schaltet ein!} par « Enclenchez ! » : avec un poste de radio ou de télévision, c'est « Allumez ! ».
\item À l'impératif pluriel, le français dit « Informez-vous ! » pour les verbes pronominaux ; l'allemand place le pronom réfléchi \emph{après} le verbe : \all{Informiert euch!} (et non \all{Euch informiert!}).
\item \all{werden + Partizip II} = passif ; \all{werden + Infinitiv} = futur. Compare : \all{Die Nachricht wird verbreitet.} « La nouvelle est diffusée » (passif) / \all{Die Nachricht wird uns überraschen.} « La nouvelle va nous surprendre » (futur).
\end{itemize}
\end{attention}

\courssub{Le champ lexical des médias dans le texte}

Classons le vocabulaire du texte en trois familles :

\begin{center}
\begin{tabularx}{\linewidth}{|X|X|X|}
\hline
\rowcolor{popblueL} \textbf{Presse (die Presse)} & \textbf{Audiovisuel (Rundfunk \& TV)} & \textbf{Numérique (Internet)} \\
\hline
die Zeitung (le journal) & der Sender (la chaîne) & das Internet (Internet) \\
\hline
die Quelle (la source) & die Sendung (l'émission) & die sozialen Netzwerke (les réseaux sociaux) \\
\hline
die Nachricht (la nouvelle) & der Rundfunk (la radio) & das Schülerforum (le forum des élèves) \\
\hline
der Journalist (le journaliste) & einschalten (allumer) & das Handy (le portable) \\
\hline
veröffentlichen (publier) & abschalten (éteindre) & teilen (partager), der Nutzer (l'utilisateur) \\
\hline
\end{tabularx}
\end{center}

La publicité (\all{die Werbung}) constitue une quatrième famille, évoquée dans le programme du jeudi.

\begin{popfigure}
\begin{tikzpicture}[mindmap, grow cyclic, every node/.style={concept, text=white}, root concept/.append style={concept color=popblue, font=\Large\bfseries}, level 1/.append style={level distance=3.4cm, sibling angle=90}, level 2/.append style={level distance=2.2cm, sibling angle=45, concept color=popmuted}]
\node[concept] {Medien}
  child[concept color=poporange] { node[concept] {Presse} child { node[concept, font=\small] {Zeitung} } child { node[concept, font=\small] {Zeitschrift} } }
  child[concept color=popgreen] { node[concept, align=center] {Rundfunk \& TV} child { node[concept, font=\small] {Sender} } child { node[concept, font=\small] {Sendung} } }
  child[concept color=poppurple] { node[concept, align=center] {Internet \& soziale Netzwerke} child { node[concept, font=\small] {Handy} } child { node[concept, font=\small] {Forum} } }
  child[concept color=poppink] { node[concept] {Werbung} child { node[concept, font=\small] {Plakat} } child { node[concept, font=\small] {Slogan} } };
\end{tikzpicture}
\legende{Carte mentale : le champ lexical des médias en allemand}
\end{popfigure}

\courssub{Le schéma de la communication}

Tout texte médiatique repose sur un modèle simple : un émetteur envoie un message à un récepteur. Notre affiche elle-même suit ce schéma : l'équipe d'organisation (émetteur) adresse des consignes (message) aux élèves (récepteurs).

\begin{popfigure}
\begin{tikzpicture}[node distance=1.6cm, every node/.style={font=\small}]
\node[draw=popblue, fill=popblueL, rounded corners, thick, minimum width=3cm, minimum height=1.1cm, align=center] (s) {\textbf{der Sender}\\ l'émetteur};
\node[draw=poporange, fill=poporangeL, rounded corners, thick, minimum width=3cm, minimum height=1.1cm, align=center, right=of s] (b) {\textbf{die Botschaft}\\ le message};
\node[draw=popgreen, fill=popgreenL, rounded corners, thick, minimum width=3cm, minimum height=1.1cm, align=center, right=of b] (e) {\textbf{der Empfänger}\\ le récepteur};
\draw[-{Stealth[length=3mm]}, very thick, popdark] (s) -- (b) node[midway, above, font=\footnotesize\itshape] {sendet};
\draw[-{Stealth[length=3mm]}, very thick, popdark] (b) -- (e) node[midway, above, font=\footnotesize\itshape] {erreicht};
\end{tikzpicture}
\legende{Le schéma de la communication : Sender – Botschaft – Empfänger}
\end{popfigure}

\begin{exercice}[À toi de jouer]
\begin{enumerate}
\item Relève dans le texte trois impératifs à particule séparable et traduis-les.
\item Transforme en phrase passive : \all{Die Schüler verbreiten die Nachricht.} (Les élèves diffusent la nouvelle.)
\item Qui est l'émetteur, quel est le message, qui est le récepteur dans l'affiche de la Medienwoche ?
\item Classe ces mots dans les trois familles (presse / audiovisuel / numérique) : \all{der Artikel, der Fernseher, das Passwort, die Schlagzeile, der Bildschirm}.
\end{enumerate}
\end{exercice}

\begin{corrige}[Exercice « À toi de jouer »]
\begin{enumerate}
\item \all{Schaltet ein!} « Allumez ! » ; \all{schaltet sie ab!} « éteignez-les ! » ; \all{Bringt eure Handys mit!} « Apportez vos portables ! » ; \all{diskutiert mit!} « discutez avec nous ! ».
\item \all{Die Nachricht wird von den Schülern verbreitet.} « La nouvelle est diffusée par les élèves. » (sujet au nominatif, auxiliaire \all{wird}, participe II \all{verbreitet} en fin de phrase, agent introduit par \all{von} + datif).
\item Émetteur : l'équipe d'organisation ; message : les consignes et le programme de la semaine des médias ; récepteurs : les élèves du lycée.
\item Presse : \all{der Artikel} (l'article), \all{die Schlagzeile} (le gros titre) ; audiovisuel : \all{der Fernseher} (le téléviseur) ; numérique : \all{das Passwort} (le mot de passe), \all{der Bildschirm} (l'écran).
\end{enumerate}
\end{corrige}

\begin{savaistu}[Les semaines des médias dans les pays germanophones]
En Allemagne, chaque \all{Bundesland} (Land fédéré) organise dans les écoles des semaines de l'éducation aux médias, les \all{Medienkompetenz-Wochen}, où les élèves apprennent à repérer les \all{Fake News} et à vérifier les sources. La Suisse alémanique, elle, publie des consignes de \all{Netiquette} (nétiquette, les règles de politesse sur Internet) en versions trilingues : allemand, français et italien. Au Cameroun, les clubs de jeunes reporters des lycées de Yaoundé et de Douala s'inspirent de ces modèles pour leurs journaux scolaires bilingues — une belle occasion de réutiliser les slogans allemands appris en classe !
\end{savaistu}

\begin{aretenir}[L'essentiel de la section]
\begin{itemize}
\item Un texte injonctif (affiche, règlement, slogan) vise à \textbf{faire agir} : il combine impératif, infinitif normatif et présent de vérité générale.
\item Lecture en 3 étapes : émetteur/destinataire $\rightarrow$ verbes d'action et modes $\rightarrow$ intention (informer, ordonner, persuader).
\item Le texte support illustre les trois outils grammaticaux du chapitre : l'\textbf{impératif} (\all{Überprüft die Quellen!}), le \textbf{passif} (\all{Fake News werden verbreitet}) et le \textbf{discours indirect au Konjunktiv I} (\all{die Medienwoche werde eröffnet}).
\item Le vocabulaire des médias se classe en quatre familles : presse, audiovisuel, numérique, publicité.
\end{itemize}
\end{aretenir}

\coursec{Lexique thématique : médias et communication}

Dans la section précédente, nous avons lu et commenté un texte de presse sur le rôle des médias dans la vie quotidienne des jeunes Camerounais. Nous y avons rencontré des mots comme \all{die Zeitung}, \all{die Nachrichten} ou \all{die sozialen Netzwerke}. Pour pouvoir lire, traduire et produire des textes sur les médias — affiches, règlements, articles, slogans —, il faut d'abord maîtriser ce vocabulaire de manière précise : article, pluriel, prononciation, emploi. C'est l'objectif de cette section : construire, champ par champ, un lexique solide et actif, que vous réinvestirez ensuite dans les exercices de thème, de version et de production d'affiches.

\courssub{Observer d'abord : un court texte de presse}

Lisons d'abord ce court article original. Soulignez mentalement tous les mots qui appartiennent au champ lexical des médias.

\begin{textebox}[Ein Blick in die Medienwelt Kameruns]
Jeden Morgen erscheint in Yaoundé die neue Ausgabe der Zeitung. Schon um sechs Uhr verkaufen junge Händler die Zeitungen an den Kreuzungen der Stadt. Die Schlagzeile auf der ersten Seite informiert die Leser über das wichtigste Ereignis des Tages. Ein Journalist der Redaktion hat den Artikel am Abend vorher geschrieben und die Quellen sorgfältig geprüft.

Aber die gedruckte Presse hat heute starke Konkurrenz: Viele Kameruner schalten morgens das Radio ein, um die Nachrichten zu hören. Der Moderator einer bekannten Sendung berichtet jeden Tag über Politik, Sport und Kultur. Am Abend überträgt das Fernsehen die wichtigsten Ereignisse, und viele Familien schauen gemeinsam die Nachrichten an, bevor sie den Fernseher wieder abschalten.

Die jungen Leute informieren sich meistens im Internet. Sie besuchen Websites, lesen Beiträge in den sozialen Netzwerken, posten Fotos, teilen Videos und liken die Seiten ihrer Lieblingsstars. Doch Vorsicht: Nicht alles, was man online liest, ist wahr. Fake News verbreiten sich sehr schnell, besonders vor Wahlen. Deshalb ist es wichtig, die Quelle einer Information zu prüfen, bevor man sie kommentiert oder weiterschickt. Auch die Werbung im Internet ist nicht immer ehrlich. Wer gut informiert sein will, muss also kritisch denken — ob in der Zeitung, im Radio oder im Netz.
\end{textebox}

\textbf{Glossaire :}

\begin{center}
\begin{tabularx}{\linewidth}{|X|X|}
\hline
\rowcolor{popblueL} \textbf{Deutsch} & \textbf{Français} \\
\hline
die Ausgabe, -n & l'édition, le numéro (d'un journal) \\
\hline
der Händler, - & le vendeur, le marchand \\
\hline
die Kreuzung, -en & le carrefour \\
\hline
das Ereignis, -se & l'événement \\
\hline
die Quelle, -n & la source \\
\hline
die Konkurrenz & la concurrence \\
\hline
die Wahl, -en & l'élection \\
\hline
weiterschicken & transmettre, faire suivre \\
\hline
ehrlich & honnête \\
\hline
kritisch denken & réfléchir de manière critique \\
\hline
\end{tabularx}
\end{center}

\textbf{Questions de compréhension :}
\begin{enumerate}
\item À quelle heure vend-on les journaux à Yaoundé ?
\item Que fait le journaliste avant la parution de l'article ?
\item Pourquoi les jeunes doivent-ils être prudents sur Internet ?
\item Citez trois médias différents mentionnés dans le texte.
\end{enumerate}

\courssub{Champ 1 : la presse écrite — die Presse}

\begin{definition}[die Schlagzeile et die Überschrift]
\all{Die Schlagzeile, -n} désigne le \cle{titre d'un article de journal}, celui qui « frappe » le lecteur (du verbe \all{schlagen} = frapper). \all{Die Überschrift, -en} est le terme général pour tout titre : titre d'un texte, d'un chapitre, d'une affiche. Toute \all{Schlagzeile} est une \all{Überschrift}, mais l'inverse n'est pas vrai.
\end{definition}

Voici le vocabulaire essentiel de la presse écrite :

\begin{center}
\begin{tabularx}{\linewidth}{|X|X|X|}
\hline
\rowcolor{popblueL} \textbf{Deutsch} & \textbf{Français} & \textbf{Exemple} \\
\hline
die Zeitung, -en & le journal & \all{Ich lese jeden Tag die Zeitung.} \\
\hline
die Zeitschrift, -en & la revue, le magazine & \all{Diese Zeitschrift erscheint monatlich.} \\
\hline
der Artikel, - & l'article & \all{Der Artikel ist sehr interessant.} \\
\hline
die Schlagzeile, -n & le titre (de journal) & \all{Die Schlagzeile steht auf der ersten Seite.} \\
\hline
der Journalist, -en / die Journalistin, -nen & le/la journaliste & \all{Die Journalistin arbeitet bei CRTV.} \\
\hline
die Redaktion, -en & la rédaction & \all{Er schickt den Text an die Redaktion.} \\
\hline
erscheinen (erscheint, erschien, ist erschienen) & paraître & \all{Die Zeitung erscheint täglich.} \\
\hline
\end{tabularx}
\end{center}

\begin{exemplebox}[Exemples traduits]
\all{Der Artikel erscheint morgen.} « L'article paraît demain. »

\all{Die Zeitschrift erscheint einmal pro Woche.} « Le magazine paraît une fois par semaine. »

\all{Der Journalist hat die Quelle geprüft.} « Le journaliste a vérifié la source. »

\all{Die Redaktion druckt die Zeitung in Douala.} « La rédaction imprime le journal à Douala. »
\end{exemplebox}

\begin{attention}[erscheinen : auxiliaire « sein »]
Le verbe \all{erscheinen} se conjugue au \cle{Perfekt} avec l'auxiliaire \all{sein} : \all{Die Zeitung ist gestern erschienen.} « Le journal est paru hier. » Ne dites jamais \all{hat erschienen}. Attention aussi au genre : on dit \all{der Artikel} (masculin), contrairement au français « l'article » qui peut induire en erreur. Notez enfin la déclinaison de \all{der Journalist} : c'est un nom faible — \all{den Journalisten} à l'accusatif, \all{dem Journalisten} au datif.
\end{attention}

\courssub{Champ 2 : l'audiovisuel — Radio und Fernsehen}

\begin{center}
\begin{tabularx}{\linewidth}{|X|X|X|}
\hline
\rowcolor{popblueL} \textbf{Deutsch} & \textbf{Français} & \textbf{Exemple} \\
\hline
das Fernsehen & la télévision & \all{Ich sehe abends fern.} \\
\hline
der Fernseher, - & le téléviseur (l'appareil) & \all{Der Fernseher ist kaputt.} \\
\hline
das Radio, -s & la radio & \all{Sie hört morgens Radio.} \\
\hline
der Sender, - & la chaîne, la station & \all{Dieser Sender sendet auf Französisch.} \\
\hline
die Sendung, -en & l'émission & \all{Die Sendung beginnt um 20 Uhr.} \\
\hline
die Nachrichten (pl.) & les informations, les nouvelles & \all{Die Nachrichten kommen um 19 Uhr.} \\
\hline
der Moderator, -en / die Moderatorin, -nen & le/la présentateur(-trice) & \all{Der Moderator stellt Fragen.} \\
\hline
der Zuschauer, - / die Zuschauerin, -nen & le/la téléspectateur(-trice) & \all{Die Zuschauer schalten ein.} \\
\hline
senden / übertragen (überträgt, übertrug, hat übertragen) & diffuser / retransmettre & \all{Das Fernsehen überträgt das Spiel live.} \\
\hline
einschalten / abschalten & allumer / éteindre & \all{Schalte bitte das Radio ein!} \\
\hline
\end{tabularx}
\end{center}

\begin{propriete}[Distinction importante : das Fernsehen / der Fernseher]
\all{Das Fernsehen} désigne la \cle{télévision} comme média (les programmes), tandis que \all{der Fernseher} désigne l'\cle{appareil}. Comparez : \all{Ich sehe fern.} « Je regarde la télévision. » / \all{Schalte den Fernseher aus!} « Éteins le téléviseur ! » Le français « la télé » recouvre les deux sens ; l'allemand les distingue. Remarquez aussi le verbe séparable \all{fernsehen} : \all{ich sehe fern, ich sah fern, ich habe ferngesehen}.
\end{propriete}

\begin{exemplebox}[Exemples traduits]
\all{Der Sender überträgt das Fußballspiel live aus Yaoundé.} « La chaîne retransmet le match de football en direct de Yaoundé. »

\all{Meine Mutter schaltet jeden Abend das Radio ein.} « Ma mère allume la radio chaque soir. »

\all{Die Moderatorin begrüßt die Zuschauer.} « La présentatrice salue les téléspectateurs. »
\end{exemplebox}

\courssub{Champ 3 : le numérique — die digitalen Medien}

\begin{center}
\begin{tabularx}{\linewidth}{|X|X|X|}
\hline
\rowcolor{popblueL} \textbf{Deutsch} & \textbf{Français} & \textbf{Exemple} \\
\hline
das Internet & l'Internet & \all{Ich surfe im Internet.} \\
\hline
die Website, -s & le site web & \all{Die Website der Schule ist neu.} \\
\hline
das soziale Netzwerk, -e & le réseau social & \all{Facebook ist ein soziales Netzwerk.} \\
\hline
der Beitrag, -e & la publication, le post & \all{Sein Beitrag hat viele Likes.} \\
\hline
posten & publier, poster & \all{Sie postet jeden Tag Fotos.} \\
\hline
teilen & partager & \all{Er teilt das Video mit Freunden.} \\
\hline
liken & liker & \all{Viele Schüler liken die Seite.} \\
\hline
herunterladen (lädt herunter, lud herunter, hat heruntergeladen) & télécharger & \all{Ich lade die App herunter.} \\
\hline
das Passwort, -e & le mot de passe & \all{Vergiss dein Passwort nicht!} \\
\hline
die Fake News (pl.) & les fausses informations & \all{Fake News verbreiten sich schnell.} \\
\hline
\end{tabularx}
\end{center}

\begin{exemplebox}[Exemples traduits]
\all{Sie hat die Nachricht auf Facebook geteilt.} « Elle a partagé l'information sur Facebook. »

\all{Die Schüler laden eine Lern-App herunter.} « Les élèves téléchargent une application d'apprentissage. »

\all{Prüfe die Quelle, bevor du den Beitrag teilst!} « Vérifie la source avant de partager la publication ! »
\end{exemplebox}

\begin{attention}[Verbes anglicisés : posten, liken]
Les verbes \all{posten} et \all{liken} sont des verbes \cle{faibles réguliers} : \all{ich poste, ich postete, ich habe gepostet} ; \all{ich like, ich likte, ich habe gelikt}. Ne dites pas \all{gepostten} ni \all{geliken}. En revanche, \all{herunterladen} est un verbe fort à particule séparable : \all{ich lade herunter, ich lud herunter, ich habe heruntergeladen}. Au Perfekt, le \all{ge-} se place entre la particule et le radical : \all{herunter\textbf{ge}laden}.
\end{attention}

\courssub{Champ 4 : communiquer et informer — kommunizieren und informieren}

\begin{center}
\begin{tabularx}{\linewidth}{|X|X|X|}
\hline
\rowcolor{popblueL} \textbf{Deutsch} & \textbf{Français} & \textbf{Exemple} \\
\hline
sich informieren (über + Akk.) & s'informer (sur) & \all{Ich informiere mich über die Wahl.} \\
\hline
berichten (über + Akk. / von + Dat.) & rapporter, faire un reportage & \all{Die Zeitung berichtet über den Markt.} \\
\hline
verbreiten & diffuser, répandre & \all{Er verbreitet Fake News.} \\
\hline
veröffentlichen & publier & \all{Die Redaktion veröffentlicht den Artikel.} \\
\hline
kommentieren & commenter & \all{Viele Leser kommentieren den Beitrag.} \\
\hline
die Quelle, -n & la source & \all{Die Quelle ist zuverlässig.} \\
\hline
die Meinung, -en & l'avis, l'opinion & \all{Meiner Meinung nach ist das falsch.} \\
\hline
die Werbung & la publicité & \all{Im Fernsehen kommt viel Werbung.} \\
\hline
\end{tabularx}
\end{center}

\begin{propriete}[sich informieren : verbe réfléchi + préposition]
\all{Sich informieren} est un verbe \cle{réfléchi} suivi de la préposition \all{über} + accusatif : \all{Wir informieren uns über die aktuellen Nachrichten.} « Nous nous informons sur les actualités. » Ne traduisez pas « s'informer de » par \all{von} : la préposition correcte est \all{über}. Autre construction utile : \all{Meiner Meinung nach...} « À mon avis... » (avec le datif, la préposition \all{nach} se place après le nom).
\end{propriete}

\begin{attention}[Faux ami : kommentieren / die Bemerkung]
\all{Kommentieren} signifie bien « commenter » (un article, un événement) : \all{Der Journalist kommentiert das Spiel.} Mais « un commentaire » au sens de \emph{remarque} se dit \all{die Bemerkung, -en} : \all{Der Lehrer macht eine Bemerkung.} « Le professeur fait une remarque. » Le nom correspondant à \all{kommentieren} est \all{der Kommentar, -e} (commentaire publié, éditorial).
\end{attention}

\courssub{Verbes à particule séparable fréquents}

\begin{propriete}[Rappel : la particule séparable]
Au présent et à l'impératif, la particule séparable va \cle{en fin de phrase} : \all{Ich schalte das Radio ein.} Au Perfekt, elle s'intègre au participe : \all{eingeschaltet}. Dans une subordonnée, le verbe se ressoude : \all{..., weil ich das Radio einschalte.}
\end{propriete}

\begin{center}
\begin{tabularx}{\linewidth}{|X|X|X|}
\hline
\rowcolor{popblueL} \textbf{Verbe} & \textbf{Français} & \textbf{Exemple} \\
\hline
einschalten & allumer & \all{Schalte bitte den Fernseher ein!} \\
\hline
abschalten & éteindre & \all{Er schaltet das Handy ab.} \\
\hline
aufnehmen (nimmt auf, nahm auf, hat aufgenommen) & enregistrer & \all{Der Sender nimmt die Sendung auf.} \\
\hline
mitteilen & communiquer, faire savoir & \all{Der Direktor teilt uns die Nachricht mit.} \\
\hline
nachschlagen (schlägt nach, schlug nach, hat nachgeschlagen) & consulter (un dictionnaire) & \all{Schlag das Wort im Wörterbuch nach!} \\
\hline
\end{tabularx}
\end{center}

\begin{exemplebox}[Exemples traduits]
\all{Die Redaktion teilt den Lesern die Ergebnisse mit.} « La rédaction communique les résultats aux lecteurs. »

\all{Ich habe das Interview gestern aufgenommen.} « J'ai enregistré l'interview hier. »

\all{Wenn du ein Wort nicht kennst, schlag es nach!} « Si tu ne connais pas un mot, consulte-le (dans le dictionnaire) ! »
\end{exemplebox}

\courssub{Deux pièges majeurs : die Medien et les faux amis}

\begin{attention}[« die Medien » est toujours pluriel]
En allemand, \all{die Medien} est un \cle{pluriel} : \all{Die Medien spielen eine wichtige Rolle.} « Les médias jouent un rôle important. » Le singulier est \all{das Medium, die Medien} : \all{Das Fernsehen ist ein Medium.} On ne dit \textbf{jamais} \all{das Media}. Le français « le média » au singulier ne doit pas vous tromper.
\end{attention}

\begin{attention}[Faux amis : die Nachricht et aktuell]
\begin{itemize}
\item \all{Die Nachricht, -en} signifie « une information, une nouvelle (transmise), un message » : \all{Ich habe eine Nachricht von meinem Bruder bekommen.} « J'ai reçu un message de mon frère. » Au pluriel, \all{die Nachrichten} = « les informations » (journal télévisé). Pour « une nouvelle » au sens d'événement personnel, l'allemand utilise plutôt \all{die Neuigkeit, -en}.
\item \all{Aktuell} signifie « actuel, d'actualité » : \all{die aktuellen Nachrichten} « les nouvelles actuelles ». Mais « actuel » au sens de « réel, véritable » se traduit par \all{wirklich} ou \all{tatsächlich} : \all{Die wirkliche Zahl ist höher.} « Le chiffre réel est plus élevé. »
\end{itemize}
\end{attention}

\courssub{Carte mentale du lexique}

\begin{popfigure}
\begin{center}
\begin{tikzpicture}[
  root/.style={rectangle, rounded corners, draw=popink, fill=popblueL, font=\bfseries, align=center, minimum width=3.4cm, minimum height=0.9cm},
  br/.style={rectangle, rounded corners, draw=popink, fill=poporangeL, font=\bfseries\small, align=center, minimum width=2.9cm},
  mot/.style={font=\small, align=left, text width=3.1cm}
]
\node[root, align=center] (r){Medien und\\Kommunikation};
\node[br, above left=1.1cm and 0.2cm of r] (p) {die Presse};
\node[br, above right=1.1cm and 0.2cm of r] (a) {das Audiovisuelle};
\node[br, below left=1.1cm and 0.2cm of r] (n) {das Internet};
\node[br, below right=1.1cm and 0.2cm of r] (k) {informieren};
\draw[popline, thick] (r) -- (p);
\draw[popline, thick] (r) -- (a);
\draw[popline, thick] (r) -- (n);
\draw[popline, thick] (r) -- (k);
\node[mot, above=0.15cm of p, align=center]{die Zeitung\\der Artikel\\die Schlagzeile\\der Journalist\\die Redaktion\\erscheinen};
\node[mot, above=0.15cm of a, align=center]{das Fernsehen\\das Radio\\der Sender\\die Sendung\\die Nachrichten\\senden / übertragen};
\node[mot, below=0.15cm of n, align=center]{die Website\\das Netzwerk\\der Beitrag\\posten, teilen\\herunterladen\\das Passwort};
\node[mot, below=0.15cm of k, align=center]{sich informieren\\berichten\\verbreiten\\veröffentlichen\\kommentieren\\die Quelle};
\end{tikzpicture}
\end{center}
\legende{Carte mentale : les quatre champs lexicaux des médias et de la communication.}
\end{popfigure}

\begin{popfigure}
\begin{center}
\begin{tikzpicture}[
  col/.style={rectangle, rounded corners, draw=popink, fill=popgreenL, font=\bfseries\small, align=center, minimum width=3.2cm, minimum height=0.8cm},
  cell/.style={font=\small, align=center, text width=3.2cm}
]
\node[col] (c1) {Presse};
\node[col, right=0.4cm of c1] (c2) {Audiovisuel};
\node[col, right=0.4cm of c2] (c3) {Numérique};
\node[col, right=0.4cm of c3] (c4) {Communiquer};
\node[cell, below=0.2cm of c1, align=center]{die Zeitung\\die Zeitschrift\\der Artikel\\die Schlagzeile\\der Journalist\\die Redaktion};
\node[cell, below=0.2cm of c2, align=center]{das Fernsehen\\das Radio\\der Sender\\die Sendung\\die Nachrichten\\der Moderator};
\node[cell, below=0.2cm of c3, align=center]{das Internet\\die Website\\der Beitrag\\das Passwort\\die Fake News\\das Netzwerk};
\node[cell, below=0.2cm of c4, align=center]{berichten\\verbreiten\\veröffentlichen\\kommentieren\\die Quelle\\die Meinung};
\end{tikzpicture}
\end{center}
\legende{Tableau à quatre colonnes : six mots-clés par champ lexical.}
\end{popfigure}

\begin{savaistu}[D'où vient le mot « Zeitung » ?]
\all{Die Zeitung} vient de \all{die Zeit} (« le temps ») : le journal est littéralement « la chose du temps », celle qui rapporte ce qui se passe \emph{maintenant}. Le même mot \all{Zeit} se retrouve dans \all{die Zeitschrift} (la revue, « l'écrit du temps »). Au Cameroun, dans les exercices de traduction, on rend souvent le titre du quotidien national \emph{Cameroon Tribune} par \all{Kameruner Tribune} — un bon entraînement pour le thème !
\end{savaistu}

\begin{exoresolu}[Complétez avec le mot juste]
Complétez les phrases avec : \all{Schlagzeile, Sender, Beitrag, Quelle, erscheint, abschalten}.

\begin{enumerate}
\item Die Zeitung \ldots\ jeden Morgen um sechs Uhr.
\item Die \ldots\ auf der ersten Seite ist sehr groß.
\item Dieser \ldots\ überträgt die Nachrichten auf Englisch.
\item Prüfe immer die \ldots\ einer Information!
\item Sein \ldots\ auf Facebook hat 500 Likes.
\item Bitte \ldots\ den Fernseher, es ist Mitternacht!
\end{enumerate}

\tcblower

\textbf{Solution détaillée :}

\begin{enumerate}
\item \all{Die Zeitung \textbf{erscheint} jeden Morgen um sechs Uhr.} — verbe \all{erscheinen} (paraître), 3\up{e} personne du singulier.
\item \all{Die \textbf{Schlagzeile} auf der ersten Seite ist sehr groß.} — le titre d'un journal = \all{die Schlagzeile}.
\item \all{Dieser \textbf{Sender} überträgt die Nachrichten auf Englisch.} — \all{dieser} impose un nom masculin : \all{der Sender} (la chaîne).
\item \all{Prüfe immer die \textbf{Quelle} einer Information!} — on vérifie la \emph{source}.
\item \all{Sein \textbf{Beitrag} auf Facebook hat 500 Likes.} — une publication sur un réseau social = \all{der Beitrag}.
\item Attention, piège ! La particule séparable se détache à l'impératif : la forme correcte est \all{Bitte \textbf{schalte} den Fernseher \textbf{ab}, es ist Mitternacht!} « Éteins le téléviseur, il est minuit ! » On ne peut donc pas écrire \all{Bitte abschalten den Fernseher}.
\end{enumerate}
\end{exoresolu}

\begin{exercice}[À vous de jouer]
\begin{enumerate}
\item Traduisez en allemand : « Le journaliste publie l'article demain. » / « Elle a partagé la vidéo sur les réseaux sociaux. » / « Allume la radio, s'il te plaît ! »
\item Trouvez l'intrus : \all{die Zeitung — die Sendung — die Zeitschrift — die Redaktion}.
\item Complétez : \all{Ich informiere mich \ldots\ die aktuellen Nachrichten.} (quelle préposition ?)
\end{enumerate}
\end{exercice}

\begin{corrige}[Exercice « À vous de jouer »]
\begin{enumerate}
\item \all{Der Journalist veröffentlicht den Artikel morgen.} / \all{Sie hat das Video in den sozialen Netzwerken geteilt.} / \all{Schalte bitte das Radio ein!} (particule \all{ein} en fin de phrase).
\item L'intrus est \all{die Sendung} (l'émission) : elle appartient à l'audiovisuel, les trois autres à la presse écrite.
\item \all{Ich informiere mich \textbf{über} die aktuellen Nachrichten.} — \all{sich informieren über} + accusatif.
\end{enumerate}
\end{corrige}

\begin{aretenir}[L'essentiel du lexique]
\begin{itemize}
\item Quatre champs : la presse (\all{die Zeitung, der Artikel, erscheinen}), l'audiovisuel (\all{das Fernsehen, der Sender, die Sendung}), le numérique (\all{das Internet, der Beitrag, posten, teilen}), communiquer (\all{sich informieren über, berichten, verbreiten, die Quelle}).
\item \all{Die Medien} est toujours pluriel ; le singulier est \all{das Medium}.
\item Faux amis : \all{die Nachricht} (information/message), \all{aktuell} (d'actualité, pas « réel »), \all{kommentieren} (commenter) mais \all{die Bemerkung} (remarque).
\item Verbes à particule : la particule va en fin de phrase au présent et à l'impératif : \all{Schalte das Radio ein!}
\end{itemize}
\end{aretenir}

Ce lexique est désormais votre boîte à outils. Dans la section suivante, nous verrons comment ces mots s'organisent dans des textes injonctifs grâce à l'impératif et aux autres formes d'injonction.

\coursec{Grammaire 1 : l'impératif et les formes d'injonction}

Dans la section précédente, nous avons constitué le lexique des médias et de la communication : \all{die Zeitung}, \all{der Artikel}, \all{die Quelle}, \all{die Nachrichten}, \all{das Handy}... Or, ce vocabulaire ne vit jamais seul : dans une affiche, un règlement de salle multimédia, un slogan publicitaire ou une consigne de modérateur sur un réseau social, il est presque toujours accompagné d'un verbe qui \emph{ordonne}, \emph{conseille} ou \emph{invite}. C'est précisément le rôle de l'\cle{impératif} et des autres \cle{formes d'injonction}. Les maîtriser est indispensable pour traduire les textes injonctifs au programme de Terminale : affiches, règlements, consignes et slogans.

\courssub{Observation : un règlement de salle multimédia}

Lisons d'abord ce court règlement affiché dans la salle multimédia d'un lycée de Yaoundé, en partenariat avec un établissement allemand. Soulignons mentalement la place et la forme des verbes.

\begin{textebox}[Règlement de la salle multimédia — \all{Regeln für den Medienraum}]
\all{Liebe Schülerinnen und Schüler! Willkommen im Medienraum unserer Schule! Lest die Regeln aufmerksam und befolgt sie! Schaltet eure Handys aus, bevor ihr den Raum betretet! Benutzt die Computer nur für die Recherche und öffnet keine privaten Seiten! Überprüft jede Information, bevor ihr sie teilt: Vergleicht die Quellen und glaubt nicht alles, was ihr im Internet lest! Seid kritisch gegenüber den Nachrichten in den sozialen Netzwerken! Sprich leise mit deinem Nachbarn und störe die anderen nicht! Iss und trink nichts im Medienraum! Räumt euren Platz auf, bevor ihr geht! Bei Problemen fragt den Betreuer — er hilft euch gern. Und denkt daran: Informiert euch täglich, aber informiert euch richtig!}\par
\emph{Glossaire :} \all{befolgen} = respecter (une règle) ; \all{ausschalten} (trennbar) = éteindre ; \all{die Recherche, -n} = la recherche ; \all{überprüfen} = vérifier ; \all{die Quelle, -n} = la source ; \all{teilen} = partager ; \all{stören} = déranger ; \all{aufräumen} (trennbar) = ranger ; \all{der Betreuer, -} = le responsable, l'encadrant.
\end{textebox}

\textbf{Questions d'observation :}
\begin{enumerate}
\item Où se trouve le verbe dans chaque consigne ? (Réponse : en \cle{position 1}.)
\item Quelles sont les trois « personnes » auxquelles on s'adresse dans ce texte ? (Réponse : \all{du} — \all{Sprich leise!} ; \all{ihr} — \all{Lest! Schaltet aus!} ; le \all{Sie} poli apparaîtrait dans un règlement pour adultes.)
\item Que devient la particule de \all{ausschalten} dans \all{Schaltet eure Handys aus!} ? (Réponse : elle se détache et va en fin de phrase.)
\end{enumerate}

\courssub{Formation de l'impératif : les trois formes}

L'allemand possède \textbf{trois formes d'impératif}, car il distingue le tutoiement singulier (\all{du}), le tutoiement pluriel (\all{ihr}) et le vouvoiement (\all{Sie}). Le français, lui, n'en a que deux (tu / nous-vous) : c'est une première différence à bien intégrer. En thème, la première question à se poser est donc toujours : « À qui parle-t-on ? »

\begin{propriete}[Les trois formes de l'impératif]
\begin{itemize}
\item \textbf{Impératif \all{du}} : \cle{radical du présent} (+ \all{-e} optionnelle, sauf pour les verbes à changement vocalique). \all{Komm!} \all{Lies!} \all{Überprüf(e)!}
\item \textbf{Impératif \all{ihr}} : \cle{forme du présent de la 2\up{e} personne du pluriel}, sans le pronom \all{ihr}. \all{Kommt!} \all{Lest!} \all{Überprüft!}
\item \textbf{Impératif \all{Sie}} : \cle{infinitif + Sie} (le pronom est obligatoire). \all{Kommen Sie!} \all{Lesen Sie!} \all{Überprüfen Sie!}
\end{itemize}
Dans les trois cas, le verbe est \textbf{toujours en position 1} et la phrase se termine par un point d'exclamation (ou un point dans les consignes écrites neutres).
\end{propriete}

\begin{center}
\begin{tabular}{|l|l|l|l|}\hline
\rowcolor{popblueL} \textbf{Infinitif} & \textbf{du} & \textbf{ihr} & \textbf{Sie} \\ \hline
\all{kommen} (venir) & \all{Komm!} & \all{Kommt!} & \all{Kommen Sie!} \\ \hline
\all{lesen} (lire) & \all{Lies!} & \all{Lest!} & \all{Lesen Sie!} \\ \hline
\all{sein} (être) & \all{Sei!} & \all{Seid!} & \all{Seien Sie!} \\ \hline
\all{ausschalten} (éteindre) & \all{Schalt(e) aus!} & \all{Schaltet aus!} & \all{Schalten Sie aus!} \\ \hline
\all{überprüfen} (vérifier) & \all{Überprüf(e)!} & \all{Überprüft!} & \all{Überprüfen Sie!} \\ \hline
\all{zuhören} (écouter) & \all{Hör(e) zu!} & \all{Hört zu!} & \all{Hören Sie zu!} \\ \hline
\end{tabular}
\end{center}

\begin{popfigure}
\begin{tikzpicture}[font=\small]
\node[draw=popblue, fill=popblueL, rounded corners, align=center, minimum width=3.4cm] (titre) at (0,0) {\textbf{IMPÉRATIF}\\ verbe en position 1};
\node[draw=popgreen, fill=popgreenL, rounded corners, align=center, minimum width=4.2cm] (du) at (-5.4,-2.6) {\textbf{du} : radical (+ e)\\ \all{Komm! Lies! Sei!}\\ \all{Schalt(e) aus!}};
\node[draw=poporange, fill=poporangeL, rounded corners, align=center, minimum width=4.2cm] (ihr) at (0,-2.6) {\textbf{ihr} : présent sans \all{ihr}\\ \all{Kommt! Lest! Seid!}\\ \all{Schaltet aus!}};
\node[draw=poppurple, fill=poppurpleL, rounded corners, align=center, minimum width=4.2cm] (sie) at (5.4,-2.6) {\textbf{Sie} : infinitif + Sie\\ \all{Kommen Sie! Lesen Sie!}\\ \all{Seien Sie! Schalten Sie aus!}};
\draw[-{Stealth}, thick, popblue] (titre) -- (du);
\draw[-{Stealth}, thick, popblue] (titre) -- (ihr);
\draw[-{Stealth}, thick, popblue] (titre) -- (sie);
\end{tikzpicture}
\legende{Les trois formes de l'impératif allemand : du / ihr / Sie}
\end{popfigure}

\courssub{Les verbes à changement vocalique et le verbe \all{sein}}

\begin{propriete}[Changement vocalique à l'impératif \all{du}]
Les verbes forts qui changent \cle{e en i ou ie} au présent (\all{du liest}, \all{du gibst}) conservent ce changement à l'impératif \all{du}, \textbf{sans jamais ajouter de -e} :
\begin{itemize}
\item \all{lesen} → \all{Lies den Artikel!} (« Lis l'article ! »)
\item \all{geben} → \all{Gib mir die Zeitung!} (« Donne-moi le journal ! »)
\item \all{nehmen} → \all{Nimm Platz!} (« Assieds-toi ! », litt. « Prends place ! »)
\item \all{sprechen} → \all{Sprich leise!} (« Parle doucement ! »)
\item \all{helfen} → \all{Hilf mir!} (« Aide-moi ! »)
\end{itemize}
En revanche, les verbes en \cle{a → ä} (\all{fahren} : \all{du fährst}) \textbf{ne changent pas} à l'impératif : \all{Fahr!} et jamais \all{Fähr!}.
\end{propriete}

\begin{attention}[Le piège du tréma]
Les francophones qui connaissent la conjugaison au présent (\all{du fährst}, \all{du lädst}) ont tendance à écrire \all{Fähr!} ou \all{Läd!}. C'est faux : le changement \all{a → ä} est une simple \emph{marque du présent}, il ne passe pas à l'impératif. On écrit : \all{Fahr vorsichtig!} (« Conduis prudemment ! »), \all{Lad die Datei herunter!} (« Télécharge le fichier ! »). Seuls les changements \all{e → i/ie} subsistent : \all{Lies!}, \all{Gib!}, \all{Nimm!}, \all{Hilf!}
\end{attention}

\begin{propriete}[Le verbe \all{sein} : un impératif irrégulier]
\all{sein} possède des formes d'impératif entièrement irrégulières, à mémoriser :
\begin{itemize}
\item \all{Sei ruhig!} = « Sois calme ! » (du)
\item \all{Seid kritisch!} = « Soyez critiques ! » (ihr)
\item \all{Seien Sie vorsichtig!} = « Soyez prudent ! » (Sie)
\end{itemize}
Très fréquent dans les consignes sur les médias : \all{Sei skeptisch gegenüber Fake News!} (« Sois sceptique face aux fausses informations ! »).
\end{propriete}

\courssub{Impératif et verbes à particule séparable}

\begin{attention}[La particule va en fin de phrase]
À l'impératif, le verbe à particule (\all{trennbares Verb}) \textbf{se sépare} : le verbe conjugué reste en position 1 et la particule va \textbf{en toute fin de phrase}. C'est l'une des erreurs les plus fréquentes des francophones.
\begin{itemize}
\item \all{Schaltet den Fernseher ein!} = « Allumez la télévision ! »
\item \all{Schalten Sie das Handy aus!} = « Éteignez le portable ! »
\item \all{Hör gut zu!} = « Écoute bien ! »
\item \all{Räumt euren Platz auf!} = « Rangez votre place ! »
\end{itemize}
Ne dites jamais \all{Schaltet ein den Fernseher!} : la particule \all{ein} doit fermer la phrase.
\end{attention}

\courssub{L'infinitif normatif : la langue des panneaux et des règlements}

Dans les règlements officiels, les panneaux, les notices et les consignes d'examen, l'allemand préfère souvent l'\cle{infinitif normatif} (\all{Infinitiv mit Imperativcharakter}) : l'infinitif seul, placé en fin de groupe, sans sujet. C'est l'équivalent exact du français « Ne pas fumer », « Se garer à l'arrière ».

\begin{exemplebox}[L'infinitif normatif en contexte médiatique]
\begin{itemize}
\item \all{Nicht rauchen!} = « Ne pas fumer ! »
\item \all{Handys ausschalten!} = « Éteindre les portables ! » (consigne de salle d'examen ou de cinéma)
\item \all{Nicht berühren!} = « Ne pas toucher ! »
\item \all{Bitte anschnallen!} = « Prière d'attacher sa ceinture ! »
\item \all{Keine Werbung einwerfen!} = « Ne pas insérer de publicité ! » (sur les boîtes aux lettres allemandes)
\item \all{Quellen angeben!} = « Indiquer ses sources ! » (consigne de devoir)
\end{itemize}
Remarquez : à l'infinitif normatif, la particule des verbes séparables \textbf{reste collée} au verbe : \all{ausschalten}, \all{angeben}.
\end{exemplebox}

\begin{methode}[Traduire une consigne française en allemand]
\begin{enumerate}
\item \textbf{Identifier le type de consigne.} « Ne pas + infinitif » ou « Infinitif » impersonnel (panneau, règlement) → infinitif normatif allemand.
\item \textbf{Traduire le verbe à l'infinitif}, en gardant la particule collée : « Ne pas toucher » → \all{Nicht berühren!}
\item \textbf{Consigne polie avec « veuillez »} → impératif \all{Sie} + \all{bitte} : « Veuillez patienter » → \all{Bitte warten Sie!} ; « Veuillez éteindre votre portable » → \all{Bitte schalten Sie Ihr Handy aus!}
\item \textbf{Vérifier l'ordre des mots} : verbe en position 1 à l'impératif, particule en fin de phrase, majuscule à tous les noms.
\end{enumerate}
\end{methode}

\courssub{Autres tournures injonctives : \all{sollen} et les particules modales}

L'injonction ne passe pas toujours par l'impératif. Deux outils sont précieux pour nuancer :

\begin{definition}[\all{sollen} + infinitif : l'ordre rapporté ou moral]
Le modal \all{sollen} exprime une obligation venue d'un tiers ou d'une règle : \all{Du sollst die Quellen prüfen.} = « Tu dois vérifier les sources. » / \all{Man soll nicht alles glauben, was man liest.} = « On ne doit pas croire tout ce qu'on lit. » C'est la forme préférée des règlements rédigés au style déclaratif.
\end{definition}

\begin{definition}[Les \all{Modalpartikeln} : adoucir l'ordre]
Les petites particules \all{mal} et \all{doch} transforment un ordre sec en invitation amicale :
\begin{itemize}
\item \all{Schau mal!} = « Regarde donc ! » (invitation, curiosité)
\item \all{Komm doch mit!} = « Viens donc avec nous ! » (insistance chaleureuse)
\item \all{Lies mal diesen Artikel!} = « Lis donc cet article ! »
\end{itemize}
Sans équivalent direct en français, elles se rendent par « donc », « un peu », le ton de la phrase.
\end{definition}

\courssub{L'impératif dans les slogans publicitaires}

Le slogan publicitaire allemand adore l'impératif : phrases courtes, rythme binaire, verbe en position 1 qui interpelle directement le consommateur. On y trouve surtout les formes \all{du} (proximité, jeunesse) et \all{ihr} (effet de groupe).

\begin{exemplebox}[Slogans et accroches à l'impératif]
\begin{itemize}
\item \all{Probier's mal mit Gemütlichkeit!} = « Essaie donc la douceur de vivre ! »
\item \all{Entdecke die Welt der Nachrichten!} = « Découvre le monde des informations ! »
\item \all{Bleib informiert!} = « Reste informé ! » (accroche d'un journal en ligne)
\item \all{Teilt eure Meinung!} = « Partagez votre opinion ! » (réseau social)
\end{itemize}
\end{exemplebox}

\begin{savaistu}[Les slogans allemands célèbres]
Attention, tous les slogans ne sont pas à l'impératif ! Le plus célèbre d'Allemagne, \all{Haribo macht Kinder froh} (« Haribo rend les enfants joyeux »), est une phrase déclarative au présent — mais son rythme et sa rime (\all{froh}) en font un slogan inoubliable. De même, \all{Geiz ist geil!} (« L'avarice, c'est génial ! »), slogan de la chaîne d'électronique Saturn, joue sur la rime entre deux adjectifs et a marqué toute une génération de consommateurs allemands. Leçon à retenir pour vos productions : un bon slogan est court, rythmé, souvent rimé — l'impératif y est fréquent mais pas obligatoire.
\end{savaistu}

\courssub{Comparaison français / allemand et récapitulatif}

\begin{center}
\begin{tabularx}{\linewidth}{|X|X|X|}\hline
\rowcolor{popblueL} \textbf{Français} & \textbf{Deutsch} & \textbf{Remarque} \\ \hline
Lis l'article ! & \all{Lies den Artikel!} & changement \all{e→ie} au \all{du} \\ \hline
Éteignez le portable ! (à plusieurs amis) & \all{Schaltet das Handy aus!} & particule \all{aus} en fin de phrase \\ \hline
Soyez critique ! (vouvoiement) & \all{Seien Sie kritisch!} & \all{sein} irrégulier : \all{seien Sie} \\ \hline
Ne pas fumer. & \all{Nicht rauchen!} & infinitif normatif = infinitif français \\ \hline
Veuillez patienter. & \all{Bitte warten Sie!} & « veuillez » = \all{bitte} + impératif \all{Sie} \\ \hline
Conduis prudemment ! & \all{Fahr vorsichtig!} & pas de tréma : jamais \all{Fähr!} \\ \hline
Tu dois vérifier les sources. & \all{Du sollst die Quellen prüfen.} & \all{sollen} = obligation rapportée \\ \hline
\end{tabularx}
\end{center}

\begin{aretenir}[L'essentiel de l'impératif]
\begin{itemize}
\item Trois formes : \all{du} = radical (\all{Komm! Lies!}) ; \all{ihr} = présent sans pronom (\all{Kommt! Lest!}) ; \all{Sie} = infinitif + \all{Sie} (\all{Kommen Sie!}).
\item Le verbe est toujours en \textbf{position 1} ; la particule des verbes séparables va en \textbf{fin de phrase}.
\item Changement vocalique \all{e→i/ie} conservé au \all{du} (\all{Gib! Nimm!}) ; jamais de tréma (\all{Fahr!}).
\item \all{sein} : \all{Sei! Seid! Seien Sie!}
\item Panneaux et règlements : \textbf{infinitif normatif} (\all{Handys ausschalten!}).
\item Pour adoucir : \all{mal}, \all{doch} ; pour l'obligation : \all{sollen} + infinitif.
\end{itemize}
\end{aretenir}

\begin{exoresolu}[Thème : traduis les consignes suivantes en allemand]
Énoncé. Traduis en allemand : 1. « Lis l'article et compare les sources ! » (tu) 2. « Éteignez vos portables ! » (vouvoiement) 3. « Ne pas toucher ! » (panneau) 4. « Soyez critiques face aux informations ! » (à un groupe d'élèves) 5. « Veuillez indiquer vos sources ! » (consigne de devoir)
\tcblower
Solution détaillée.\par
1. « Lis » = impératif \all{du} de \all{lesen} avec changement \all{e→ie} : \all{Lies den Artikel und vergleiche die Quellen!} (l'article = accusatif masculin \all{den Artikel}).\par
2. Vouvoiement → infinitif + \all{Sie} ; \all{ausschalten} se sépare, particule en fin : \all{Schalten Sie Ihre Handys aus!}\par
3. Panneau → infinitif normatif : \all{Nicht berühren!}\par
4. Groupe d'élèves → \all{ihr} ; verbe \all{sein} irrégulier : \all{Seid kritisch gegenüber den Informationen!}\par
5. « Veuillez » → \all{bitte} + impératif \all{Sie} ; \all{angeben} se sépare : \all{Bitte geben Sie Ihre Quellen an!}
\end{exoresolu}

\begin{exercice}[À toi de jouer]
Mets les verbes entre parenthèses à l'impératif demandé, puis traduis chaque phrase en français :
\begin{enumerate}
\item (\all{überprüfen}, du) \all{\ldots\ jede Nachricht, bevor du sie teilst!}
\item (\all{ausschalten}, ihr) \all{\ldots\ bitte eure Handys \ldots!}
\item (\all{sein}, Sie) \all{\ldots\ \ldots\ vorsichtig mit den sozialen Netzwerken!}
\item (\all{zuhören}, du) \all{\ldots\ gut \ldots, wenn der Betreuer spricht!}
\item Traduis en allemand : « Ne pas déranger ! » et « Veuillez ranger votre place ! »
\end{enumerate}
\end{exercice}

\begin{corrige}[Exercice « À toi de jouer »]
1. \all{Überprüf(e) jede Nachricht, bevor du sie teilst!} = « Vérifie chaque information avant de la partager ! »\par
2. \all{Schaltet bitte eure Handys aus!} = « Éteignez vos portables, s'il vous plaît ! » (particule \all{aus} en fin).\par
3. \all{Seien Sie vorsichtig mit den sozialen Netzwerken!} = « Soyez prudent avec les réseaux sociaux ! »\par
4. \all{Hör gut zu, wenn der Betreuer spricht!} = « Écoute bien quand le responsable parle ! »\par
5. \all{Nicht stören!} ; \all{Bitte räumen Sie Ihren Platz auf!}
\end{corrige}

Ces formes injonctives maîtrisées, nous pourrons dans la section suivante les articuler avec le discours indirect et le passif, deux outils essentiels pour rapporter et reformuler l'information dans les médias.

\coursec{Grammaire 2 : discours indirect et passif}

Dans la section précédente, nous avons vu comment les médias \emph{donnent des ordres} au public grâce à l'impératif et aux formes d'injonction (\all{Schalten Sie ein! Lesen Sie jetzt!}). Mais les médias font aussi le contraire : ils \emph{rapportent} ce que d'autres ont dit et ils \emph{décrivent des événements} sans toujours nommer les acteurs. Pour cela, la langue journalistique allemande dispose de deux outils majeurs : le \cle{discours indirect} au subjonctif I (\all{die indirekte Rede}) et le \cle{passif} (\all{das Passiv}). Ce sont deux points de grammaire très fréquents au baccalauréat, en version comme en thème.

\courssub{Observation : un article de presse}

Lisons d'abord une courte dépêche et soulignons mentalement les formes verbales étonnantes.

\begin{textebox}[Dépêche d'agence : Pressefreiheit in Kamerun]
Yaoundé. Auf einer Medienkonferenz in der kamerunischen Hauptstadt erklärte der Kommunikationsminister, die Pressefreiheit sei ein wichtiger Teil der Demokratie. Er sagte, die Regierung werde neue Radiosender unterstützen, und fügte hinzu, dass junge Journalisten besser ausgebildet werden müssten. Ein Vertreter der Journalisten meinte, viele Kollegen hätten zu wenig Geld und zu wenig Material. Er gab an, dass in den Provinzen oft keine Zeitungen gedruckt werden könnten.

Am Rande der Konferenz wurde ein neuer Preis für investigative Journalisten angekündigt. Der Preis wird jedes Jahr von einer unabhängigen Jury vergeben. Im vergangenen Jahr wurden über zweihundert Artikel und Reportagen eingereicht. Die Gewinner wurden im Fernsehen und in den sozialen Netzwerken bekannt gegeben. Ein junger Journalist aus Douala sagte, er habe sich sehr über die Auszeichnung gefreut. Die nächste Konferenz werde in Bafoussam stattfinden, wurde am Ende mitgeteilt.
\end{textebox}

\textbf{Glossaire :} \all{die Pressefreiheit} la liberté de la presse ; \all{der Radiosender, -} la station de radio ; \all{ausbilden} former ; \all{der Vertreter, -} le représentant ; \all{angeben (gab an, hat angegeben)} indiquer, déclarer ; \all{drucken} imprimer ; \all{ankündigen} annoncer ; \all{vergeben (vergab, hat vergeben)} attribuer ; \all{einreichen} soumettre ; \all{die Auszeichnung, -en} la récompense ; \all{mitteilen} communiquer.

\textbf{Questions de compréhension :}
\begin{enumerate}
\item Que déclare le ministre de la Communication ?
\item Qui attribue le nouveau prix, et à quelle fréquence ?
\item Relevez trois verbes conjugués au subjonctif I et trois formes du passif.
\end{enumerate}

\textbf{Observation grammaticale.} Remarquez : \all{die Pressefreiheit \textbf{sei} wichtig}, \all{die Regierung \textbf{werde} unterstützen}, \all{Kollegen \textbf{hätten} zu wenig Geld}, \all{er \textbf{habe} sich gefreut}. Ces formes ne sont ni l'indicatif (\all{ist, wird, haben, hat}) ni le subjonctif II de l'irréel : c'est le \cle{subjonctif I}, le temps-mode du propos rapporté. De même : \all{wurde angekündigt}, \all{wird vergeben}, \all{wurden eingereicht} : c'est le \cle{passif}.

\courssub{Le discours indirect : fonction et marqueurs}

\begin{definition}[Le discours indirect]
Le discours indirect (\all{die indirekte Rede}) sert à \textbf{rapporter les paroles de quelqu'un} sans les citer mot pour mot. C'est l'outil de base du journaliste, qui doit rester neutre et ne pas garantir la vérité de ce qu'il rapporte. En allemand, il se construit avec le \textbf{subjonctif I} (\all{der Konjunktiv I}).
\end{definition}

Les \textbf{verbes introducteurs} (\all{Verben des Sagens}) les plus utiles dans les textes médiatiques :

\begin{center}
\begin{tabularx}{\linewidth}{|X|X|}
\hline
\rowcolor{popblueL} Deutsch & Français \\
\hline
\all{berichten} & rapporter, faire un reportage \\
\hline
\all{erklären} & déclarer, expliquer \\
\hline
\all{meinen} & estimer, penser \\
\hline
\all{behaupten} & affirmer, prétendre \\
\hline
\all{mitteilen} & communiquer, annoncer \\
\hline
\all{angeben (gab an, hat angegeben)} & indiquer, déclarer \\
\hline
\all{hinzufügen} & ajouter \\
\hline
\all{betonen} & souligner \\
\hline
\all{warnen vor + Dativ} & mettre en garde contre \\
\hline
\all{ankündigen} & annoncer \\
\hline
\end{tabularx}
\end{center}

\begin{exemplebox}[Marqueurs en contexte]
\all{Der Journalist \textbf{berichtet}, dass viele Jugendliche keine Zeitung mehr lesen.} « Le journaliste rapporte que beaucoup de jeunes ne lisent plus de journal. »

\all{Die Sprecherin \textbf{erklärte}, der Sender werde das Programm ändern.} « La porte-parole a déclaré que la chaîne allait modifier le programme. »

\all{Der Experte \textbf{meint}, soziale Netzwerke seien gefährlich für junge Menschen.} « L'expert estime que les réseaux sociaux sont dangereux pour les jeunes. »
\end{exemplebox}

\courssub{Le subjonctif I : formation}

\begin{propriete}[Formation du subjonctif I]
Le subjonctif I se forme sur le \textbf{radical de l'infinitif} + les terminaisons :

\begin{center}
\begin{tabular}{|l|l|}
\hline
\rowcolor{popgreenL} Personne & Terminaison \\
\hline
ich & -e \\
\hline
du & -est \\
\hline
er/sie/es & \textbf{-e} \\
\hline
wir & -en \\
\hline
ihr & -et \\
\hline
sie/Sie & -en \\
\hline
\end{tabular}
\end{center}

Exemple avec \all{sagen} : ich sage, du sagest, er \textbf{sage}, wir sagen, ihr saget, sie sagen.

\textbf{Seule exception :} \all{sein} est irrégulier (pas de \emph{-e} à la 1\textsuperscript{re} et à la 3\textsuperscript{e} personne du singulier) : ich \textbf{sei}, du seiest/seist, er \textbf{sei}, wir seien, ihr seiet, sie seien.

En pratique journalistique, on utilise presque uniquement la \textbf{3\textsuperscript{e} personne du singulier} (\all{er sei, er habe, er werde, er gebe}) et du pluriel.
\end{propriete}

\begin{popfigure}
\begin{center}
\begin{tikzpicture}
\node[draw=popblue, fill=popblueL, rounded corners, align=center, font=\bfseries] (titre) at (0,0) {Konjunktiv I\\ (3. Person)};
\node[draw=popgreen, fill=popgreenL, rounded corners, align=center] (sein) at (-4.5,-2) {\textbf{sein}\\ er/sie/es \textbf{sei}\\ sie \textbf{seien}};
\node[draw=poporange, fill=poporangeL, rounded corners, align=center] (haben) at (-1.5,-2) {\textbf{haben}\\ er \textbf{habe}\\ sie \textbf{haben}};
\node[draw=poppurple, fill=poppurpleL, rounded corners, align=center] (werden) at (1.5,-2) {\textbf{werden}\\ er \textbf{werde}\\ sie \textbf{werden}};
\node[draw=poppink, fill=poppinkL, rounded corners, align=center] (koennen) at (4.5,-2) {\textbf{können}\\ er \textbf{könne}\\ sie \textbf{können}};
\draw[-{Stealth}, thick, popblue] (titre) -- (sein);
\draw[-{Stealth}, thick, popblue] (titre) -- (haben);
\draw[-{Stealth}, thick, popblue] (titre) -- (werden);
\draw[-{Stealth}, thick, popblue] (titre) -- (koennen);
\node[draw=popgold, fill=popgoldL, rounded corners, align=center] (modal) at (0,-3.8) {Autres modaux : er \textbf{müsse}, er \textbf{solle}, er \textbf{dürfe}, er \textbf{wolle}, er \textbf{möge}};
\draw[-{Stealth}, thick, popgold] (titre.south) .. controls (0,-1.5) .. (modal.north);
\end{tikzpicture}
\legende{Les formes du subjonctif I les plus utiles au baccalauréat}
\end{center}
\end{popfigure}

\begin{aretenir}[Les quatre formes à connaître par cœur]
\begin{itemize}
\item \all{sein} $\rightarrow$ er \textbf{sei} : \all{Er sagte, die Lage sei ernst.} « Il a dit que la situation était grave. »
\item \all{haben} $\rightarrow$ er \textbf{habe} : \all{Sie erklärte, sie habe keine Zeit.} « Elle a déclaré qu'elle n'avait pas le temps. »
\item \all{werden} $\rightarrow$ er \textbf{werde} : \all{Der Sender teilte mit, die Sendung werde ausfallen.} « La chaîne a annoncé que l'émission serait annulée. »
\item \all{es gibt} $\rightarrow$ es \textbf{gebe} : \all{Der Artikel sagt, es gebe zu viele Fake News.} « L'article dit qu'il y a trop de fausses informations. »
\end{itemize}
\end{aretenir}

\courssub{Discours indirect avec \all{dass} ou avec inversion}

Le propos rapporté peut être introduit de deux façons, toutes deux correctes :

\begin{enumerate}
\item \textbf{Avec \all{dass}} : le verbe conjugué va à la fin de la subordonnée.\\
\all{Er sagte, \textbf{dass} die Sendung um 20 Uhr \textbf{beginne}.} « Il a dit que l'émission commençait à 20 h. »
\item \textbf{Sans \all{dass}} (inversion) : le verbe conjugué occupe la deuxième place, comme dans une principale.\\
\all{Er sagte, die Sendung \textbf{beginne} um 20 Uhr.}
\end{enumerate}

\begin{attention}[Ordre des mots]
Sans \all{dass}, la phrase rapportée garde l'ordre de la phrase principale (verbe en position 2). Erreur fréquente du francophone : écrire \all{Er sagte, die Sendung um 20 Uhr beginne} sans \all{dass} : c'est impossible. Soit \all{dass} + verbe final, soit pas de \all{dass} + verbe en position 2.
\end{attention}

\courssub{Le passé du discours indirect et la substitution}

Pour rapporter un propos qui parlait du \textbf{passé}, on utilise le subjonctif I du \textbf{parfait} : \all{habe/sei} + participe II.

\begin{itemize}
\item Direct : \all{„Ich habe den Fehler gesehen.“} $\rightarrow$ Indirect : \all{Er sagte, er \textbf{habe} den Fehler \textbf{gesehen}.} « Il a dit qu'il avait vu l'erreur. »
\item Direct : \all{„Sie ist nach Douala gefahren.“} $\rightarrow$ Indirect : \all{Man berichtete, sie \textbf{sei} nach Douala \textbf{gefahren}.} « On a rapporté qu'elle était partie à Douala. »
\end{itemize}

\begin{propriete}[Substitution par le subjonctif II]
Quand la forme du subjonctif I est \textbf{identique à l'indicatif} (donc ambiguë), on la remplace par le \textbf{subjonctif II}. C'est le cas surtout au pluriel et à la 1\textsuperscript{re} personne :

\begin{itemize}
\item \all{sie sagten, sie \textbf{hätten} keine Informationen} (et non \all{sie haben}, qui serait l'indicatif) « elles ont dit qu'elles n'avaient pas d'informations »
\item \all{er sagte, die Leute \textbf{gäben} nicht genug Acht} (remplace \all{geben}, identique à l'indicatif) « il a dit que les gens ne faisaient pas assez attention »
\item \all{der Minister erklärte, die Sender \textbf{müssten} moderner werden} (remplace \all{müssen}, identique à l'indicatif) « le ministre a déclaré que les chaînes devaient se moderniser »
\end{itemize}

En revanche, pas de substitution quand la forme du subjonctif I est déjà distincte : \all{sie \textbf{seien} müde} est correct au pluriel, car l'indicatif serait \all{sind}.
\end{propriete}

\begin{popfigure}
\begin{center}
\begin{tikzpicture}
\node[draw=popblue, fill=popblueL, rounded corners, align=center] (d1) at (0,0) {\textbf{Discours direct}\\ „Die Sendung \textbf{beginnt} um 20 Uhr.“\\ (présent)};
\node[draw=popgreen, fill=popgreenL, rounded corners, align=center] (i1) at (8.5,0) {\textbf{Discours indirect}\\ Er sagte, die Sendung \textbf{beginne} um 20 Uhr.\\ (Konjunktiv I)};
\draw[-{Stealth}, very thick, poporange] (d1) -- (i1) node[midway, above, align=center, font=\small]{présent $\rightarrow$ KI};
\node[draw=popblue, fill=popblueL, rounded corners, align=center] (d2) at (0,-3) {\textbf{Discours direct}\\ „Ich \textbf{habe} den Artikel \textbf{geschrieben}.“\\ (passé)};
\node[draw=popgreen, fill=popgreenL, rounded corners, align=center] (i2) at (8.5,-3) {\textbf{Discours indirect}\\ Sie sagte, sie \textbf{habe} den Artikel \textbf{geschrieben}.\\ (KI du parfait : habe/sei + P.II)};
\draw[-{Stealth}, very thick, poporange] (d2) -- (i2) node[midway, above, align=center, font=\small]{passé $\rightarrow$ habe/sei + Partizip II};
\end{tikzpicture}
\legende{Transformation du discours direct en discours indirect}
\end{center}
\end{popfigure}

\begin{attention}[Pas de concordance des temps en allemand]
Le français impose une concordance : « il a dit qu'il \emph{viendrait} » (conditionnel). L'allemand, lui, garde simplement le \textbf{subjonctif I du présent}, même après un verbe introducteur au passé. Ne calquez donc jamais le conditionnel français par \all{würde} + infinitif dans un texte journalistique :

« Il a dit qu'il viendrait. » = \all{Er sagte, er \textbf{werde} kommen.} (et non \all{er würde kommen}, qui sonne comme un doute ou un irréel).

De même : « Le ministre a déclaré que la liberté de la presse était importante. » = \all{Der Minister erklärte, die Pressefreiheit \textbf{sei} wichtig.}
\end{attention}

\courssub{Le passif (das Vorgangspassiv)}

\begin{definition}[Le passif]
Le passif met en avant \textbf{l'action ou son résultat}, et non l'auteur de l'action. C'est la construction favorite des dépêches d'agence : \all{Die Fake News wurden schnell verbreitet.} « Les fausses informations ont été diffusées rapidement. »
\end{definition}

\begin{propriete}[Formation du Vorgangspassiv]
\textbf{Passif = auxiliaire \all{werden} conjugué + participe II} (placé en fin de phrase).

\begin{center}
\begin{tabular}{|l|l|l|}
\hline
\rowcolor{poporangeL} Temps & Auxiliaire & Exemple \\
\hline
Présent & wird/werden + P.II & \all{Die Nachrichten \textbf{werden gesendet}.} \\
\hline
Prétérit & wurde/wurden + P.II & \all{Die Nachrichten \textbf{wurden gesendet}.} \\
\hline
Parfait & ist/sind + P.II + worden & \all{Die Nachrichten \textbf{sind gesendet worden}.} \\
\hline
Futur & wird + P.II + werden & \all{Der Film \textbf{wird gezeigt werden}.} \\
\hline
\end{tabular}
\end{center}

Conjugaison de l'auxiliaire au prétérit : ich wurde, du wurdest, er wurde, wir wurden, ihr wurdet, sie wurden. Attention au participe du passif au parfait : c'est \all{\textbf{worden}} (et non \all{geworden}).
\end{propriete}

\begin{exemplebox}[Le passif dans les textes médiatiques]
\all{Der Artikel \textbf{wird} morgen von der Journalistin \textbf{geschrieben}.} « L'article sera écrit demain par la journaliste. » (futur proche rendu par le présent + complément de temps)

\all{Die Nachricht \textbf{wurde} von der Redaktion \textbf{veröffentlicht}.} « L'information a été publiée par la rédaction. »

\all{In Kamerun \textbf{werden} viele Programme auf Englisch und Französisch \textbf{gesendet}.} « Au Cameroun, beaucoup d'émissions sont diffusées en anglais et en français. »

\all{Der Gewinner \textbf{wurde} im Fernsehen \textbf{bekannt gegeben}.} « Le gagnant a été annoncé à la télévision. »
\end{exemplebox}

\textbf{L'agent} se rend par \all{\textbf{von} + datif} : \all{von der Redaktion}, \all{von dem Minister}, \all{von den Journalisten}. Mais dans les textes médiatiques, l'agent est \textbf{très souvent omis}, parce qu'il est inconnu, évident ou sans importance : \all{Es wurde ein neuer Preis angekündigt.} « Un nouveau prix a été annoncé. » (par qui ? le lecteur s'en moque : c'est l'événement qui compte).

\begin{attention}[Place des éléments au passif]
Au passif, c'est \all{werden} qui porte toute la conjugaison (personne, temps, position 2 dans la phrase) ; le participe II se place \textbf{en toute fin de phrase}. Ne dites jamais \all{Der Artikel wird geschrieben morgen} : le complément de temps se place au milieu : \all{Der Artikel \textbf{wird} morgen \textbf{geschrieben}.} Autre piège : ne confondez pas \all{werden} (auxiliaire du passif) et \all{sein} : \all{die Tür ist geschlossen} décrit un \emph{état} (la porte est fermée), \all{die Tür wird geschlossen} décrit une \emph{action} (on ferme la porte).
\end{attention}

\courssub{Traduire le français « on » : passif ou \all{man}}

\begin{methode}[Comment traduire « on » en allemand]
\begin{enumerate}
\item \textbf{L'action prime, l'agent est sans importance} $\rightarrow$ utilisez le \textbf{passif} : « On diffuse les infos à 20 h. » = \all{Die Nachrichten \textbf{werden} um 20 Uhr \textbf{gesendet}.}
\item \textbf{L'agent est indéfini mais présent à l'esprit} (les gens, quelqu'un) $\rightarrow$ utilisez \textbf{\all{man} + actif} : « On rapporte que... » = \all{\textbf{Man} berichtet, dass...}
\item \textbf{« On » = « nous »} dans la conversation $\rightarrow$ \all{wir} : « On va au marché demain. » = \all{\textbf{Wir} gehen morgen auf den Markt.}
\end{enumerate}
Dans un texte journalistique, la solution 1 (passif) est presque toujours la meilleure.
\end{methode}

\begin{center}
\begin{tabularx}{\linewidth}{|X|X|X|}
\hline
\rowcolor{popblueL} Français & Deutsch & Remarque \\
\hline
On diffuse les infos à 20 h. & \all{Die Nachrichten werden um 20 Uhr gesendet.} & passif : l'action prime \\
\hline
On a annoncé le résultat. & \all{Das Ergebnis wurde bekannt gegeben.} & prétérit passif \\
\hline
On rapporte que le ministre démissionnera. & \all{Man berichtet, dass der Minister zurücktreten werde.} & \all{man} + discours indirect \\
\hline
On a beaucoup parlé de ce reportage. & \all{Über diese Reportage wurde viel gesprochen.} & sujet français « on » disparaît \\
\hline
Il a dit qu'il viendrait. & \all{Er sagte, er werde kommen.} & pas de conditionnel ! \\
\hline
\end{tabularx}
\end{center}

\courssub{Exercice résolu et entraînement}

\begin{exoresolu}[Thème : du français à l'allemand]
Traduisez en allemand : « Le présentateur a déclaré que l'émission était très populaire. Les résultats ont été publiés hier par la rédaction. On a aussi annoncé qu'un nouveau journal serait créé à Yaoundé. »
\tcblower
\textbf{Solution détaillée :}

1. « Le présentateur a déclaré que l'émission était très populaire. » Verbe introducteur au passé + propos rapporté $\rightarrow$ subjonctif I : \all{Der Moderator erklärte, die Sendung \textbf{sei} sehr beliebt.} (ou : \all{..., dass die Sendung sehr beliebt sei.}) Pas de concordance des temps : \all{sei} suffit.

2. « Les résultats ont été publiés hier par la rédaction. » Passif au prétérit + agent avec \all{von} + datif : \all{Die Ergebnisse \textbf{wurden} gestern \textbf{von der Redaktion veröffentlicht}.} Le participe \all{veröffentlicht} reste en fin de phrase.

3. « On a aussi annoncé qu'un nouveau journal serait créé à Yaoundé. » « On » $\rightarrow$ passif impersonnel ; « serait créé » $\rightarrow$ subjonctif I du passif (\all{werde} + participe II + \all{werden}) : \all{Es \textbf{wurde} auch \textbf{angekündigt}, dass eine neue Zeitung in Yaoundé \textbf{gegründet werde}.} Variante acceptée avec substitution : \all{..., dass eine neue Zeitung in Yaoundé gegründet \textbf{würde}.} — mais la forme en subjonctif I est préférable dans un texte journalistique.
\end{exoresolu}

\begin{exercice}[Version et thème]
\textbf{A. Version.} Traduisez en français : \all{Der Sprecher teilte mit, dass die Pressekonferenz verschoben worden sei. Viele Journalisten seien nicht informiert worden, behauptete ein Reporter. Die Fotos wurden von einer Agentur verbreitet.}

\textbf{B. Thème.} Traduisez en allemand : 1. « Le journaliste a écrit que la situation était difficile. » 2. « On diffuse les nouvelles chaque soir à 20 h. » 3. « Ils ont dit qu'ils n'avaient pas d'argent. » 4. « L'interview a été enregistrée hier à Douala. »
\end{exercice}

\begin{corrige}[Exercice : version et thème]
\textbf{A.} « Le porte-parole a communiqué que la conférence de presse avait été reportée. Beaucoup de journalistes n'avaient pas été informés, a affirmé un reporter. Les photos ont été diffusées par une agence. »

\textbf{B.} 1. \all{Der Journalist schrieb, die Lage sei schwierig.} (ou \all{..., dass die Lage schwierig sei.}) 2. \all{Die Nachrichten werden jeden Abend um 20 Uhr gesendet.} 3. \all{Sie sagten, sie hätten kein Geld.} (substitution : \all{haben} serait identique à l'indicatif, donc subjonctif II \all{hätten}). 4. \all{Das Interview wurde gestern in Douala aufgenommen.}
\end{corrige}

\begin{savaistu}[Les agences de presse et la langue neutre]
Les grandes agences de presse germanophones — la \textbf{dpa} (Deutsche Presse-Agentur) en Allemagne, l'\textbf{APA} en Autriche et \textbf{Keystone-ATS} en Suisse — rédigent leurs dépêches massivement au \textbf{passif} et au \textbf{subjonctif I}. Pourquoi ? Parce qu'une agence ne doit ni nommer des coupables sans preuve, ni garantir la vérité des déclarations qu'elle rapporte. Le subjonctif I signale au lecteur : « c'est ce que X a dit, à vous de juger ». C'est une écriture de la \emph{neutralité professionnelle}. Au Cameroun, les journalistes formés à l'allemand retrouvent exactement la même démarche dans les rédactions bilingues de Yaoundé et de Douala.
\end{savaistu}

\begin{aretenir}[L'essentiel de la section]
\begin{itemize}
\item Discours indirect = verbe introducteur (\all{berichten, erklären, meinen...}) + \textbf{subjonctif I} : \all{er sei, er habe, er werde, es gebe}.
\item Passé rapporté : \all{habe/sei} + participe II ; forme ambiguë $\rightarrow$ subjonctif II (\all{sie hätten}).
\item Pas de concordance des temps : jamais \all{würde} pour calquer le conditionnel français.
\item Passif = \all{werden} conjugué + participe II en fin de phrase ; agent avec \all{von} + datif, souvent omis.
\item « On » français $\rightarrow$ passif allemand (action) ou \all{man} + actif (agent indéfini).
\end{itemize}
\end{aretenir}

\coursec{Méthode de traduction : thème et version}

Après avoir révisé l'impératif, le discours indirect et le passif, nous disposons de tous les outils grammaticaux nécessaires pour aborder l'épreuve reine du chapitre : la \cle{traduction}. Au Baccalauréat A, deux exercices complémentaires évaluent cette compétence : la \cle{version} (traduire de l'allemand vers le français) et le \cle{thème} (traduire du français vers l'allemand). Les textes proposés appartiennent souvent à la famille des textes injonctifs : affiches, règlements, consignes, slogans publicitaires liés aux médias. Cette section vous donne une méthode rigoureuse, étape par étape, pour ne plus jamais traduire « au feeling ».

\courssub{Observer avant de traduire : deux exemples}

\begin{textebox}[Deux textes injonctifs à observer]
\textbf{Texte 1 (panneau officiel) :}\\
„Nicht hinauslehnen! Das Rauchen ist in der Redaktion verboten. Handys bitte ausschalten!“\\[0.3em]
\textbf{Texte 2 (slogan d'une campagne médiatique) :}\\
„Überprüft immer eure Quellen, bevor ihr einen Artikel teilt. Informiert euch, bevor ihr urteilt!“
\end{textebox}

Observons : le texte 1 utilise l'\cle{infinitif normatif} (\all{hinauslehnen}, \all{ausschalten}) et une tournure passive (\all{ist ... verboten}) ; le texte 2 utilise l'\cle{impératif} (\all{überprüft}, \all{informiert euch}). En français, on rendra le premier par « Ne pas se pencher au dehors ! Il est interdit de fumer dans la rédaction. Éteindre les portables, s'il vous plaît ! » et le second par « Vérifiez toujours vos sources avant de partager un article. Informez-vous avant de juger ! ». La traduction n'est donc pas un transfert mot à mot : c'est d'abord un \cle{choix de structure} adapté au type de texte.

\begin{aretenir}[Le principe d'or du traducteur]
On ne traduit pas des mots, on traduit des \cle{structures} et des \cle{intentions}. Une consigne allemande à l'infinitif normatif appelle en français « ne pas + infinitif » ; un impératif allemand appelle un impératif français ; un passif allemand appelle un passif ou une tournure impersonnelle française. Identifier le type de texte (affiche, règlement, slogan) est donc la première décision du traducteur.
\end{aretenir}

\courssub{La version : de l'allemand vers le français}

\begin{methode}[Les 5 étapes de la version]
\begin{enumerate}
\item \textbf{Lire tout le texte} une première fois sans écrire, pour saisir le sens global et identifier le type de texte (affiche, règlement, slogan, article).
\item \textbf{Repérer les structures} : impératif ou infinitif normatif ? Passif (\all{werden} + participe II) ? Discours indirect (Konjunktiv I) ? Verbe à particule dont la particule attend en fin de phrase ?
\item \textbf{Traduire par phrases entières}, jamais mot à mot : reformuler en français naturel.
\item \textbf{Choisir l'équivalent français} de la structure injonctive : impératif allemand → impératif français ; infinitif normatif → « ne pas + infinitif » ou interdiction nominale.
\item \textbf{Relire} la version française : est-ce du français correct et idiomatique ? Aucun calque de l'allemand ?
\end{enumerate}
\end{methode}

\begin{propriete}[Traduire les formes d'injonction en version]
\begin{itemize}
\item Un \cle{impératif} allemand se traduit par un impératif français : \all{Schaltet das Handy aus!} = « Éteignez le portable ! »
\item Un \cle{infinitif normatif} allemand se traduit par « ne pas + infinitif » ou par une interdiction : \all{Nicht hinauslehnen!} = « Ne pas se pencher au dehors ! » ; \all{Bitte nicht stören!} = « Prière de ne pas déranger ! »
\item Un \cle{passif} se traduit par un passif ou une tournure active impersonnelle : \all{Das Rauchen ist verboten.} = « Il est interdit de fumer. » / « Fumer est interdit. »
\end{itemize}
\end{propriete}

\begin{attention}[Le verbe à particule : le piège classique de la version]
En allemand, la particule séparable attend à la \cle{fin de la phrase}. Ne traduisez jamais avant d'avoir lu la dernière ligne ! \all{Der Moderator schaltet das Mikrofon ein.} : si vous traduisez \all{schaltet} seul, vous écrivez « commute » ; or \all{einschalten} = « allumer ». La phrase signifie : « Le présentateur allume le micro. » Autres exemples : \all{Sie gibt die Nachrichten bekannt.} (« Elle annonce les nouvelles »), \all{Er hört mit dem Rauchen auf.} (« Il arrête de fumer »).
\end{attention}

\courssub{Le thème : du français vers l'allemand}

\begin{methode}[Les 5 étapes du thème]
\begin{enumerate}
\item \textbf{Lire et comprendre} la phrase française en entier ; identifier le type de texte et le destinataire.
\item \textbf{Repérer sujet, verbe, compléments} et déterminer le temps et le mode français (indicatif, impératif, passif...).
\item \textbf{Choisir la structure allemande} équivalente : impératif ou infinitif normatif pour une consigne, \all{werden} + participe II pour un passif, Konjunktiv I pour une parole rapportée.
\item \textbf{Construire la phrase allemande} en respectant l'ordre des mots : verbe conjugué en 2\textsuperscript{e} position (ou en 1\textsuperscript{re} à l'impératif), particule et participe en fin de phrase, verbe en position finale dans la subordonnée.
\item \textbf{Relire} : cas (Nominatif, Akkusativ, Dativ), majuscules à tous les noms, orthographe (ß/ss, Umlaute), accords.
\end{enumerate}
\end{methode}

\begin{popfigure}
\begin{tikzpicture}[node distance=0.55cm,
box/.style={rectangle, rounded corners=3pt, draw=popblue, fill=popblueL, align=center, inner sep=5pt, font=\small},
arr/.style={-{Stealth[length=2.5mm]}, thick, poporange}]
\node[box, fill=poporangeL, draw=poporange, align=center] (p){Phrase française\\« Vérifiez vos sources. »};
\node[box, below=of p, align=center] (a){1. Analyse\\S / V / C + temps et mode};
\node[box, below=of a, align=center] (c){2. Choix du mode\\impératif ? infinitif ? passif ?};
\node[box, below=of c, align=center] (k){3. Construction allemande\\verbe en 2\textsuperscript{e} position, cas, particule finale};
\node[box, below=of k, fill=popgreenL, draw=popgreen, align=center] (r){4. Relecture\\majuscules, cas, ß/ss, ordre des mots};
\draw[arr] (p) -- (a);
\draw[arr] (a) -- (c);
\draw[arr] (c) -- (k);
\draw[arr] (k) -- (r);
\end{tikzpicture}
\legende{Organigramme du thème : de la phrase française à la phrase allemande relue.}
\end{popfigure}

\begin{exemplebox}[Thème commenté pas à pas]
Phrase à traduire : « Vérifiez toujours vos sources avant de partager un article. »
\begin{enumerate}
\item Analyse : impératif (2\textsuperscript{e} pers. plur.) ; subordonnée « avant de + infinitif ».
\item Choix : impératif allemand du pluriel \all{überprüft} ; « avant de partager » → subordonnée avec \all{bevor} + verbe conjugué (\all{bevor ihr ... teilt}), car l'allemand n'a pas de construction « avant de + infinitif » avec sujet exprimé.
\item Construction : verbe en 1\textsuperscript{re} position (impératif), verbe de la subordonnée en position finale.
\item Résultat : \all{Überprüft immer eure Quellen, bevor ihr einen Artikel teilt.}
\end{enumerate}
Version commentée : \all{Das Rauchen ist in der Redaktion verboten.} → passif d'état avec \all{sein} + participe II ; équivalent français impersonnel : « Il est interdit de fumer dans la rédaction. »
\end{exemplebox}

\begin{attention}[« Il est interdit de... » : le faux pas du candidat]
Ne traduisez JAMAIS « Il est interdit de fumer » par \all{Er ist verboten} (cela voudrait dire « Lui est interdit » !). Deux solutions correctes :
\begin{itemize}
\item Tournure impersonnelle : \all{Es ist verboten, hier zu rauchen.}
\item Infinitif normatif (style règlement) : \all{Hier nicht rauchen!} ou avec le substantif : \all{Das Rauchen ist verboten.}
\end{itemize}
De même, « avant de + infinitif » ne se rend JAMAIS par \all{vor zu} : utilisez \all{bevor} + proposition conjuguée (\all{bevor ihr teilt}) ou \all{vor} + nom (\all{vor dem Teilen eines Artikels}).
\end{attention}

\courssub{Traduire le passif et les temps}

\begin{propriete}[Passif français → passif allemand]
« être + participe passé » se traduit par \all{werden} + participe II (passif d'action) : « L'article est publié » = \all{Der Artikel wird veröffentlicht.} Attention au temps : le passé composé français « a été publié » se rend souvent par un prétérit allemand : \all{Der Artikel wurde veröffentlicht.} En version, repérez l'agent introduit par \all{von} (+ Dativ) : \all{Die Sendung wird von Millionen Zuschauern gesehen.} = « L'émission est regardée par des millions de téléspectateurs. »
\end{propriete}

\begin{propriete}[Rappel : les formes du passif aux temps principaux]
\begin{center}
\begin{tabular}{|l|l|l|}
\hline
\rowcolor{popblueL} Temps & Allemand & Français \\
\hline
Présent & Der Artikel \textbf{wird} veröffentlicht. & L'article est publié. \\
\hline
Prétérit & Der Artikel \textbf{wurde} veröffentlicht. & L'article a été publié. \\
\hline
Perfekt & Der Artikel \textbf{ist} veröffentlicht \textbf{worden}. & L'article a été publié. \\
\hline
Futur & Der Artikel \textbf{wird} veröffentlicht \textbf{werden}. & L'article sera publié. \\
\hline
\end{tabular}
\end{center}
Aux temps composés, le participe II de \all{werden} devient \all{worden} (et non \all{geworden}) dans la construction passive.
\end{propriete}

\begin{center}
\begin{tabularx}{\linewidth}{|X|X|X|}
\hline
\rowcolor{popblueL} Français & Deutsch & Remarque \\
\hline
Il est interdit de fumer. & Es ist verboten, zu rauchen. / Rauchen verboten! & jamais \all{er ist verboten} \\
\hline
Éteignez vos portables ! & Schaltet eure Handys aus! & particule \all{aus} en finale \\
\hline
Ne pas se pencher au dehors ! & Nicht hinauslehnen! & infinitif normatif \\
\hline
Vérifiez vos sources. & Überprüft eure Quellen. & \all{die Quelle} = la source \\
\hline
Il devient journaliste. & Er wird Journalist. & « devenir » = \all{werden} \\
\hline
L'article a été publié hier. & Der Artikel wurde gestern veröffentlicht. & passé composé → prétérit \\
\hline
Il reçoit un prix. & Er bekommt einen Preis. & \all{bekommen} = recevoir \\
\hline
Les informations actuelles & die aktuellen Nachrichten & « maintenant » = \all{jetzt} \\
\hline
\end{tabularx}
\end{center}

\begin{attention}[Pièges lexicaux : les faux amis du chapitre]
\begin{itemize}
\item \all{bekommen} = recevoir ; « devenir » = \all{werden}.
\item \all{die Quelle, -n} = la source ; « la qualité » = \all{die Qualität}.
\item \all{aktuell} = actuel (qui concerne le présent) ; « actuellement » au sens de « maintenant » = \all{jetzt}, \all{im Moment}.
\item \all{die Nachricht, -en} = la nouvelle, l'information ; « le message » = \all{die Botschaft} ou \all{die Mitteilung}.
\end{itemize}
\end{attention}

\begin{exoresolu}[Thème guidé : trois phrases sur les médias]
Énoncé : traduisez en allemand. 1) « Ne touchez pas aux câbles ! » (règlement du studio). 2) « Lisez les nouvelles avant de les commenter. » 3) « Le reportage a été regardé par beaucoup de téléspectateurs. »
\tcblower
Solution détaillée :
\begin{enumerate}
\item Consigne de règlement officiel → infinitif normatif : \all{Die Kabel nicht anfassen!} ou \all{Nicht an die Kabel fassen!} (verbe à particule \all{anfassen}, infinitif en finale).
\item Impératif pluriel + « avant de » → \all{bevor} + subordonnée, verbe en finale : \all{Lest die Nachrichten, bevor ihr sie kommentiert.} (verbe fort \all{lesen} : impératif pluriel \all{lest}).
\item Passé composé passif → prétérit du passif : \all{Die Reportage wurde von vielen Zuschauern gesehen.} (\all{werden} au prétérit : \all{wurde} ; agent avec \all{von} + Dativ).
\end{enumerate}
\end{exoresolu}

\begin{methode}[La relecture systématique : la check-list du candidat]
\begin{enumerate}
\item \textbf{Majuscules} : tous les noms allemands prennent une majuscule (\all{die Quelle}, \all{das Handy}).
\item \textbf{Ordre des mots} : verbe conjugué en 2\textsuperscript{e} position dans la principale, en 1\textsuperscript{re} à l'impératif, en position finale dans toute subordonnée (\all{bevor}, \all{weil}, \all{dass}...).
\item \textbf{Cas} : Akkusativ après \all{überprüfen}, \all{teilen} ; Dativ après \all{von}, \all{mit}, \all{vor} (localisation).
\item \textbf{Particules séparables} : renvoyées en fin de phrase (\all{schaltet ... aus}).
\item \textbf{Orthographe} : ß après voyelle longue (\all{die Straße}), ss après voyelle brève (\all{müssen}) ; Umlaute ä, ö, ü non oubliés.
\end{enumerate}
\end{methode}

\begin{exercice}[À vous de jouer : thème et version]
Thème : traduisez. 1) « Informez-vous avant de partager cette vidéo. » 2) « Il est interdit de photographier dans le studio. » 3) « La journaliste reçoit un prix pour son reportage. »\\
Version : traduisez. \all{„Bitte die Handys ausschalten! Die Sendung wird live übertragen. Überprüft eure Quellen!“}
\end{exercice}

\begin{corrige}[Exercice : thème et version]
Thème : 1) \all{Informiert euch, bevor ihr dieses Video teilt.} 2) \all{Es ist verboten, im Studio zu fotografieren.} (ou \all{Im Studio nicht fotografieren!}) 3) \all{Die Journalistin bekommt einen Preis für ihre Reportage.}\\
Version : « Veuillez éteindre les portables ! L'émission est retransmise en direct. Vérifiez vos sources ! »
\end{corrige}

\begin{savaistu}[Le thème au Baccalauréat A]
Au Baccalauréat A, le thème porte généralement sur 4 à 5 phrases courtes en lien direct avec le chapitre étudié — ici, les médias et la communication. Chaque erreur de cas, d'ordre des mots ou de particule est sanctionnée : une relecture méthodique de deux minutes peut vous faire gagner plusieurs points. Les correcteurs valorisent aussi les tournures idiomatiques : préférez \all{Es ist verboten, ... zu} à un calque maladroit du français. Entraînez-vous à traduire chaque semaine quelques consignes et slogans : c'est le meilleur investissement pour le jour de l'examen.
\end{savaistu}

\coursec{Production guidée : affiche, règlement, slogan}

Dans la section précédente, nous avons appris à \cle{traduire} des textes courts du français vers l'allemand (thème) et de l'allemand vers le français (version), en appliquant la méthode en cinq étapes : lire, repérer, analyser, traduire, relire. La traduction et la production écrite sont deux faces d'une même compétence : dans les deux cas, il faut fabriquer un texte allemand correct, clair et adapté à son destinataire. Nous allons maintenant passer à l'acte : \textbf{produire} nous-mêmes les trois grands types de textes injonctifs du chapitre « Média und Kommunikation » : l'\cle{affiche} (\all{das Plakat, -e}), le \cle{règlement} (\all{die Hausordnung, -en}, \all{die Regel, -n}) et le \cle{slogan} (\all{der Slogan, -s}, \all{der Werbespruch, -sprüche}).

\courssub{Observer d'abord : une affiche authentique}

Avant d'écrire, observons. Voici une affiche rédigée pour le club des médias d'un lycée de Yaoundé. Lisez-la attentivement et repérez ses différentes parties.

\begin{textebox}[Affiche du club médias du lycée]
\textbf{MEDIENCLUB JUGENDSTIMME}\\[0.3em]
\textbf{Deine Stimme zählt – informier dich, prüf nach, sprich mit!}\\[0.5em]
Du interessierst dich für Radio, Presse und soziale Netzwerke? Du willst lernen, wie man Nachrichten schreibt, Interviews führt und Fake News erkennt? Dann komm zu uns!\\[0.5em]
Im Medienclub lernst du:\\
-- Artikel für die Schülerzeitung zu schreiben\\
-- Interviews zu führen und zu filmen\\
-- Informationen zu prüfen, bevor du sie teilst\\[0.5em]
\textbf{Wann:} jeden Mittwoch, 15 Uhr\\
\textbf{Wo:} Raum 12, Gymnasium Yaoundé\\
\textbf{Kontakt:} Frau Ngo Bell, Lehrerin für Deutsch\\
\textbf{Eintritt frei – alle sind willkommen!}
\end{textebox}

\begin{definition}[Glossaire de l'affiche]
\all{die Jugendstimme, -n} : la voix des jeunes ; \all{die Stimme, -n} : la voix ; \all{zählen} : compter ; \all{die Schülerzeitung, -en} : le journal des élèves ; \all{führen} : mener, conduire ; \all{erkennen (erkannte, hat erkannt)} : reconnaître ; \all{die Fake News} (pl.) : les fausses informations ; \all{teilen} : partager ; \all{der Eintritt, -e} : l'entrée ; \all{willkommen} : bienvenu.
\end{definition}

\textbf{Questions d'observation :}
\begin{enumerate}
\item Où se trouve le titre ? Comment est-il présenté ?
\item Quelle phrase joue le rôle de slogan ? Quel mode verbal utilise-t-elle ?
\item Quelles informations pratiques sont données à la fin ?
\item Pourquoi le texte est-il découpé en courts paragraphes ?
\end{enumerate}

\courssub{Les caractéristiques de l'affiche (das Plakat)}

\begin{propriete}[Structure d'une affiche efficace]
Une affiche allemande comprend quatre blocs, toujours dans le même ordre :
\begin{enumerate}
\item \textbf{Le titre accrocheur} (\all{die Schlagzeile, -n}) : court, en gros caractères, souvent un nom fort ou une formule choc : \all{MEDIENCLUB JUGENDSTIMME}.
\item \textbf{Le slogan} (\all{der Slogan}) : une phrase à l'\cle{impératif} qui interpelle le lecteur : \all{Deine Stimme zählt – informier dich!}
\item \textbf{Le corps} : les consignes ou arguments, souvent en infinitifs avec \all{zu} : \all{Artikel zu schreiben, Interviews zu führen}.
\item \textbf{Les informations essentielles} : \all{Wann?} (quand), \all{Wo?} (où), \all{Wer?} (qui), \all{Kontakt} (contact).
\end{enumerate}
La mise en page doit être \cle{aérée} : peu de texte, des tirets, des blancs, des mots-clés en gras.
\end{propriete}

\begin{popfigure}
\begin{tikzpicture}[
  block/.style={rectangle, rounded corners=4pt, draw=popblue, fill=popblueL, text width=4.2cm, align=center, minimum height=1cm, font=\small},
  arr/.style={-{Stealth[length=3mm]}, thick, poporange}]
\node[block, fill=poppinkL, draw=poppink, align=center] (titre){\textbf{1. TITRE}\\ \all{die Schlagzeile}\\ accrocheur, en gras};
\node[block, below=0.5cm of titre, fill=popgoldL, draw=popgold, align=center] (slogan){\textbf{2. SLOGAN}\\ \all{der Slogan}\\ impératif, rythmé};
\node[block, below=0.5cm of slogan, fill=popgreenL, draw=popgreen, align=center] (corps){\textbf{3. CORPS}\\ consignes, arguments\\ infinitifs avec \all{zu}};
\node[block, below=0.5cm of corps, fill=poppurpleL, draw=poppurple, align=center] (infos){\textbf{4. INFOS}\\ \all{Wann? Wo? Wer? Kontakt?}};
\draw[arr] (titre) -- (slogan);
\draw[arr] (slogan) -- (corps);
\draw[arr] (corps) -- (infos);
\end{tikzpicture}
\legende{Structure d'une affiche : du titre accrocheur aux informations pratiques.}
\end{popfigure}

\begin{attention}[L'ordre des questions en allemand]
En français on dit « quoi, qui, quand, où » ; en allemand, les rubriques d'une affiche sont presque toujours introduites par \all{Wann:}, \all{Wo:}, \all{Wer:}, \all{Kontakt:}, suivies de deux-points, sans verbe. N'écrivez pas \all{Wann ist es?} sur une affiche : la rubrique télégraphique suffit.
\end{attention}

\courssub{Les caractéristiques du règlement (die Hausordnung)}

Le règlement est un texte injonctif \cle{impersonnel} : il ne s'adresse pas à un « tu » ou un « vous » précis, il énonce des normes valables pour tous. Observons un extrait du règlement d'une webradio scolaire.

\begin{exemplebox}[Un règlement de studio radio]
\all{1. Handys sind während der Sendung auszuschalten.}\\
« Les portables doivent être éteints pendant l'émission. »\\[0.3em]
\all{2. Nicht filmen!}\\
« Ne pas filmer ! »\\[0.3em]
\all{3. Es ist verboten, im Studio zu essen und zu trinken.}\\
« Il est interdit de manger et de boire dans le studio. »\\[0.3em]
\all{4. Das Mikrofon ist nach der Sendung auszuschalten.}\\
« Le micro doit être éteint après l'émission. »\\[0.3em]
\all{5. Es ist erlaubt, Musikwünsche per E-Mail zu schicken.}\\
« Il est permis d'envoyer ses souhaits musicaux par courriel. »
\end{exemplebox}

\begin{propriete}[Les trois formes de l'injonction impersonnelle]
\begin{enumerate}
\item \textbf{L'infinitif normatif} : infinitif seul (ou précédé de \all{nicht} ou d'un complément), placé en fin de phrase : \all{Nicht filmen!} « Ne pas filmer ! » C'est la forme la plus brève, typique des panneaux et consignes.
\item \textbf{La tournure \all{sein + zu + infinitif}} (passif d'obligation) : \all{Handys sind auszuschalten.} « Les portables doivent être éteints. » Attention : l'infinitif à particule se colle avec \all{zu} : \all{auszuschalten}, \all{herunterzufahren}.
\item \textbf{La tournure \all{Es ist + adjectif, ... zu + infinitif}} : \all{Es ist verboten/erlaubt/wichtig, ... zu ...} « Il est interdit/permis/important de... »
\end{enumerate}
Dans les trois cas, le ton reste \cle{neutre} et impersonnel : pas de \all{du}, pas de \all{Sie}.
\end{propriete}

\begin{center}
\begin{tabularx}{\linewidth}{|X|X|X|}
\hline
\rowcolor{popblueL} Français & Deutsch & Remarque \\
\hline
Ne pas fumer. & \all{Nicht rauchen!} & infinitif normatif \\
\hline
Il est interdit de filmer. & \all{Es ist verboten zu filmen.} & tournure impersonnelle \\
\hline
Les portables doivent être éteints. & \all{Handys sind auszuschalten.} & \all{sein + zu} + infinitif \\
\hline
Il est permis de poser des questions. & \all{Es ist erlaubt, Fragen zu stellen.} & virgule avant \all{zu} si groupe long \\
\hline
Interdiction d'entrer. & \all{Betreten verboten!} & nom + \all{verboten} \\
\hline
\end{tabularx}
\end{center}

\begin{attention}[Pas de « du » dans un règlement]
L'erreur classique du francophone est de traduire « vous devez éteindre » par \all{Du musst ausschalten} ou \all{Sie müssen ausschalten}. Dans un règlement allemand, on évite la personne : on préfère l'\cle{infinitif normatif} ou \all{Es ist + adjectif, ... zu}. Réservez \all{müssen} aux explications orales, pas au texte affiché.
\end{attention}

\courssub{Les caractéristiques du slogan (der Slogan)}

\begin{exemplebox}[Slogans de campagnes médiatiques]
\all{Schaltet ein – die Wahrheit wartet nicht!}\\
« Allumez vos postes – la vérité n'attend pas ! »\\[0.3em]
\all{Teile klug!}\\
« Partage intelligemment ! »\\[0.3em]
\all{Informier dich – prüf nach – denk mit!}\\
« Informe-toi – vérifie – réfléchis avec nous ! »\\[0.3em]
\all{Lesen macht stark.}\\
« Lire rend fort. »
\end{exemplebox}

\begin{propriete}[Rappel : les trois formes de l'impératif]
\begin{center}
\begin{tabular}{|l|l|l|}
\hline
\rowcolor{popblueL} Personne & Formation & Exemple \\
\hline
\all{du} & radical (+ \all{-e}) & \all{Informier dich!} / \all{Informiere dich!} \\
\hline
\all{ihr} & radical + \all{-t} & \all{Informiert euch!} \\
\hline
\all{Sie} & infinitif + \all{Sie} & \all{Informieren Sie sich!} \\
\hline
\end{tabular}
\end{center}
Exceptions : les verbes à vocalisme \all{e} qui devient \all{i} ou \all{ie} au radical prennent la voyelle changée à l'impératif \all{du}, sans \all{-e} : \all{Lies!} (lire), \all{Gib!} (donner), \all{Sprich!} (parler). Le verbe \all{sein} est irrégulier : \all{Sei ruhig!}, \all{Seid ruhig!}, \all{Seien Sie ruhig!}
\end{propriete}

\begin{methode}[Fabriquer un bon slogan allemand]
\begin{enumerate}
\item \textbf{Compter les mots} : un bon slogan fait \cle{3 à 7 mots}. Au-delà, il n'est plus mémorisable.
\item \textbf{Choisir le moteur} : soit un verbe à l'\cle{impératif} (\all{Teile klug!}), soit un \cle{nom fort} (\all{Lesen macht stark.}), soit un présent affirmatif.
\item \textbf{Créer un rythme} : répétition, triple impératif (\all{informier dich – prüf nach – denk mit!}), rime ou allitération (\all{klug} / \all{gut}).
\item \textbf{Soigner la chute} : le dernier mot doit frapper : \all{die Wahrheit wartet nicht!}
\item \textbf{Vérifier la grammaire} : impératif correct, majuscules aux noms, particule séparable en fin de phrase.
\end{enumerate}
\end{methode}

\begin{attention}[La majuscule, même dans un slogan]
En allemand, \textbf{tous les noms communs} prennent la majuscule, même dans un slogan ou un titre : \all{Teile deine Meinung!} et jamais \all{teile deine meinung}. De même, l'impératif d'un verbe réfléchi exige le pronom réfléchi : \all{Informier dich!} et non \all{Informier!}, car le verbe est \all{sich informieren}.
\end{attention}

\begin{savaistu}[Des campagnes réelles dans le monde germanophone]
En Allemagne, des campagnes comme \all{\#fairteilen} (« partageons équitablement ») encouragent le partage responsable sur les réseaux sociaux : vérifier une information avant de la diffuser. En Suisse, pays multilingue, les élèves créent chaque année des affiches \cle{bilingues} allemand-français pour les semaines des médias à l'école : un excellent modèle pour nos lycées camerounais, où l'allemand côtoie le français et l'anglais au quotidien !
\end{savaistu}

\courssub{Plan de production guidée : de l'idée à l'affiche}

Suivez ces quatre étapes en classe, en groupes de trois ou quatre.

\begin{methode}[Produire une affiche complète en quatre étapes]
\begin{enumerate}
\item \textbf{Brainstormer le lexique} : en cinq minutes, listez sur une feuille tous les mots du chapitre utiles au thème choisi (radio, réseaux sociaux, presse) : \all{die Nachricht, der Beitrag, senden, teilen, prüfen, melden, das Mikrofon, die Redaktion...}
\item \textbf{Rédiger trois slogans à l'impératif} : un pour la \all{du}-form, un pour la \all{ihr}-form, un pour la \all{wir}-form (\all{Informieren wir uns richtig!}). Choisissez ensuite le meilleur.
\item \textbf{Rédiger cinq articles de règlement} : deux à l'infinitif normatif (\all{Nicht filmen!}), deux avec \all{sein + zu} + infinitif (\all{Handys sind auszuschalten.}), un avec \all{Es ist verboten/erlaubt, ... zu}.
\item \textbf{Assembler l'affiche} : titre en gras, slogan en dessous, corps avec les arguments, rubriques \all{Wann / Wo / Wer / Kontakt} en bas. Aérez la mise en page.
\end{enumerate}
\end{methode}

\begin{exoresolu}[Rédiger un mini-règlement de salle informatique]
Énoncé : vous êtes responsable de la salle multimédia du lycée de Douala. Rédigez quatre articles de règlement : deux interdictions à l'infinitif normatif, une obligation avec \all{sein + zu}, une permission avec \all{Es ist erlaubt}. Thème : l'usage responsable d'Internet.
\tcblower
\textbf{Solution détaillée :}\\[0.3em]
\all{1. Nicht auf unbekannte Links klicken!}\\
« Ne pas cliquer sur des liens inconnus ! » (infinitif normatif négatif ; \all{klicken auf} + accusatif)\\[0.3em]
\all{2. Keine persönlichen Daten veröffentlichen!}\\
« Ne publier aucune donnée personnelle ! » (infinitif normatif avec \all{keine} ; \all{die Daten} = les données)\\[0.3em]
\all{3. Die Computer sind nach dem Unterricht herunterzufahren.}\\
« Les ordinateurs doivent être éteints après les cours. » (\all{sein + zu} + infinitif ; la particule séparable \all{herunterfahren} devient \all{herunterzufahren})\\[0.3em]
\all{4. Es ist erlaubt, für die Schularbeit im Internet zu recherchieren.}\\
« Il est permis de faire des recherches sur Internet pour le travail scolaire. » (tournure impersonnelle ; \all{die Schularbeit} = le devoir)\\[0.3em]
\textbf{Points de méthode :} aucun \all{du} ni \all{Sie} ; chaque article est numéroté ; les noms (\all{Links}, \all{Daten}, \all{Computer}, \all{Unterricht}, \all{Internet}) portent la majuscule ; les particules séparables sont bien en finale ou soudées à \all{zu}.
\end{exoresolu}

\courssub{La grille de relecture}

Avant de rendre votre production, relisez-la avec cette grille. C'est exactement celle que le correcteur du Bac applique mentalement.

\begin{center}
\begin{tabularx}{\linewidth}{|X|X|}
\hline
\rowcolor{popblueL} Question de relecture & Exemple de correction \\
\hline
Le verbe conjugué est-il en 2\up{e} position ? & \all{Jeden Mittwoch treffen wir uns.} (et non \all{Jeden Mittwoch wir treffen uns.}) \\
\hline
Tous les noms ont-ils une majuscule ? & \all{die Meinung}, \all{das Plakat} \\
\hline
Les cas après \all{von, mit, für} sont-ils corrects ? & \all{mit dem Handy} (datif), \all{für die Schüler} (accusatif) \\
\hline
La particule séparable est-elle en finale ? & \all{Schaltet ein!} (et non \all{Einschaltet!}) \\
\hline
L'impératif est-il bien formé ? & \all{Informier dich!} / \all{Informiert euch!} / \all{Informieren Sie sich!} \\
\hline
Le règlement reste-t-il impersonnel ? & \all{Es ist verboten...} et non \all{Du darfst nicht...} \\
\hline
\end{tabularx}
\end{center}

\begin{experience}[Atelier d'écriture en classe]
Par groupes de quatre : choisissez un thème (la webradio du lycée, la campagne contre les fausses informations, le festival des médias de Yaoundé). Suivez le plan en quatre étapes : brainstorm de lexique (5 min), trois slogans (10 min), cinq articles de règlement (10 min), assemblage de l'affiche sur une feuille (15 min). Chaque groupe présente son affiche à la classe ; les autres groupes la relisent avec la grille et proposent une correction. Votez ensuite pour le meilleur slogan de la classe.
\end{experience}

\courssub{Complément Bac : le paragraphe d'opinion}

\textbf{À retenir pour l'examen.} Au baccalauréat A, la production écrite demande souvent un texte de \cle{80 à 120 mots} sur un sujet du chapitre, par exemple : \all{Soziale Netzwerke – Fluch oder Segen?} (« Les réseaux sociaux : malédiction ou bénédiction ? »). Pour viser l'excellence, votre paragraphe doit contenir au moins \textbf{une phrase au passif} et \textbf{une phrase au discours indirect}.

\begin{exemplebox}[Modèle de paragraphe d'opinion (environ 100 mots)]
\all{Meiner Meinung nach sind soziale Netzwerke sehr nützlich, aber auch gefährlich. Viele Jugendliche werden jeden Tag von falschen Nachrichten beeinflusst. Ein Journalist hat letzte Woche gesagt, dass junge Leute ihre Informationen nicht mehr prüften. Deshalb müssen wir lernen, Quellen zu vergleichen. Ich finde es wichtig, dass Medienkompetenz in der Schule unterrichtet wird. Wer klug teilt, schützt sich und die anderen.}\\[0.3em]
« À mon avis, les réseaux sociaux sont très utiles, mais aussi dangereux. Beaucoup de jeunes sont influencés chaque jour par de fausses informations. Un journaliste a dit la semaine dernière que les jeunes ne vérifiaient plus leurs informations. C'est pourquoi nous devons apprendre à comparer les sources. Je trouve important que l'éducation aux médias soit enseignée à l'école. Celui qui partage intelligemment se protège et protège les autres. »
\end{exemplebox}

Repérez dans ce modèle : le passif (\all{werden beeinflusst}, \all{unterrichtet wird}), le discours indirect au Konjunktiv I (\all{hat gesagt, dass ... nicht mehr prüften} — ici le Konjunktiv I \all{prüften} a la même forme que le Präteritum, ce qui est fréquent ; on pourrait aussi écrire \all{prüfen würden}), les connecteurs d'opinion (\all{Meiner Meinung nach}, \all{Deshalb}, \all{Ich finde es wichtig, dass...}). Réutilisez ce squelette pour vos propres productions.

\begin{aretenir}[L'essentiel de la section]
\begin{itemize}
\item L'\cle{affiche} : titre accrocheur + slogan à l'impératif + corps aéré + rubriques \all{Wann / Wo / Wer / Kontakt}.
\item Le \cle{règlement} : articles numérotés, infinitif normatif (\all{Nicht filmen!}), \all{sein + zu} + infinitif, ou \all{Es ist verboten/erlaubt, ... zu} ; jamais de \all{du}.
\item Le \cle{slogan} : 3 à 7 mots, impératif ou nom fort, rythme, chute percutante, majuscules aux noms.
\item La \cle{relecture} : place du verbe, majuscules, cas, particule en finale.
\item Au Bac : 80-120 mots, avec une phrase au passif et une au discours indirect.
\end{itemize}
\end{aretenir}

\begin{exercice}[À vous de produire]
\begin{enumerate}
\item Rédigez trois slogans à l'impératif pour une campagne « Utilisons bien nos téléphones » au lycée.
\item Rédigez le règlement (5 articles) de la page Facebook du club d'allemand : deux infinitifs normatifs, deux tournures avec \all{sein + zu}, une avec \all{Es ist verboten}.
\item Assemblez une affiche complète pour la « Semaine des médias » de votre établissement (titre, slogan, trois arguments, date, lieu, contact).
\item Sujet Bac : rédigez un paragraphe de 80 à 120 mots : \all{Sollten Handys in der Schule verboten werden?} avec une phrase au passif et une au discours indirect.
\end{enumerate}
\end{exercice}

\newpage

\coursec{Activité d'intégration}

\courssub{Objectifs de l'activité}

À l'issue de cette activité d'intégration, tu seras capable de :

\begin{itemize}
\item rédiger en allemand une \cle{affiche} attractive pour un événement scolaire, avec un slogan à l'\cle{impératif} et des consignes claires ;
\item formuler les articles d'un \cle{règlement} en utilisant l'infinitif normatif et le \cle{passif} ;
\item traduire du français vers l'allemand (thème) des phrases contenant du \cle{discours indirect} au Konjunktiv I ;
\item traduire de l'allemand vers le français (version) un court texte injonctif authentique ;
\item relire et corriger ta production à l'aide d'une grille : ordre des mots, place du verbe, cas, majuscules des noms, orthographe.
\end{itemize}

\courssub{Situation}

Ton lycée, le \all{Gymnasium Yaoundé}, organise en partenariat avec le \all{Heinrich-Heine-Gymnasium} de Hambourg une semaine des médias, la \all{Medienwoche}. Le club des jeunes reporters (\all{der Jugendreporter-Club}), dont tu es membre, est chargé de la communication. Le proviseur a annoncé que la semaine commencerait lundi et que chaque classe serait représentée par un groupe de quatre élèves. Votre mission : produire une affiche bilingue, un mini-règlement du club et des traductions qui seront envoyées aux partenaires allemands.

Voici le courriel officiel reçu de Hambourg :

\begin{textebox}[E-Mail aus Hamburg]
Liebe Freunde aus Yaoundé,\\
wir freuen uns sehr auf unsere gemeinsame Medienwoche! Bitte schickt uns bis Freitag euer Plakat mit einem Slogan und vier Regeln für den Workshop. Denkt daran: Kurze Sätze wirken stärker als lange! Unsere Reporter-AG wird eure Texte lesen und kommentieren.\\
Herzliche Grüße\\
Frau Berger, Leiterin der Reporter-AG
\end{textebox}

\begin{definition}[Glossaire du courriel]
\all{das Plakat, -e} : l'affiche ; \all{der Slogan, -s} : le slogan ; \all{die Regel, -n} : la règle ; \all{der Workshop, -s} : l'atelier ; \all{wirken} : produire un effet ; \all{die Reporter-AG} : le club des reporters (\all{AG} = \all{Arbeitsgemeinschaft}, club scolaire) ; \all{der Gruß, ⸚e} : la salutation ; \all{gemeinsam} : commun, en commun ; \all{der Leiter, - / die Leiterin, -nen} : le/la responsable.
\end{definition}

\courssub{Consignes}

Par groupes de quatre, vous réalisez les cinq tâches suivantes. Répartissez les rôles : un secrétaire, un correcteur, un traducteur, un rapporteur.

\textbf{Tâche 1 — L'affiche (production écrite).} Rédigez en allemand une affiche pour la \all{Medienwoche}. Elle doit contenir : un titre accrocheur, un \cle{slogan à l'impératif} (2\up{e} personne du pluriel), la date et le lieu, et \cle{quatre consignes} pour les participants (au moins deux à l'infinitif normatif, ex. \all{Handys ausschalten!}).

\textbf{Tâche 2 — Le mini-règlement (production écrite).} Rédigez \cle{trois articles} du règlement du club des jeunes reporters. Chaque article commence par \all{Artikel 1 : …}. Au moins une phrase doit être au \cle{passif} (ex. \all{Die Artikel werden vor der Veröffentlichung korrigiert.}).

\textbf{Tâche 3 — Thème (français → allemand).} Traduisez en allemand les trois phrases suivantes, dictées par l'enseignant :
\begin{enumerate}
\item Le directeur a annoncé que la semaine des médias commencerait lundi.
\item On nous a demandé d'éteindre nos téléphones pendant l'atelier.
\item Les articles seront publiés dans le journal du lycée.
\end{enumerate}

\textbf{Tâche 4 — Version (allemand → français).} Traduisez en français cet extrait d'affiche du lycée partenaire :

\begin{textebox}[Ausschnitt aus dem Hamburger Plakat]
„Medien machen Meinung!“ — Kommt am Montag um 9 Uhr in die Aula! Filmt eure Interviews nur mit Erlaubnis. Alle Beiträge werden von unserer Redaktion geprüft, bevor sie online gestellt werden. Seid kritisch, aber fair!
\end{textebox}

\textbf{Tâche 5 — Relecture et vote.} Affichez votre production au tableau. La classe vote pour le meilleur slogan. Puis chaque groupe relit et corrige son texte avec la grille : verbe en 2\up{e} position, particule séparable en fin de phrase, cas après préposition, majuscule à tous les noms, Konjunktiv I au discours indirect.

\courssub{Ressources à mobiliser}

\begin{aretenir}[Boîte à outils grammaticale]
\begin{itemize}
\item \cle{Impératif} (2\up{e} pers. pluriel) : radical + \all{-t} → \all{Kommt! Informiert euch! Seid kritisch!} (verbe en 1\up{re} position).
\item \cle{Infinitif normatif} (consignes impersonnelles) : infinitif à la fin → \all{Handys ausschalten! Nicht ohne Erlaubnis filmen!}
\item \cle{Passif} (Präsens) : sujet + \all{werden} conjugué + participe II en fin → \all{Die Beiträge werden geprüft.} ; au Futur : \all{werden} + participe II + \all{werden} → \all{Die Artikel werden veröffentlicht werden.}
\item \cle{Discours indirect} : Konjunktiv I → \all{Der Direktor sagte, die Woche beginne am Montag.} (si le KI est identique à l'indicatif, on prend le KII : \all{sie begänne}).
\item Subordonnée : verbe conjugué en dernière position → \all{…, bevor sie online gestellt werden.}
\end{itemize}
\end{aretenir}

\begin{propriete}[Rappel : les trois formes de l'impératif]
\begin{center}
\begin{tabular}{|l|l|l|}
\hline
\rowcolor{popblueL} Personne & Formation & Exemple \\
\hline
\all{du} (2\up{e} pers. sg.) & radical (+\all{-e}) & \all{Komm pünktlich!} \\
\hline
\all{ihr} (2\up{e} pers. pl.) & radical + \all{-t} & \all{Kommt pünktlich!} \\
\hline
\all{Sie} (politesse) & infinitif + \all{Sie} & \all{Kommen Sie pünktlich!} \\
\hline
\end{tabular}
\end{center}
Exceptions : \all{sein} → \all{Sei! Seid! Seien Sie!} ; les verbes à vocalisme \all{e→i/ie} prennent la voyelle changée au singulier sans \all{-e} : \all{Lies! Nimm! Gib!} (mais \all{ihr} garde le \all{e} : \all{Lest! Nehmt!}).
\end{propriete}

\begin{propriete}[Rappel : le Konjunktiv I du discours indirect]
Le Konjunktiv I se forme sur le radical de l'infinitif : \all{ich beginne, du beginnest, er/sie/es beginne, wir beginnen, ihr beginnet, sie beginnen}. Seul \all{sein} est irrégulier : \all{ich sei, du seist, er sei, wir seien, ihr seiet, sie seien}. Quand la forme du KI est identique à l'indicatif présent (ex. \all{sie beginnen}), on la remplace par le Konjunktiv II : \all{sie begännen}, \all{sie wären}. Exemple : \all{Der Direktor sagte, die Woche beginne am Montag und jede Klasse sei vertreten.}
\end{propriete}

\begin{definition}[Lexique utile]
\all{die Medienwoche, -n} : la semaine des médias ; \all{die Aula, Aulen} : la salle de fêtes ; \all{die Erlaubnis, -se} : l'autorisation ; \all{der Beitrag, ⸚e} : la contribution, l'article ; \all{die Redaktion, -en} : la rédaction ; \all{prüfen} : vérifier ; \all{veröffentlichen} : publier ; \all{die Veröffentlichung, -en} : la publication ; \all{die Meinung, -en} : l'opinion ; \all{kritisch} : critique ; \all{fair} : équitable ; \all{ausschalten (schaltet aus, schaltete aus, hat ausgeschaltet)} : éteindre ; \all{ankündigen} : annoncer ; \all{die Schülerzeitung, -en} : le journal des élèves ; \all{pünktlich} : ponctuel, à l'heure.
\end{definition}

\begin{attention}[Pièges fréquents]
Ne traduis pas mot à mot ! « Éteignez vos téléphones » = \all{Schaltet eure Handys aus!} (particule \all{aus-} en fin de phrase, pas \all{*ausschaltet eure Handys}). Au discours indirect, « commencerait » (conditionnel français) se rend par le Konjunktiv : \all{beginne} ou \all{begänne} ; la forme \all{würde beginnen} existe dans la langue courante, mais elle est à éviter dans un style soigné et à l'examen. Attention aussi au futur passif : « seront publiés » = \all{werden … veröffentlicht werden} (deux fois \all{werden} !). Et n'oublie pas la majuscule : \all{die Medienwoche}, \all{das Handy}.
\end{attention}

\courssub{Proposition de production et corrigé commenté}

\begin{exoresolu}[Modèle de production du groupe 1]
\textbf{Affiche :}\\
„MEDIENWOCHE AM GYMNASIUM YAOUNDÉ — Entdecke die Welt der Nachrichten!“\\
Montag bis Freitag, Aula, 9 Uhr. Mit dem Heinrich-Heine-Gymnasium Hamburg.\\
1. Kommt pünktlich! 2. Handys ausschalten! 3. Stellt kritische Fragen! 4. Nicht ohne Erlaubnis filmen!\\[0.3em]
\textbf{Règlement :}\\
Artikel 1 : Jeder Reporter schreibt einen Artikel pro Woche.\\
Artikel 2 : Alle Artikel werden vor der Veröffentlichung von der Redaktion korrigiert.\\
Artikel 3 : Interviews dürfen nur mit Erlaubnis gefilmt werden.\\[0.3em]
\textbf{Thème :}\\
1. Der Direktor hat angekündigt, dass die Medienwoche am Montag beginne.\\
2. Man hat uns gebeten, unsere Handys während des Workshops auszuschalten.\\
3. Die Artikel werden in der Schülerzeitung veröffentlicht werden.\\[0.3em]
\textbf{Version :}\\
« Les médias font l'opinion ! » — Venez lundi à 9 h à la salle de fêtes ! Ne filmez vos interviews qu'avec autorisation. Toutes les contributions sont vérifiées par notre rédaction avant d'être mises en ligne. Soyez critiques, mais justes !
\tcblower
\textbf{Commentaire des choix de langue :}
\begin{itemize}
\item \textbf{Slogan à l'impératif :} \all{Entdecke die Welt der Nachrichten!} — impératif singulier (radical de \all{entdecken} + \all{-e}), verbe en 1\up{re} position, ton direct et accrocheur. Le titre joue sur l'assonance \all{Medien / Nachrichten}.
\item \textbf{Consignes variées :} deux impératifs pluriels (\all{Kommt pünktlich!}, \all{Stellt kritische Fragen!}) et deux infinitifs normatifs impersonnels (\all{Handys ausschalten!}, \all{Nicht ohne Erlaubnis filmen!}) : c'est exactement le mélange exigé par une affiche authentique. Note la particule séparable \all{aus-} rejetée en fin de consigne.
\item \textbf{Passif dans le règlement :} \all{Alle Artikel werden … korrigiert} (Präsens passif : \all{werden} + participe II) ; \all{dürfen … gefilmt werden} (passif avec modal : infinitif \all{werden} en toute fin). L'agent est introduit par \all{von} + datif (\all{von der Redaktion}).
\item \textbf{Discours indirect :} « commencerait » est rendu par le Konjunktiv I \all{beginne} dans la subordonnée introduite par \all{dass}, verbe en dernière position. La forme \all{begänne} (KII) serait aussi correcte, car \all{beginne} est identique à l'indicatif.
\item \textbf{Thème, phrase 2 :} « On nous a demandé de… » devient \all{Man hat uns gebeten, … auszuschalten} : proposition infinitive avec \all{zu} (\all{auszuschalten} : le \all{zu} s'insère entre la particule et le radical). \all{während des Workshops} : génitif après \all{während}.
\item \textbf{Thème, phrase 3 :} « seront publiés » est un futur passif : \all{werden} (auxiliaire de futur, conjugué en 2\up{e} position) + participe II \all{veröffentlicht} + infinitif \all{werden} (auxiliaire de passif) en toute fin. Ne pas oublier le second \all{werden} !
\item \textbf{Version :} \all{Medien machen Meinung} est rendu par une formule française naturelle, « Les médias font l'opinion », et non par un calque ; \all{bevor sie online gestellt werden} (passif en subordonnée) devient « avant d'être mises en ligne » : le français préfère l'infinitif passif.
\end{itemize}
\end{exoresolu}

\courssub{Bilan de l'activité}

Cette tâche d'intégration t'a permis de réinvestir l'ensemble du chapitre dans une situation de communication authentique :

\begin{itemize}
\item \textbf{Production :} tu as rédigé une affiche et un règlement, deux genres injonctifs typiques de la communication scolaire, en choisissant à bon escient entre impératif et infinitif normatif.
\item \textbf{Grammaire :} tu as mobilisé l'impératif, le passif (présent, futur, avec modal) et le discours indirect au Konjunktiv I.
\item \textbf{Traduction :} tu as pratiqué le thème (attention au conditionnel français rendu par le Konjunktiv et au futur passif à double \all{werden}) et la version (équivalences naturelles plutôt que calques).
\item \textbf{Relecture :} la grille de correction (place du verbe, particules séparables, cas, majuscules) t'a rendu autonome dans l'amélioration de ton texte.
\end{itemize}

\begin{methode}[Pour aller plus loin]
Pour préparer l'évaluation : 1) choisis une affiche réelle d'un établissement germanophone et traduis-la ; 2) réécris ton règlement en le mettant entièrement au passif ; 3) transforme le courriel de Frau Berger en discours indirect (\all{Frau Berger schrieb, dass …}). Ces trois exercices reprennent exactement les compétences du chapitre.
\end{methode}

\newpage

\coursec{Méthodes et exercices résolus}

\courssub{Savoir-faire 1 : Rédiger un texte injonctif (affiche, règlement, consignes, slogan)}

\begin{methode}[Rédiger une affiche, un règlement ou un slogan]
\begin{enumerate}
\item \textbf{Identifier le type de texte} : une \cle{affiche} (das Plakat, die Plakate) informe et invite ; un \cle{règlement} (die Ordnung, die Ordnungen ; das Verbot, die Verbote) interdit ou oblige ; un \cle{slogan} (der Slogan, die Slogans) est court, rythmé, mémorable.
\item \textbf{Choisir la forme d'injonction adaptée} :
\begin{itemize}
\item impératif 2\ieme{} personne du pluriel : \all{Informiert euch!} (informez-vous !) ;
\item infinitif d'injonction (panneaux, règlements) : \all{Nicht rauchen!} (ne pas fumer !), \all{Handys ausschalten!} (éteindre les portables !) ;
\item tournure impersonnelle : \all{Das Rauchen ist verboten.} (il est interdit de fumer), \all{Man muss ...} (il faut ...) ;
\item passif avec \all{werden} : \all{Fotos werden nicht geduldet.} (les photos ne sont pas tolérées).
\end{itemize}
\item \textbf{Structurer} : titre accrocheur en majuscules ou en gros caractères, 3 à 6 consignes courtes, une phrase de conclusion (contact, date, lieu).
\item \textbf{Soigner le slogan} : phrase brève, rythme, éventuelle rime : \all{Lesen lohnt sich!} (lire en vaut la peine !).
\item \textbf{Relire} : majuscule à tous les noms, verbe conjugué en position 2, particule séparable rejetée en fin de phrase à l'impératif : \all{Schaltet eure Handys aus!} (éteignez vos portables !).
\end{enumerate}
\end{methode}

\begin{propriete}[Rappel : les trois formes de l'impératif]
\begin{center}
\begin{tabular}{|l|l|l|}
\hline
\rowcolor{popblueL} Personne & Formation & Exemple \\
\hline
\all{du} & radical (+ \all{-e} facultatif) & \all{Komm! Lies den Artikel!} \\
\hline
\all{ihr} & radical + \all{-t} (= présent sans \all{ihr}) & \all{Kommt! Lest den Artikel!} \\
\hline
\all{Sie} & infinitif + \all{Sie} & \all{Kommen Sie! Lesen Sie den Artikel!} \\
\hline
\end{tabular}
\end{center}
Attention : les verbes forts à vocalisme \all{e} qui change en \all{i/ie} au présent gardent ce changement à l'impératif \all{du} : \all{lesen} $\rightarrow$ \all{Lies!}, \all{sprechen} $\rightarrow$ \all{Sprich!}. En revanche, pas d'Umlaut à l'impératif : \all{fahren} $\rightarrow$ \all{Fahr!} (et non \all{Fähr!}).
\end{propriete}

\begin{exoresolu}[Rédaction d'un règlement de salle multimédia]
\textbf{Énoncé (type Bac)} : Vous êtes délégué(e) de classe. Rédigez en allemand le règlement (5 consignes minimum) de la salle multimédia de votre lycée à Yaoundé. Utilisez l'infinitif d'injonction ou l'impératif.
\tcblower
\textbf{Raisonnement} : un règlement allemand utilise typiquement l'\emph{infinitif d'injonction} (verbe à la fin, sans sujet) ou l'impératif pluriel. Les verbes à particule séparable restent soudés à l'infinitif : \all{ausschalten} $\rightarrow$ \all{Handys ausschalten!}. On évite les phrases longues : une consigne = une ligne.

\textbf{Réponse rédigée} :

\textbf{ORDNUNG FÜR DEN MEDIENRAUM}

\all{1. Handys ausschalten!} (éteindre les portables !)

\all{2. Nicht essen und nicht trinken!} (ne pas manger ni boire !)

\all{3. Leise sprechen!} (parler doucement !)

\all{4. Die Computer nach der Arbeit herunterfahren!} (éteindre les ordinateurs après le travail !)

\all{5. Den Raum sauber verlassen!} (laisser la salle propre !)

\all{Der Lehrer / Die Lehrerin}

\textbf{Justifications grammaticales} : à l'infinitif d'injonction, l'infinitif se place en \emph{dernière position} ; les verbes séparables restent soudés à l'infinitif (\all{ausschalten}, \all{herunterfahren}) ; \all{verlassen} est inséparable (préfixe \all{ver-}) et prend ici un accusatif : \all{den Raum}. Tous les noms prennent la majuscule : \all{Handys}, \all{Computer}, \all{Raum}. La négation d'une action générale se fait avec \all{nicht} (pas \all{kein}, réservé à la négation d'un nom).

\textbf{Conclusion} : le texte est court, impersonnel, chaque consigne tient sur une ligne : c'est exactement le format attendu d'un règlement allemand.
\end{exoresolu}

\courssub{Savoir-faire 2 : Traduire en version (allemand $\rightarrow$ français) un texte sur les médias}

\begin{methode}[Réussir une version]
\begin{enumerate}
\item \textbf{Lire tout le texte} une première fois pour dégager le thème (ici : les médias, la communication) et le temps dominant.
\item \textbf{Segmenter chaque phrase} : repérer le verbe conjugué (position 2), le sujet, puis les compléments ; attention au verbe renvoyé en fin de proposition subordonnée.
\item \textbf{Identifier les pièges} : verbes séparables (\all{schaltet ... ein}), passif (\all{wird gedruckt} = est imprimé), discours indirect au Konjunktiv I (\all{er sagte, er habe ...} = il a dit qu'il avait ...), faux amis (\all{aktuell} = actuel, pas « actif » ; \all{die Zeitung} = le journal).
\item \textbf{Traduire en bon français} : pas de calque de l'ordre allemand ; le français remet le verbe près du sujet.
\item \textbf{Relire} : vérifier qu'aucun mot allemand n'a été oublié et que le français est naturel.
\end{enumerate}
\end{methode}

\begin{exoresolu}[Version : un court article de presse]
\textbf{Énoncé (type Bac)} : Traduisez en français le texte suivant.

\all{Immer mehr Jugendliche in Kamerun informieren sich über das Handy. Eine Lehrerin aus Yaoundé sagte, ihre Schüler läsen kaum noch Zeitungen, weil die Nachrichten im Internet schneller verbreitet würden. Trotzdem wird das Radio auf dem Land noch oft eingeschaltet.}
\tcblower
\textbf{Raisonnement mot à mot} :
\begin{itemize}
\item \all{Immer mehr Jugendliche} = de plus en plus de jeunes (\all{immer mehr} = de plus en plus) ;
\item \all{informieren sich über das Handy} = s'informent par le portable (verbe réfléchi \all{sich informieren}) ;
\item \all{ihre Schüler läsen kaum noch Zeitungen} : discours indirect ; le Konjunktiv I de \all{lesen} à la 3\ieme{} personne du pluriel (\all{lasen}) serait identique à l'indicatif, on emploie donc le Konjunktiv II \all{läsen} ; le français rend le discours rapporté par l'indicatif après « dire que » : ses élèves ne lisaient presque plus de journaux ;
\item \all{weil die Nachrichten ... verbreitet würden} : passif au Konjunktiv II (discours rapporté) dans une subordonnée \all{weil}, verbe en fin = parce que les informations étaient diffusées plus rapidement sur Internet ;
\item \all{wird ... eingeschaltet} : passif présent, verbe séparable \all{einschalten} = on allume encore souvent la radio / la radio est encore souvent allumée ;
\item \all{auf dem Land} = à la campagne (faux ami : pas « sur le terrain »).
\end{itemize}

\textbf{Traduction de référence} : « De plus en plus de jeunes au Cameroun s'informent par le téléphone portable. Une enseignante de Yaoundé a déclaré que ses élèves ne lisaient presque plus de journaux, parce que les informations étaient diffusées plus rapidement sur Internet. Pourtant, à la campagne, on allume encore souvent la radio. »

\textbf{Conclusion} : les trois difficultés (discours indirect, passif, verbe séparable en fin de phrase) ont été repérées avant de traduire ; le français final est fluide et complet.
\end{exoresolu}

\courssub{Savoir-faire 3 : Traduire en thème (français $\rightarrow$ allemand)}

\begin{methode}[Réussir un thème]
\begin{enumerate}
\item \textbf{Analyser la phrase française} : sujet, verbe, temps, compléments ; ne jamais traduire mot à mot.
\item \textbf{Poser le squelette allemand} : verbe conjugué en position 2 ; si subordonnée, verbe conjugué à la fin.
\item \textbf{Choisir les cas} : après \all{mit} (+ datif), \all{für, ohne} (+ accusatif), \all{über} (+ accusatif pour un thème), vérifier l'article et l'adjectif.
\item \textbf{Traiter les verbes séparables} : à l'indicatif, la particule va en fin de phrase : \all{Il allume la radio} $\rightarrow$ \all{Er schaltet das Radio ein.}
\item \textbf{Choisir le temps} : récit au passé $\rightarrow$ Perfekt (oral) ou Präteritum (\all{sein, haben} et les verbes modaux s'emploient de préférence au Präteritum).
\item \textbf{Relire} : majuscules des noms, Umlaute, accord sujet-verbe, place du verbe.
\end{enumerate}
\end{methode}

\begin{exoresolu}[Thème : phrases sur les médias]
\textbf{Énoncé (type Bac)} : Traduisez en allemand.

1. Allumez la télévision, le journal commence !

2. Le présentateur a dit que les nouvelles seraient diffusées à vingt heures.

3. On ne doit pas croire tout ce qu'on lit sur Internet.
\tcblower
\textbf{Phrase 1} : impératif pluriel \all{ihr} (consigne collective) + verbe séparable \all{einschalten} : la particule \all{ein} va en fin de phrase. « Le journal » télévisé = \all{die Nachrichten} (pluriel) ; « la télévision » (l'appareil) = \all{den Fernseher} (accusatif masculin).

\all{Schaltet den Fernseher ein, die Nachrichten beginnen!}

\textbf{Phrase 2} : discours indirect au passé $\rightarrow$ subordonnée \all{dass} avec le verbe conjugué en fin ; pour le futur rapporté, on emploie \all{würden} + infinitif (forme standard du futur au discours indirect). \all{aussenden} (diffuser) est séparable : participe II \all{ausgesendet}.

\all{Der Moderator sagte, dass die Nachrichten um zwanzig Uhr ausgesendet würden.}

\textbf{Phrase 3} : « on » = \all{man} ; « tout ce que » = \all{alles, was} (relatif \all{was} obligatoire après \all{alles}) ; subordonnée relative $\rightarrow$ verbe \all{liest} en fin ; modal \all{darf} + négation \all{nicht} ; « sur Internet » = \all{im Internet}.

\all{Man darf nicht alles glauben, was man im Internet liest.}

\textbf{Conclusion} : chaque phrase illustre un réflexe du thème : particule en fin à l'impératif, \all{würden} + infinitif pour le futur rapporté, verbe final dans la relative.
\end{exoresolu}

\courssub{Savoir-faire 4 : Transformer des phrases (actif/passif, direct/indirect, impératif)}

\begin{methode}[Réussir les transformations grammaticales]
\begin{enumerate}
\item \textbf{Actif $\rightarrow$ passif} : le complément d'objet direct devient sujet ; \all{werden} au temps de la phrase active + participe II ; l'ancien sujet devient \all{von} + datif (ou disparaît). \all{Man druckt die Zeitung.} $\rightarrow$ \all{Die Zeitung wird gedruckt.}
\item \textbf{Direct $\rightarrow$ indirect} : introduire \all{dass} + Konjunktiv I (ou Konjunktiv II / \all{würde} + infinitif si le Konjunktiv I est identique à l'indicatif) ; adapter les pronoms et les marqueurs temporels.
\item \textbf{Phrase $\rightarrow$ impératif} : repérer la personne visée (\all{du / ihr / Sie}) ; radical pour \all{du} (\all{Lies!}), radical + \all{-t} pour \all{ihr} (\all{Lest!}), infinitif + \all{Sie} (\all{Lesen Sie!}).
\item \textbf{Infinitif d'injonction} : supprimer le sujet et mettre l'infinitif à la fin : \all{Ihr dürft hier nicht fotografieren.} $\rightarrow$ \all{Hier nicht fotografieren!}
\item \textbf{Vérifier} la place du verbe après chaque transformation.
\end{enumerate}
\end{methode}

\begin{exoresolu}[Transformations sur le thème des médias]
\textbf{Énoncé (type Bac)} : Transformez les phrases selon la consigne entre parenthèses.

1. \all{Die Journalisten schreiben die Artikel.} (au passif)

2. \all{Der Chefredakteur sagte: „Wir drucken die Zeitung morgen.“} (au discours indirect)

3. \all{Ihr sollt eure Handys ausschalten.} (à l'impératif)
\tcblower
\textbf{1. Passif} : COD \all{die Artikel} devient sujet (nominatif pluriel) ; présent de \all{werden} : \all{werden} ; participe II de \all{schreiben} : \all{geschrieben} ; agent introduit par \all{von} + datif : \all{von den Journalisten}.

\all{Die Artikel werden von den Journalisten geschrieben.}

\textbf{2. Discours indirect} : \all{sagen} + \all{dass} ; le Konjunktiv I de \all{drucken} à la 3\ieme{} personne du pluriel serait identique à l'indicatif (\all{drucken}), on prend donc la forme de substitution \all{würden drucken} ; \all{wir} devient \all{sie} ; le verbe conjugué passe en fin de subordonnée ; \all{morgen} devient \all{am nächsten Tag} (comme en français « demain » $\rightarrow$ « le lendemain »).

\all{Der Chefredakteur sagte, dass sie die Zeitung am nächsten Tag drucken würden.}

\textbf{3. Impératif} : la personne visée est \all{ihr} $\rightarrow$ impératif en \all{-t} ; verbe séparable \all{ausschalten} : la particule \all{aus} va en fin de phrase.

\all{Schaltet eure Handys aus!}

\textbf{Conclusion} : à chaque transformation, une seule chose change (la construction), le sens reste identique ; la place du verbe est le critère de correction décisif.
\end{exoresolu}

\courssub{Savoir-faire 5 : Relire et corriger sa production}

\begin{methode}[La relecture en 5 points de contrôle]
\begin{enumerate}
\item \textbf{Majuscules} : tout nom allemand prend une majuscule, même en pleine phrase : \all{das Handy, die Zeitung, der Artikel}.
\item \textbf{Place du verbe} : position 2 dans la phrase principale, position finale dans toute subordonnée (\all{weil, dass, wenn, was ...}).
\item \textbf{Cas} : vérifier chaque déterminant après préposition (\all{mit dem, für den, im Internet}) et chaque COD à l'accusatif (\all{den Artikel}).
\item \textbf{Orthographe} : Umlaute (\all{ä, ö, ü}), \all{ß} après voyelle longue (\all{die Straße}), particules séparables soudées à l'infinitif (\all{ausschalten}).
\item \textbf{Faux amis et calques} : « actuel » = \all{aktuell} ; « le journal » = \all{die Zeitung} ; « regarder la télé » = \all{fernsehen} (séparable : \all{Ich sehe abends fern.}).
\end{enumerate}
\end{methode}

\begin{exoresolu}[Correction d'une production d'élève]
\textbf{Énoncé (type Bac)} : La production suivante d'un élève contient 5 erreurs. Repérez-les et corrigez-les.

\all{1. die schüler lesen die zeitung nicht, weil sie haben keine Zeit. 2. Schaltet aus eure Handys! 3. Ich informiere mich auf dem Internet.}
\tcblower
\textbf{Erreurs 1 et 2 — majuscules} : \all{die schüler} $\rightarrow$ \all{die Schüler} ; \all{die zeitung} $\rightarrow$ \all{die Zeitung}. Tout nom allemand prend la majuscule.

\textbf{Erreur 3 — place du verbe dans la subordonnée} : \all{weil sie haben keine Zeit} $\rightarrow$ \all{weil sie keine Zeit haben}. Après \all{weil}, le verbe conjugué va en dernière position.

\textbf{Erreur 4 — particule séparable} : \all{Schaltet aus eure Handys!} $\rightarrow$ \all{Schaltet eure Handys aus!} À l'impératif, la particule \all{aus} se détache et va en fin de phrase.

\textbf{Erreur 5 — préposition calquée du français} : \all{auf dem Internet} $\rightarrow$ \all{im Internet} (= \all{in dem Internet}). « Sur Internet » se dit \all{im Internet}.

\textbf{Phrases corrigées} : \all{Die Schüler lesen die Zeitung nicht, weil sie keine Zeit haben.} — \all{Schaltet eure Handys aus!} — \all{Ich informiere mich im Internet.}

\textbf{Conclusion} : une relecture méthodique (majuscules $\rightarrow$ verbe $\rightarrow$ cas $\rightarrow$ orthographe $\rightarrow$ faux amis) permet de récupérer facilement un à deux points sur la copie.
\end{exoresolu}

\begin{attention}[Les 5 erreurs qui coûtent des points au Bac]
\begin{enumerate}
\item \textbf{Oublier la majuscule aux noms} : \all{die zeitung} est fautif ; écrivez \all{die Zeitung}. C'est la faute la plus sanctionnée en production écrite.
\item \textbf{Laisser le verbe en position 2 après \all{weil} ou \all{dass}} : on ne dit pas \all{weil sie hat keine Zeit} mais \all{weil sie keine Zeit hat}. Le verbe conjugué va toujours en fin de subordonnée.
\item \textbf{Mal placer la particule séparable} : \all{Schaltet aus eure Handys!} est fautif ; la particule va en fin de phrase : \all{Schaltet eure Handys aus!}
\item \textbf{Traduire les faux amis au pied de la lettre} : « actuel » = \all{aktuell} (\all{aktiv} = actif) ; « le journal » = \all{die Zeitung} ; « une nouvelle » = \all{eine Nachricht} ; « assister à » = \all{an etwas teilnehmen} (\all{assistieren} = aider).
\item \textbf{Calquer les prépositions françaises} : « sur Internet » = \all{im Internet} (pas \all{auf dem Internet}) ; « à la radio » = \all{im Radio} ; « à la télévision » = \all{im Fernsehen}. Apprenez ces combinaisons figées par cœur.
\end{enumerate}
\end{attention}

\newpage

\coursec{Exercices}

\courssub{Je vérifie mes connaissances}

\begin{exercice}[QCM : grammaire des textes injonctifs]
Pour chaque question, une seule réponse est correcte. Entoure la lettre correspondante.
\begin{enumerate}
\item L'impératif à la 2\ieme{} personne du singulier de \all{lesen} est :
\begin{enumerate}
\item[a)] \all{les!} \quad b) \all{lies!} \quad c) \all{lest!} \quad d) \all{lesen Sie!}
\end{enumerate}
\item Au discours indirect, « \all{Die Zeitung erscheint täglich.} » devient :
\begin{enumerate}
\item[a)] \all{Er sagte, die Zeitung erscheint täglich.}
\item[b)] \all{Er sagte, die Zeitung erscheine täglich.}
\item[c)] \all{Er sagte, die Zeitung erschien täglich.}
\end{enumerate}
\item « \all{Der Artikel wird von der Redaktion geschrieben.} » est une phrase :
\begin{enumerate}
\item[a)] au passif Präsens \quad b) au passif Präteritum \quad c) au discours indirect
\end{enumerate}
\item Pour interdire poliment dans un règlement, on écrit :
\begin{enumerate}
\item[a)] \all{Handys benutzt nicht!}
\item[b)] \all{Benutzen Sie bitte keine Handys!}
\item[c)] \all{Handys nicht benutzen Sie!}
\end{enumerate}
\end{enumerate}
\end{exercice}

\begin{exercice}[Vrai ou faux ? Justifie ta réponse]
Dis si les affirmations suivantes sont vraies (V) ou fausses (F) et corrige les phrases fausses.
\begin{enumerate}
\item À l'impératif, le verbe se place toujours en première position de la phrase.
\item \all{„Sei ruhig!“} est l'impératif correct du verbe \all{sein} à la 2\ieme{} personne du singulier.
\item Au passif Präteritum, on utilise l'auxiliaire \all{wurde} + participe II.
\item Au discours indirect, le verbe conjugué reste en deuxième position dans la subordonnée introduite par \all{dass}.
\item Dans une affiche allemande, les noms communs ne prennent pas la majuscule.
\end{enumerate}
\end{exercice}

\begin{exercice}[Vocabulaire : les médias]
Complète chaque phrase avec le mot allemand qui convient, choisi dans la liste : \all{die Quelle} – \all{die Sendung} – \all{die Schlagzeile} – \all{der Zuschauer} – \all{die Pressefreiheit} – \all{verbreiten}.
\begin{enumerate}
\item Die \dots{} der Zeitung lautet heute: „Wahlen in Kamerun“.
\item Jeden Abend um 20 Uhr sehe ich mir die \dots{} im Fernsehen an.
\item Man soll immer die \dots{} überprüfen, bevor man eine Nachricht teilt.
\item Der \dots{} schaltet den Fernseher ein und schaut die Nachrichten.
\item Ohne \dots{} können Journalisten nicht frei arbeiten.
\item Fake News \dots{} sich im Internet sehr schnell.
\end{enumerate}
\end{exercice}

\begin{exercice}[Conjugaison : l'impératif]
Mets les verbes suivants à l'impératif aux trois formes (\all{du}, \all{ihr}, \all{Sie}). Rédige ensuite des phrases complètes sur le thème des médias.
\begin{center}
\begin{tabular}{|l|l|l|l|}
\hline
\rowcolor{popblueL} Verbe & du & ihr & Sie \\ \hline
lesen & \dots & \dots & \dots \\ \hline
einschalten & \dots & \dots & \dots \\ \hline
sein & \dots & \dots & \dots \\ \hline
nehmen & \dots & \dots & \dots \\ \hline
überprüfen & \dots & \dots & \dots \\ \hline
\end{tabular}
\end{center}
Exemple : \all{informieren} → \all{Informier dich! / Informiert euch! / Informieren Sie sich!}
\end{exercice}

\courssub{Je m'entraîne}

\begin{exercice}[Thème : phrases à traduire]
Traduis les phrases suivantes en allemand. Attention à l'impératif, au passif et à l'ordre des mots.
\begin{enumerate}
\item Informez-vous, mais ne croyez pas tout ce que vous lisez sur Internet !
\item Les nouvelles sont diffusées chaque soir à vingt heures.
\item Ne pas utiliser les portables pendant l'émission !
\item Vérifiez les sources avant de partager une information !
\item L'article a été publié hier par la rédaction.
\end{enumerate}
\end{exercice}

\begin{exercice}[Transformation : de l'actif au passif]
Mets les phrases suivantes au passif, en conservant le temps indiqué.
\begin{enumerate}
\item \all{Die Redaktion veröffentlicht den Artikel.} (Präsens)
\item \all{Der Sender zeigt jeden Abend die Nachrichten.} (Präsens)
\item \all{Die Journalisten interviewten den Minister.} (Präteritum)
\item \all{Man verbreitete die Fake News im Internet.} (Präteritum)
\item \all{Die Schüler gestalten eine Plakatwand für die Medienwoche.} (Präsens)
\end{enumerate}
\end{exercice}

\begin{exercice}[Discours indirect : les déclarations du ministre]
Rapporte les déclarations suivantes au discours indirect avec le Konjunktiv I, après l'introduction : \all{Der Minister erklärte, \dots}
\begin{enumerate}
\item \all{„Die Presse ist frei.“}
\item \all{„Die Journalisten arbeiten ohne Angst.“}
\item \all{„Das Internet hat die Kommunikation verändert.“}
\item \all{„Wir werden ein neues Mediengesetz vorbereiten.“}
\item \all{„Die Bürger sollen die Quellen überprüfen.“}
\end{enumerate}
\end{exercice}

\begin{textebox}[Règlement de la salle multimédia — texte à compléter]
Liebe Schülerinnen und Schüler!

\dots{} bitte leise in der Mediathek! \dots{} die Computer nur für die Recherche! \dots{} eure Handys vor dem Unterricht \dots{}! \dots{} keine Getränke an die Computer! Wenn ein Problem auftritt, \dots{} euch bei der Lehrerin! Nach der Stunde \dots{} den Raum bitte ordentlich!
\end{textebox}

\begin{exercice}[Texte à trous et appariement]
\textbf{A. Texte à trous.} Complète le règlement de la salle multimédia ci-dessus avec les formes qui conviennent (impératif \all{ihr}) : \all{benutzen} – \all{ausschalten} – \all{sein} – \all{melden} – \all{verlassen} – \all{bringen}.

\textbf{B. Appariement.} Associe chaque slogan allemand (1--5) à son équivalent français (a--e).
\begin{enumerate}
\item \all{Lesen bildet!}
\item \all{Überprüfe die Quellen!}
\item \all{Schaltet ein – wir informieren euch!}
\item \all{Teile keine Fake News!}
\item \all{Bleiben Sie informiert!}
\end{enumerate}
\begin{enumerate}
\item[a)] Restez informés !
\item[b)] Ne partage pas de fausses informations !
\item[c)] La lecture instruit !
\item[d)] Branchez-vous – nous vous informons !
\item[e)] Vérifie les sources !
\end{enumerate}
\end{exercice}

\courssub{Je me prépare au Bac}

\begin{textebox}[Medienwoche am Gymnasium Yaoundé]
Liebe Schülerinnen und Schüler!

Vergleicht die Nachrichten verschiedener Sender! Das ist das Motto unserer Medienwoche, die am Montag in der Aula eröffnet wird. Der Direktor erklärte gestern, die Medienwoche werde von der Deutsch-AG organisiert und alle Klassen könnten daran teilnehmen.

Während der Woche werden Zeitungsartikel analysiert, Radiosendungen aufgenommen und Plakate gestaltet. Ein Journalist aus Douala hält am Mittwoch einen Vortrag über Fake News. Er sagte in einem Interview: „Junge Leute glauben oft alles, was sie im Internet lesen.“ Deshalb gilt unsere wichtigste Regel: Überprüft die Quellen, bevor ihr eine Nachricht teilt!

Am Freitag wird die beste Plakatwand von einer Jury ausgewählt. Der Gewinner erhält ein Abonnement einer deutschsprachigen Zeitschrift. Also: Macht mit, informiert euch und bleibt kritisch!
\end{textebox}

\begin{exercice}[Compréhension de l'écrit et questions de langue]
Lis le texte « Medienwoche am Gymnasium Yaoundé » ci-dessus, puis réponds aux questions \textbf{en allemand} par des phrases complètes.

\textbf{Questions de compréhension :}
\begin{enumerate}
\item Was ist das Motto der Medienwoche?
\item Wer organisiert die Medienwoche? (Achtung: der Text steht hier im Passiv und im indirekten Diskurs.)
\item Welche drei Aktivitäten sind für die Woche geplant?
\item Was kritisiert der Journalist aus Douala?
\item Was bekommt der Gewinner am Freitag?
\end{enumerate}

\textbf{Questions de langue :}
\begin{enumerate}
\item Relève dans le texte trois impératifs et précise à quelle forme (\all{du}, \all{ihr} ou \all{Sie}) ils sont construits.
\item Relève deux phrases au passif et remets-les à l'actif.
\item Relève une phrase au discours indirect et remets-la au discours direct.
\end{enumerate}
\end{exercice}

\begin{textebox}[Texte à traduire — thème]
Chers élèves ! Informez-vous, mais ne croyez pas tout ce que vous lisez sur Internet ! Les nouvelles sont diffusées chaque soir à vingt heures. Le journaliste a déclaré que la liberté de la presse était en danger. Pendant l'émission, les portables ne doivent pas être utilisés. Vérifiez toujours les sources !
\end{textebox}

\begin{textebox}[Text zu übersetzen — version]
Liebe Schüler! Vergleicht die Nachrichten verschiedener Sender! Der Direktor erklärte, die Medienwoche werde am Montag eröffnet. Fake News werden im Internet schnell verbreitet – überprüft die Quellen! Wer gut informiert sein will, muss mehrere Zeitungen lesen.
\end{textebox}

\begin{exercice}[Thème et version type Baccalauréat]
\textbf{A. Thème (français → allemand).} Traduis en allemand le texte « Texte à traduire — thème » ci-dessus.

\textbf{B. Version (allemand → français).} Traduis en français le texte « Text zu übersetzen — version » ci-dessus.

\textbf{Barème indicatif :} respect du sens (50 \%), grammaire : impératif, passif, discours indirect (30 \%), orthographe et majuscules (20 \%).
\end{exercice}

\begin{textebox}[Règlement fautif à corriger]
liebe schüler!

Benutzt ihr die Computer nur für die recherche! Schaltet aus eure handys vor dem Unterricht! Der Lehrer sagte, dass die Schüler müssen leise sein. Essen und Trinken wird in der Mediathek nicht erlaubt. Nach der Stunde macht zu die Fenster und aufgeräumt den Raum! Die Bücher werden von den Schülern zurück ins Regal gestellt müssen.
\end{textebox}

\begin{exercice}[Production écrite et correction guidée type Bac]
\textbf{A. Production écrite (80 à 120 mots).} Rédige en allemand une affiche ou un court article pour le journal de ton lycée intitulé : \all{„Medien und wir“}. Tu dois obligatoirement utiliser :
\begin{itemize}
\item au moins \textbf{un impératif} (forme \all{ihr} ou \all{Sie}) ;
\item au moins \textbf{une phrase au passif} ;
\item au moins \textbf{une phrase au discours indirect} (avec \all{dass} ou sans conjonction) ;
\item au moins \textbf{cinq mots} du lexique des médias étudié dans le chapitre.
\end{itemize}

\textbf{B. Correction guidée.} Le règlement « Règlement fautif à corriger » ci-dessus contient \textbf{8 erreurs} (ordre des mots, cas, majuscules, particule séparable, conjugaison). Retrouve-les et réécris le texte correctement.
\end{exercice}

\begin{methode}[Réussir l'épreuve de traduction et de production au Bac]
\begin{enumerate}
\item \textbf{Repère les structures obligatoires} avant d'écrire : y a-t-il un impératif, un passif, une phrase rapportée ? Souligne-les dans le texte de départ.
\item \textbf{Traduis phrase par phrase}, en plaçant d'abord le verbe conjugué à la bonne position (2\ieme{} position à l'indicatif, 1\iere{} à l'impératif, fin de phrase après \all{dass}).
\item \textbf{Au passif}, vérifie l'auxiliaire (\all{werden}, jamais \all{sein} pour le passif d'action) et la place du participe II en fin de phrase.
\item \textbf{Au discours indirect}, choisis le Konjunktiv I ; s'il est identique à l'indicatif, remplace-le par le Konjunktiv II (\all{sei}, \all{habe}, \all{werde}, \all{könne} ; sinon \all{arbeiteten}, \all{hätte}, \all{würde}).
\item \textbf{Relis-toi} avec la grille : majuscule à tous les noms, Umlaute, particules séparables en fin de phrase, accord de l'auxiliaire avec le sujet, cas après \all{von} (datif) et \all{während} (génitif).
\end{enumerate}
\end{methode}

\begin{attention}[Les 5 pièges fréquents des francophones]
\begin{itemize}
\item Ne traduis jamais « être + participe » français par \all{sein} + participe : « les nouvelles sont diffusées » = \all{die Nachrichten \textbf{werden} gesendet} (passif d'action = \all{werden}).
\item À l'impératif \all{du} et \all{ihr}, \textbf{pas de pronom} : \all{Überprüft die Quellen!} et non \all{Überprüft ihr die Quellen!}
\item Le verbe à particule se sépare aussi à l'impératif : \all{Schaltet die Handys \textbf{aus}!}
\item Après \all{dass}, le verbe conjugué va \textbf{à la fin} : \all{Er sagte, dass die Presse frei sei.}
\item Tous les noms prennent la majuscule, même dans les slogans : \all{die Quelle}, \all{das Handy}, \all{die Recherche}.
\end{itemize}
\end{attention}

\newpage

\coursec{Corrigés des exercices}

\begin{corrige}[Exercice 1 — QCM : grammaire des textes injonctifs]
\begin{enumerate}
\item \textbf{b)} \all{lies!} — les verbes qui changent \all{e} en \all{i/ie} au radical à la 2\ieme{} et 3\ieme{} personnes du singulier gardent ce changement à l'impératif \all{du} : \all{lesen} → \all{lies!}, \all{sprechen} → \all{sprich!}, \all{nehmen} → \all{nimm!}. La forme \all{les!} est fausse ; \all{lest!} est la forme \all{ihr} ; \all{lesen Sie!} est la forme de politesse.
\item \textbf{b)} \all{Er sagte, die Zeitung erscheine täglich.} — le discours indirect exige le Konjunktiv I (\all{erscheine}) ; sans conjonction \all{dass}, le verbe conjugué reste en deuxième position.
\item \textbf{a)} passif Präsens : \all{wird} (présent de \all{werden}) + participe II \all{geschrieben} en fin de phrase.
\item \textbf{b)} \all{Benutzen Sie bitte keine Handys!} — à la forme \all{Sie}, le verbe prend la forme de l'infinitif et le pronom \all{Sie} suit immédiatement le verbe. Les formes a) et c) placent le verbe à une position impossible.
\end{enumerate}
\end{corrige}

\begin{corrige}[Exercice 2 — Vrai ou faux ?]
\begin{enumerate}
\item \textbf{V.} À l'impératif, le verbe conjugué occupe la première position : \all{Öffnet die Bücher!} (« Ouvrez les livres ! »).
\item \textbf{V.} \all{sein} est irrégulier : \all{sei} (du), \all{seid} (ihr), \all{seien Sie} (Sie).
\item \textbf{V.} Passif Präteritum : \all{wurde(n)} + participe II, ex. \all{Der Artikel wurde geschrieben.} (« L'article fut écrit. »).
\item \textbf{F.} Avec \all{dass}, le verbe conjugué se rejette à la fin : \all{Er sagte, dass die Zeitung täglich erscheine.} C'est seulement \emph{sans} conjonction que le verbe reste en deuxième position.
\item \textbf{F.} En allemand, \textbf{tous} les noms communs prennent la majuscule, y compris sur les affiches et dans les slogans : \all{die Medienwoche}, \all{das Plakat}, \all{die Quelle}.
\end{enumerate}
\end{corrige}

\begin{corrige}[Exercice 3 — Vocabulaire : les médias]
\begin{enumerate}
\item \all{die Schlagzeile, -n} (« le gros titre ») : \all{Die Schlagzeile der Zeitung lautet heute: „Wahlen in Kamerun“.}
\item \all{die Sendung, -en} (« l'émission ») : \all{Jeden Abend um 20 Uhr sehe ich mir die Sendung im Fernsehen an.}
\item \all{die Quelle, -n} (« la source ») : \all{Man soll immer die Quelle überprüfen, bevor man eine Nachricht teilt.}
\item \all{der Zuschauer, -} (« le téléspectateur ») : \all{Der Zuschauer schaltet den Fernseher ein und schaut die Nachrichten.}
\item \all{die Pressefreiheit} (« la liberté de la presse », sans pluriel) : \all{Ohne Pressefreiheit können Journalisten nicht frei arbeiten.}
\item \all{sich verbreiten} (« se propager », verbe pronominal) : \all{Fake News verbreiten sich im Internet sehr schnell.} (ici au pluriel : \all{verbreiten sich})
\end{enumerate}
\end{corrige}

\begin{corrige}[Exercice 4 — Conjugaison : l'impératif]
\begin{center}
\begin{tabular}{|l|l|l|l|}
\hline
\rowcolor{popblueL} Verbe & du & ihr & Sie \\ \hline
lesen & lies & lest & lesen Sie \\ \hline
einschalten & schalte \dots\ ein & schaltet \dots\ ein & schalten Sie \dots\ ein \\ \hline
sein & sei & seid & seien Sie \\ \hline
nehmen & nimm & nehmt & nehmen Sie \\ \hline
überprüfen & überprüfe & überprüft & überprüfen Sie \\ \hline
\end{tabular}
\end{center}
Phrases possibles : \all{Lies die Schlagzeilen! / Lest den Artikel! / Lesen Sie die Zeitung!} — \all{Schalte das Radio ein! / Schaltet den Fernseher ein! / Schalten Sie bitte das Handy aus!} — \all{Sei kritisch! / Seid vorsichtig im Internet! / Seien Sie bitte leise!} — \all{Nimm die Zeitung! / Nehmt Platz! / Nehmen Sie das Mikrofon!} — \all{Überprüfe die Quelle! / Überprüft die Nachrichten! / Überprüfen Sie die Informationen!}

Points de vigilance : \all{einschalten} est séparable (la particule \all{ein} va en fin de phrase) ; \all{überprüfen} est inséparable (préfixe \all{über-} non accentué) ; \all{nehmen} perd le \all{h} et double le \all{m} : \all{nimm!}.
\end{corrige}

\begin{corrige}[Exercice 5 — Thème : phrases à traduire]
\begin{enumerate}
\item \all{Informieren Sie sich, aber glauben Sie nicht alles, was Sie im Internet lesen!} (ou forme \all{ihr} : \all{Informiert euch, aber glaubt nicht alles, was ihr im Internet lest!}) — verbe pronominal \all{sich informieren} ; subordonnée relative avec \all{was} et verbe en fin.
\item \all{Die Nachrichten werden jeden Abend um zwanzig Uhr gesendet.} (passif Präsens : \all{werden} + participe II \all{gesendet} en fin de phrase ; jamais \all{sind gesendet} pour un passif d'action. Remarque : \all{senden} admet deux participes, \all{gesendet} pour la diffusion radiophonique, \all{gesandt} pour l'envoi d'une lettre.)
\item \all{Handys während der Sendung nicht benutzen!} (tournure à l'infinitif, typique des consignes et des règlements ; \all{während} + génitif : \all{der Sendung})
\item \all{Überprüfen Sie die Quellen, bevor Sie eine Information teilen!} (ou \all{Überprüft die Quellen, bevor ihr eine Information teilt!}) — après \all{bevor}, verbe conjugué en fin de subordonnée.
\item \all{Der Artikel wurde gestern von der Redaktion veröffentlicht.} (passif Präteritum : \all{wurde} + participe II ; « par » = \all{von} + datif : \all{von der Redaktion})
\end{enumerate}
\end{corrige}

\begin{corrige}[Exercice 6 — Transformation : de l'actif au passif]
\begin{enumerate}
\item \all{Der Artikel wird von der Redaktion veröffentlicht.} (l'accusatif \all{den Artikel} devient sujet au nominatif)
\item \all{Die Nachrichten werden jeden Abend vom Sender gezeigt.} (\all{von dem} → contraction \all{vom} ; pluriel : \all{werden})
\item \all{Der Minister wurde von den Journalisten interviewt.} (Präteritum : \all{wurde})
\item \all{Die Fake News wurden im Internet verbreitet.} (pluriel : \all{wurden} ; le sujet \all{man} disparaît au passif)
\item \all{Eine Plakatwand für die Medienwoche wird von den Schülern gestaltet.} (singulier : \all{wird})
\end{enumerate}
\end{corrige}

\begin{corrige}[Exercice 7 — Discours indirect : les déclarations du ministre]
\begin{enumerate}
\item \all{Der Minister erklärte, die Presse sei frei.} (\all{ist} → Konjunktiv I \all{sei})
\item \all{Der Minister erklärte, die Journalisten arbeiteten ohne Angst.} (le Konjunktiv I \all{arbeiten} est identique à l'indicatif : on le remplace par le Konjunktiv II \all{arbeiteten})
\item \all{Der Minister erklärte, das Internet habe die Kommunikation verändert.} (\all{hat} → \all{habe})
\item \all{Der Minister erklärte, man werde ein neues Mediengesetz vorbereiten.} (ou : \all{sie würden ein neues Mediengesetz vorbereiten} ; le « nous » du discours direct devient \all{man} ou \all{sie} ; \all{werden} + infinitif → \all{werde} + infinitif)
\item \all{Der Minister erklärte, die Bürger sollten die Quellen überprüfen.} (\all{sollen} → \all{sollten}, Konjunktiv II de substitution, car le Konjunktiv I \all{sollen} est identique à l'indicatif)
\end{enumerate}
\end{corrige}

\begin{corrige}[Exercice 8 — Texte à trous et appariement]
\textbf{A.} \all{Seid bitte leise in der Mediathek! Benutzt die Computer nur für die Recherche! Schaltet eure Handys vor dem Unterricht aus! Bringt keine Getränke an die Computer! Wenn ein Problem auftritt, meldet euch bei der Lehrerin! Nach der Stunde verlasst den Raum bitte ordentlich!}

Points de vigilance : \all{ausschalten} et \all{sich melden} sont des verbes à particule séparable — la particule \all{aus} va en fin de phrase et le pronom réfléchi \all{euch} suit le verbe ; \all{verlassen} est inséparable (préfixe \all{ver-}) et ne se sépare jamais ; \all{sein} → \all{seid} (irrégulier).

\textbf{B.} 1--c (\all{Lesen bildet!} = « La lecture instruit ! ») ; 2--e ; 3--d ; 4--b ; 5--a.
\end{corrige}

\begin{corrige}[Exercice 9 — Compréhension de l'écrit et questions de langue]
\textbf{Compréhension :}
\begin{enumerate}
\item \all{Das Motto lautet: „Vergleicht die Nachrichten verschiedener Sender!“}
\item \all{Die Medienwoche wird von der Deutsch-AG organisiert.} (passif Präsens ; \all{von} + datif désigne l'agent)
\item \all{Geplant sind: Zeitungsartikel werden analysiert, Radiosendungen werden aufgenommen und Plakate werden gestaltet.}
\item \all{Er kritisiert, dass junge Leute oft alles glauben, was sie im Internet lesen.}
\item \all{Der Gewinner bekommt ein Abonnement einer deutschsprachigen Zeitschrift.}
\end{enumerate}
\textbf{Langue :}
\begin{enumerate}
\item \all{Vergleicht}, \all{Überprüft}, \all{Macht mit}, \all{informiert}, \all{bleibt} : tous à la forme \all{ihr} (radical + \all{-t}, sans pronom). \all{Macht mit} : particule séparable \all{mit} en fin de phrase.
\item \all{Die Medienwoche wird am Montag in der Aula eröffnet.} → actif : \all{Man eröffnet die Medienwoche am Montag in der Aula.} ; \all{Am Freitag wird die beste Plakatwand von einer Jury ausgewählt.} → actif : \all{Am Freitag wählt eine Jury die beste Plakatwand aus.}
\item \all{Der Direktor erklärte, die Medienwoche werde von der Deutsch-AG organisiert.} → discours direct : \all{Der Direktor erklärte: „Die Medienwoche wird von der Deutsch-AG organisiert.“} (Konjunktiv I \all{werde} → indicatif \all{wird})
\end{enumerate}
\end{corrige}

\begin{corrige}[Exercice 10 — Thème et version type Baccalauréat]
\textbf{A. Thème — proposition de traduction :}

\all{Liebe Schüler! Informiert euch, aber glaubt nicht alles, was ihr im Internet lest! Die Nachrichten werden jeden Abend um zwanzig Uhr gesendet. Der Journalist erklärte, die Pressefreiheit sei in Gefahr. Während der Sendung dürfen die Handys nicht benutzt werden. Überprüft immer die Quellen!}

Points de vigilance : impératif \all{ihr} sans pronom ; passif Präsens \all{werden gesendet} ; discours indirect au Konjunktiv I \all{sei} ; passif avec modal \all{dürfen nicht benutzt werden} (modal conjugué en 2\ieme{} position, \all{werden} à la fin) ; \all{während} + génitif.

\textbf{B. Version — proposition de traduction :}

« Chers élèves ! Comparez les informations de différentes chaînes ! Le directeur a déclaré que la semaine des médias serait inaugurée lundi. Les fausses informations se propagent vite sur Internet — vérifiez les sources ! Qui veut être bien informé doit lire plusieurs journaux. »

Points de vigilance : \all{werden verbreitet} = passif présent (« sont propagées / se propagent ») ; \all{werde eröffnet} = discours indirect au passif, rendu par le conditionnel (« serait inaugurée ») ; \all{Wer \dots\ will, muss \dots} = « qui veut\dots doit\dots ».

\textbf{Barème indicatif :} respect du sens 50 \% ; grammaire (impératif, passif, discours indirect, ordre des mots) 30 \% ; orthographe, majuscules et Umlaute 20 \%.
\end{corrige}

\begin{corrige}[Exercice 11 — Production écrite et correction guidée]
\textbf{A. Exemple de production (à titre indicatif, environ 100 mots) :}

\all{Medien und wir — Liebe Schülerinnen und Schüler! Jeden Tag werden wir mit Nachrichten überschwemmt: im Radio, im Fernsehen und im Internet. Aber nicht jede Information ist wahr. Fake News werden oft schneller verbreitet als die Wahrheit. Unser Deutschlehrer sagte gestern, dass wir die Quellen immer überprüfen müssten. Er hat recht! Deshalb: Lest mehrere Zeitungen, vergleicht die Schlagzeilen verschiedener Sender und teilt keine Nachricht ohne Prüfung! Bleibt kritisch und gut informiert — die Pressefreiheit braucht aufmerksame Leser!}

Cette production contient : des impératifs \all{ihr} (\all{Lest}, \all{vergleicht}, \all{teilt}, \all{Bleibt}), deux passifs (\all{werden wir überschwemmt}, \all{werden verbreitet}), une phrase au discours indirect avec \all{dass} (\all{dass wir die Quellen immer überprüfen müssten}) et le lexique des médias (\all{Nachrichten}, \all{Fake News}, \all{Quellen}, \all{Schlagzeilen}, \all{Sender}, \all{Pressefreiheit}).

\textbf{B. Les 8 erreurs et le texte corrigé :}
\begin{enumerate}
\item \all{liebe schüler} → \all{Liebe Schüler} (deux majuscules : début de phrase et nom commun) ;
\item \all{für die recherche} → \all{für die Recherche} (majuscule au nom) ;
\item \all{Benutzt ihr} → \all{Benutzt} (à l'impératif \all{ihr}, pas de pronom) ;
\item \all{Schaltet aus eure handys} → \all{Schaltet eure Handys aus} (particule séparable en fin de phrase + majuscule) ;
\item \all{dass die Schüler müssen leise sein} → \all{dass die Schüler leise sein müssen} (verbe conjugué à la fin dans la subordonnée) ;
\item \all{Essen und Trinken wird} → \all{Essen und Trinken werden} (deux sujets = pluriel) ;
\item \all{macht zu die Fenster und aufgeräumt den Raum} → \all{macht die Fenster zu und räumt den Raum auf} (particules séparables en fin de phrase ; impératif \all{räumt \dots\ auf} et non participe II) ;
\item \all{werden \dots\ gestellt müssen} → \all{müssen \dots\ gestellt werden} (passif avec modal : modal conjugué en 2\ieme{} position, \all{werden} à la fin).
\end{enumerate}
Texte corrigé :

\all{Liebe Schüler! Benutzt die Computer nur für die Recherche! Schaltet eure Handys vor dem Unterricht aus! Der Lehrer sagte, dass die Schüler leise sein müssen. Essen und Trinken werden in der Mediathek nicht erlaubt. Nach der Stunde macht die Fenster zu und räumt den Raum auf! Die Bücher müssen von den Schülern zurück ins Regal gestellt werden.}
\end{corrige}

\begin{methode}[Réussir l'épreuve de traduction et de production au Bac]
\begin{enumerate}
\item \textbf{Repère les structures obligatoires} avant d'écrire : y a-t-il un impératif, un passif, une phrase rapportée ? Souligne-les dans le texte de départ.
\item \textbf{Traduis phrase par phrase}, en plaçant d'abord le verbe conjugué à la bonne position (2\ieme{} position à l'indicatif, 1\iere{} à l'impératif, fin de phrase après \all{dass}).
\item \textbf{Au passif}, vérifie l'auxiliaire (\all{werden}, jamais \all{sein} pour le passif d'action) et la place du participe II en fin de phrase.
\item \textbf{Au discours indirect}, choisis le Konjunktiv I ; s'il est identique à l'indicatif, remplace-le par le Konjunktiv II (\all{sei}, \all{habe}, \all{werde}, \all{könne} ; sinon \all{arbeiteten}, \all{hätte}, \all{würde}).
\item \textbf{Relis-toi} avec la grille : majuscule à tous les noms, Umlaute, particules séparables en fin de phrase, accord de l'auxiliaire avec le sujet, cas après \all{von} (datif) et \all{während} (génitif).
\end{enumerate}
\end{methode}

\begin{attention}[Les 5 pièges fréquents des francophones]
\begin{itemize}
\item Ne traduis jamais « être + participe » français par \all{sein} + participe : « les nouvelles sont diffusées » = \all{die Nachrichten \textbf{werden} gesendet} (passif d'action = \all{werden}).
\item À l'impératif \all{du} et \all{ihr}, \textbf{pas de pronom} : \all{Überprüft die Quellen!} et non \all{Überprüft ihr die Quellen!}
\item Le verbe à particule se sépare aussi à l'impératif : \all{Schaltet die Handys \textbf{aus}!}
\item Après \all{dass}, le verbe conjugué va \textbf{à la fin} : \all{Er sagte, dass die Presse frei sei.}
\item Tous les noms prennent la majuscule, même dans les slogans : \all{die Quelle}, \all{das Handy}, \all{die Recherche}.
\end{itemize}
\end{attention}

\newpage

\coursec{Fiche bilan}

\begin{popfigure}
\begin{tikzpicture}[mindmap, grow cyclic, every node/.style={concept, font=\small},
root concept/.append style={concept color=popblue, font=\bfseries\small, minimum size=3.2cm},
level 1/.append style={level distance=3.6cm, sibling angle=60},
level 1 concept/.append style={minimum size=2.4cm, font=\scriptsize}]
\node[root concept]{Médias et communication : textes injonctifs}
  child[concept color=poporange]{ node {Lexique : \all{die Medien}, \all{die Presse}, \all{die Werbung}} }
  child[concept color=popgreen]{ node {Impératif : \all{Lies!}, \all{Lesen Sie!}, \all{Lasst uns lesen!}} }
  child[concept color=poppurple]{ node {Injonctions : \all{Infinitiv}, \all{Passiv}, \all{Modalverben}} }
  child[concept color=popgold]{ node {Discours indirect : \all{Konjunktiv I}, \all{dass}-Satz} }
  child[concept color=poppink]{ node {Passif : \all{werden} + Partizip II ; \all{von}/\all{durch}} }
  child[concept color=popblue!70]{ node {Traduction : thème et version, ordre des mots} }
  child[concept color=poporange!70]{ node {Production : \all{Plakat}, \all{Regeln}, \all{Slogan}} };
\end{tikzpicture}
\legende{Carte mentale de la leçon : médias, communication et textes injonctifs}
\end{popfigure}

\begin{bilan}
\textbf{1. Lexique clé (à connaître avec l'article et le pluriel)}
\begin{itemize}
\item \all{die Zeitung, -en} : le journal ; \all{die Zeitschrift, -en} : la revue ; \all{der Artikel, -} : l'article
\item \all{die Nachrichten (Pl.)} : les informations ; \all{die Schlagzeile, -n} : le gros titre
\item \all{das Fernsehen} : la télévision ; \all{das Radio, -s} : la radio ; \all{die Sendung, -en} : l'émission
\item \all{das Internet} : l'Internet ; \all{die sozialen Netzwerke (Pl.)} : les réseaux sociaux
\item \all{die Werbung} : la publicité ; \all{das Plakat, -e} : l'affiche ; \all{der Slogan, -s} : le slogan
\item \all{die Meinungsfreiheit} : la liberté d'opinion ; \all{die Pressefreiheit} : la liberté de la presse
\item \all{berichten} : rapporter, informer ; \all{veröffentlichen} : publier ; \all{verboten} : interdit ; \all{erforderlich} : obligatoire
\end{itemize}

\textbf{2. Grammaire à connaître par cœur}
\begin{center}
\begin{tabularx}{\linewidth}{|l|X|X|}
\hline
\rowcolor{popblueL} Point & Règle & Exemple \\
\hline
Impératif (2\ieme{} pers. sing.) & radical du verbe (+ e), sans pronom & \all{Lies den Artikel!} « Lis l'article ! » \\
\hline
Impératif (\all{Sie}) & infinitif + \all{Sie} (verbe en tête) & \all{Lesen Sie die Zeitung!} « Lisez le journal ! » \\
\hline
Impératif (\all{ihr}) & 2\ieme{} pers. plur. sans pronom & \all{Schaltet das Radio ein!} « Allumez la radio ! » \\
\hline
Consignes impersonnelles & infinitif en fin de phrase ou substantif & \all{Nicht hinauslehnen!} ; \all{Rauchen verboten!} \\
\hline
Discours indirect & \all{Konjunktiv I} ou \all{dass} + verbe en fin & \all{Er sagte, er lese die Zeitung.} / \all{..., dass er die Zeitung lese.} \\
\hline
Passif & \all{werden} + Partizip II ; agent avec \all{von} & \all{Der Artikel wird von der Redaktion geschrieben.} \\
\hline
Passif aux autres temps & Präteritum : \all{wurde} ; Perfekt : \all{ist ... worden} & \all{Die Sendung wurde gesendet.} \\
\hline
\end{tabularx}
\end{center}

\textbf{3. Méthodes express}
\begin{itemize}
\item \textbf{Version (allemand $\rightarrow$ français)} : repérer le verbe conjugué (2\ieme{} place), le sujet, les compléments ; traduire le cadre avant les détails ; vérifier les temps (Präsens, Perfekt, Präteritum).
\item \textbf{Thème (français $\rightarrow$ allemand)} : majuscule à tous les noms ; verbe conjugué en 2\ieme{} position ; particule séparable et participe passé en fin de phrase ; attention aux cas après prépositions.
\item \textbf{Affiche / règlement} : phrases courtes, impératif ou infinitif d'injonction, slogans percutants, vocabulaire des médias.
\item \textbf{Relecture} : orthographe, Umlaute, majuscules, ordre des mots, accord des cas.
\end{itemize}
\end{bilan}

\begin{aretenir}[Checklist avant le Bac]
\begin{itemize}
\item[$\square$] Je connais le vocabulaire des médias avec article et pluriel (\all{die Zeitung, -en}).
\item[$\square$] Je sais former les trois impératifs : \all{Lies!}, \all{Lest!}, \all{Lesen Sie!}.
\item[$\square$] Je sais rédiger une consigne impersonnelle à l'infinitif : \all{Bitte nicht stören!}.
\item[$\square$] Je place le verbe conjugué en deuxième position dans la phrase déclarative.
\item[$\square$] Je mets le verbe en fin de phrase après \all{dass}, \all{weil}, \all{ob}.
\item[$\square$] Je forme le passif : \all{werden} + Partizip II ; \all{ist ... worden} au Perfekt.
\item[$\square$] Je rends le discours indirect avec le \all{Konjunktiv I} ou une subordonnée en \all{dass}.
\item[$\square$] Je mets la majuscule à tous les noms allemands et j'écris les Umlaute (ä, ö, ü, ß).
\item[$\square$] Je sépare correctement les verbes à particule : \all{Schaltet den Fernseher ein!}.
\item[$\square$] Je sais structurer une affiche : titre, slogan, consignes, informations pratiques.
\item[$\square$] Je relis ma traduction : temps, cas, ordre des mots, sens global.
\item[$\square$] J'évite les faux amis : \all{aktualisieren} n'est pas « actualiser » au sens de « rendre actuel » sans contexte ; \all{die Information} = le renseignement, \all{die Nachrichten} = les informations télévisées.
\end{itemize}
\end{aretenir}

\findecours

\end{document}
